\documentclass[aps,prd,twocolumn,nofootinbib,floatfix,superscriptaddress]{revtex4-2}

\usepackage{amsmath,amssymb,amsthm}
\usepackage{graphicx}
\usepackage{xcolor}
\usepackage{booktabs}
\usepackage{hyperref}
\usepackage{bm}
\usepackage{fancyhdr}
\hypersetup{colorlinks=true,linkcolor=blue,citecolor=blue,urlcolor=blue}

% ---- Manuscript versioning -------------------------------------------------
% version.tex is regenerated on every CI build from git metadata; see
% .github/workflows/build-pdf.yml. For local compiles without it, fall back
% to a self-evident placeholder so referees can never mistake one for the
% other.
\IfFileExists{version.tex}{\input{version}}{%
  \newcommand{\paperversion}{local-build}%
  \newcommand{\paperdate}{\today}%
}

% ---- Footer with manuscript revision (every page) --------------------------
\pagestyle{fancy}
\fancyhf{}
\fancyfoot[C]{\thepage}
\fancyfoot[L]{\footnotesize\textit{Manuscript \paperversion\ \textperiodcentered\ \paperdate}}
\renewcommand{\headrulewidth}{0pt}
\renewcommand{\footrulewidth}{0pt}
\AtBeginDocument{\pagestyle{fancy}}

% Status-badge macros
\newcommand{\PROVEN}{{\color{black!75!green}\textsc{[proven]}}}
\newcommand{\DEMO}{{\color{black!70!blue}\textsc{[demonstrated]}}}
\newcommand{\INDIC}{{\color{black!60!yellow}\textsc{[indicated]}}}
\newcommand{\OPEN}{{\color{black!60!red}\textsc{[open]}}}
\newcommand{\DROPPED}{{\color{red!70!black}\textsc{[dropped]}}}

\newcommand{\Mpl}{M_{\rm Pl}}
\newcommand{\dKL}{D_{\rm KL}}
\newcommand{\Lm}{\mathcal{L}_m}
\newcommand{\Teff}{T^{\rm eff}}
\newcommand{\fprim}{f_{\rm prim}}

\begin{document}

\title{Information-Enhanced Gravity:\\
A Scalar-Tensor Theory Without Dark Matter, Dark Energy, or Cosmological Constant}

\author{S. Brailsford}
\affiliation{Independent AI researcher}

\date{\today}

\begin{abstract}
We present Information-Enhanced Gravity (ISST), a scalar-tensor
theory defined by the action
$S = (16\pi)^{-1}\!\int\!\Psi R\sqrt{-g}\,d^4x +
\int\!(1{+}f)\Lm\sqrt{-g}\,d^4x$,
where $\Psi$ is a scalar field with no kinetic term and $f$ is the
Kullback--Leibler divergence of the matter phase-space distribution
from a maximally uninformative reference. The theory
\emph{eliminates dark matter as a particle species}, the
cosmological constant, and dark energy --- replacing them with
information-enhanced baryonic gravity and Wiltshire's backreaction
cosmology, which we derive as the simplest physically-motivated
inhomogeneous background compatible with the field equations
(strict exclusion of LTB/Szekeres alternatives identified as open). A Branch theorem
($R\nabla\Psi = 0$) partitions spacetime into two sectors. A
parameter-free static-fluid theorem
(App.~\ref{app:static_fluid_theorem}) forces all physical stars
into the GR-recovering branch --- not as a screening limit but
pointwise --- underwriting exact PPN parameters
($\gamma = \beta = 1$), gravitational-wave speed $c_{\rm GW} = c$
(structural, not tuned), and LLR bounds. From the Standard Model
thermal history we obtain $\fprim = 5.66 \pm 0.12$, predicting
$\Omega_m/\Omega_b$ within $4.2\%$ of Planck with no fitted
parameters ($1.6$--$1.8\sigma$ tension). The field equations yield
a first-principles acceleration scale $a_{\rm crit}$ within $11\%$
of the empirical MOND $a_0$ --- a quantity neither $\Lambda$CDM nor
MOND derives. Rotation curves for 175 SPARC galaxies are fitted
with universal parameters (no per-galaxy tuning); the morphological
breakdown is structurally predicted by the same operator that gives
Newtonian dynamics for NGC~1052-DF2. The CMB peak-height ratios
under the strict action remain open: the perturbation-level
mechanism producing collisionless-like growth from $(1{+}f)\Lm$
has not been derived, and a first-pass acoustic-angle result
($-0.83\%$ on an effective-FRW proxy) carries
${\sim}\,5$--$10\%$ systematic uncertainty. The $\fprim$ derivation
chain from the KL functional to the discrete DOF sum has a specific
structural obstruction (single-epoch functional vs.\ path-additive
history). Pre-committed falsification conditions include direct
dark-matter particle detection, $c_{\rm GW} \neq c$, and DESI~DR2
$f\sigma_8$ shape.
\end{abstract}

\maketitle
\thispagestyle{fancy}

%==========================================================
\section{Introduction}
\label{sec:intro}

Modified gravity theories face a recurring problem: they fix one
scale and break another. MOND~\cite{Milgrom1983} explains galaxy
rotation curves but has no cosmology. $f(R)$
theories~\cite{Sotiriou2010} address cosmic acceleration but
require screening mechanisms that are themselves fine-tuned.
Tensor-vector-scalar theories were largely killed by the GW170817
constraint $|c_{\rm GW}/c - 1| < 10^{-15}$~\cite{LIGO2017,Creminelli2017}.
The result is a landscape where every proposal succeeds locally and
fails globally. This paper presents a theory that attempts the full
range --- from Solar System to CMB --- within a single action. It
succeeds in some places, fails honestly in others, and identifies
precisely where the open problems lie.

\textbf{The action.} ISST modifies gravity through two mechanisms:
a non-minimally coupled scalar $\Psi$ that multiplies the Ricci
scalar (with no kinetic term), and an information-theoretic
enhancement $(1{+}f)$ of the matter Lagrangian, where $f$ measures
how far the matter distribution has departed from thermodynamic
equilibrium. The theory has no cosmological constant, no
dark-matter particle, and no dark-energy field. It has one action,
two branches, and a set of tensions it does not hide.

\textbf{Structural backbone.} Five results anchor the theory, each
derived from the action with no free parameters:
\begin{enumerate}
\item The \emph{Branch theorem}: the Bianchi identity applied to
the field equation forces $R\,\nabla_\nu\Psi = 0$ pointwise ---
spacetime partitions into regions where either the Ricci scalar
vanishes (Branch~A, cosmological) or $\Psi$ is constant (Branch~B,
compact objects).

\item The \emph{static-fluid theorem}
(App.~\ref{app:static_fluid_theorem}): for any static, spherically
symmetric perfect fluid on Branch~A, central pressure is forced to
$P_c = 0$ --- unphysical. All physical stars are therefore on
Branch~B, where the field equation reduces to Einstein's equations
with $G = (1{+}f)/\Psi_0$. This is exact GR recovery, not a
screening approximation. It underwrites the PPN results, the LLR
bound, the gravitational-wave speed, the rotating-neutron-star
extension, and the DF2 prediction --- half the paper's confirmed
results trace to this single derivation.

\item The \emph{acceleration scale}: the wall-Friedmann equation
and weak-field algebra yield
$a_{\rm crit} = \tfrac{3(\sqrt{21}-3)}{4}\,cH_0/(1{+}\fprim)
\approx 1.07\times 10^{-10}\,{\rm m\,s^{-2}}$, within $11\%$ of
the empirical MOND value. $\Lambda$CDM does not predict this
number; MOND takes it as an empirical input. ISST derives it from
two algebraic identities.

\item The \emph{cosmological-floor identity}:
$(1{+}\fprim)\Omega_b = \Omega_m$ overshoots the Planck value by
$4.2\%$ from a single Standard Model thermodynamic calculation
with no fitted parameters. We report this as a $1.6$--$1.8\sigma$
tension and do not absorb the residual.

\item \emph{Gravitational-wave speed}: $c_{\rm GW} = c$ follows
structurally from the absence of a kinetic term for $\Psi$ ---
the theory sits in the GW170817-surviving corner of Horndeski
space by construction, not by tuning a coefficient. Most
scalar-tensor theories were killed or constrained by that
measurement. ISST passes it trivially.
\end{enumerate}

\textbf{On the open CMB perturbation status.} The theory's most
significant limitation is stated plainly: the CMB peak-height
ratios are not yet derived from the action. The strict
uniform-$f$ perturbation source is falsified at $>10\sigma$ at
the third peak. An environment-dependent closure provides only
percent-level corrections, not the factor-$\sim\!30$ needed. A
first-pass acoustic-angle calculation gives $-0.83\%$ agreement
with Planck, but on an effective-FRW proxy whose background
diverges from the committed Wiltshire $H(z)$ by $+76\%$ at
recombination --- the systematic uncertainty is
${\sim}\,5$--$10\%$, larger than the headline number. A modified
Boltzmann code with the true two-domain background is identified
as the critical next computation. For context: MOND's relativistic
completions (AQUAL~\cite{Bekenstein1984}, 1984;
TeVeS~\cite{BekensteinTeVeS2004}, 2004) required several more years
of independent analysis before engaging CMB peak structure. ISST engages the CMB
in its first paper, identifies the failure mode, and states the
path to resolution. The gap is real. It is also young.

\textbf{Honesty as method.} This paper openly catalogues six
claims that were made and subsequently dropped during development
--- including a wrong-sign $S_8$ prediction, a BBN-catastrophe
artefact, and the strict uniform-$f$ CMB source. Most
modified-gravity papers do not publish their failed attempts. We
do, because a theory that cannot identify its own errors cannot
be trusted when it claims success. The full audit trail --- 88
results tagged as PROVEN, DEMONSTRATED, INDICATED, OPEN, or
DROPPED --- is maintained in the companion derivation ledger and
summarised in Table~\ref{tab:audit}.

\textbf{Collaboration note.} This work was produced as a research
collaboration between S.B.\ (an AI researcher) and Anthropic's
Claude. S.B.\ provided physical direction, adversarial pressure,
and all commitments to specific numbers; the AI performed
derivations, ran adversarial self-tests, and managed the
derivation-audit system. Neither role substitutes for the other.
The methodology has obvious epistemic limitations: adversarial
testing by a system that also advocates has a conflict of
interest, and the Derivation Passport's verdicts have not yet
been independently audited by a researcher outside this project.
That external validation is identified as a methodological
priority (Sec.~\ref{sec:future}), and the keystone derivations
(static-fluid theorem, App.~\ref{app:static_fluid_theorem}; ADM
constraint analysis, App.~\ref{app:adm_dof}) are presented in
full in the appendices for independent verification.

\textbf{A note on language.} The companion site
(\href{https://lily-labs.co.uk/isst}{lily-labs.co.uk/isst})
explains ISST in accessible terms. The language there prioritises
clarity over precision. For exact claims, derivation status,
caveats, error bars, and parameter accounting, this paper is the
authoritative source. All load-bearing derivations are presented
in the body or appendices; the companion site provides
supplementary computational outputs and the interactive Passport
engine. Code and modified Boltzmann configurations will be
archived on Zenodo upon acceptance.

%==========================================================
\section{The action and field equations}
\label{sec:action}

\subsection{Action and field equations}

ISST is defined by the action
\begin{equation}
\label{eq:action}
S \;=\; \frac{1}{16\pi}\!\int\! d^4x\sqrt{-g}\,\Psi R
\;+\; \int\! d^4x\sqrt{-g}\,(1{+}f)\,\Lm,
\end{equation}
with no kinetic term for $\Psi$ and no cosmological constant.
Variation with respect to $g^{\mu\nu}$ yields
\begin{equation}
\label{eq:eom}
\Psi G_{\mu\nu} + g_{\mu\nu}\Box\Psi - \nabla_\mu\nabla_\nu\Psi
= 8\pi\,\Teff_{\mu\nu},
\end{equation}
where $\Teff_{\mu\nu} = (1{+}f)T_{\mu\nu} + \Delta^f_{\mu\nu}$.
Under uniform $f$ (the scope of this paper) and on perfect-fluid
matter, $\Delta^f_{\mu\nu}$ is negligible at perfect-fluid scope
(bound $|\Delta^f|/[(1{+}f)T] \lesssim \alpha(\sigma/c)^2 \sim 10^{-8}$
galactic, $\sim 10^{-5}$ cosmological-dust; the bound vanishes
identically under the ratchet postulate (Sec.~\ref{sec:ratchet}),
which is the operative assumption for this paper's perfect-fluid
applications; see Sec.~\ref{sec:psi_variational} below). The action
is varied with respect to $g^{\mu\nu}$ and matter fields only;
$\Psi$ is treated as trace-determined (metrically slaved through
the trace of \eqref{eq:eom}), with no independent Euler--Lagrange
equation (see Sec.~\ref{sec:psi_variational}). The status of this
construction is \PROVEN\ at all foundational layers and consistent
with the structural commitments enforced by the
Derivation Passport (App.~\ref{app:passport}).

\subsection{Variational status of $\Psi$}
\label{sec:psi_variational}

The scalar $\Psi$ is treated as \emph{trace-determined}: the action
\eqref{eq:action} is varied with respect to $g^{\mu\nu}$ and matter
fields $\chi_i$, but \emph{not} with respect to $\Psi$. There is
therefore no Euler--Lagrange equation $\delta S/\delta\Psi = 0$ as
an independent extremisation; the pointwise constraint
$R\nabla_\nu\Psi = 0$ (Sec.~\ref{sec:branch}) emerges from the
\emph{Bianchi identity} applied to the metric equation
\eqref{eq:eom}, combined with the matter Noether identity
$\nabla^\mu T^{\rm eff}_{\mu\nu} = 0$ that holds off-shell by the
diffeomorphism invariance of $S_m$. \emph{The trace of
\eqref{eq:eom}, $\Box\Psi = (8\pi/3)(1{+}f)T$, is a sourced
second-order PDE: $\Psi$ propagates in every PDE sense.} The
``trace-determined'' framing refers specifically to the variational
status: $\Psi$ is not independently extremised, and its evolution
is entirely determined by the metric and matter equations through
the trace constraint. Whether this constraint-propagation
constitutes an independent dynamical degree of freedom is a
Hamiltonian-constraint question (App.~\ref{app:adm_dof}), not a
PDE-classification question on the trace equation. The Hamiltonian
answer is that the $(\pi_\Psi, R=0)$ pair is first-class, generating
the gauge transformation $\delta\Psi = \varepsilon$ on the
constraint surface, and contributes zero propagating DOF beyond the
two graviton polarisations of GR. The role of $\Psi$ is closer to a
\emph{constraint field} or a metrically slaved state variable than
to a Lagrange multiplier in the standard variational sense (where
``Lagrange multiplier'' usually denotes an auxiliary field whose
EOM is itself the constraint); here the EOM-of-the-multiplier role
is played by the trace identity, not by $\delta S/\delta\Psi$. The
closest gravitational analogue is mimetic
gravity~\cite{ChamseddineMukhanov2013}, where a constraint scalar
enforces a structural condition without an extremisation principle.
ADM analysis (App.~\ref{app:adm_dof}) confirms zero scalar
propagating modes consistent with this interpretation. Both
branches~A and~B are admissible solutions of the metric equation;
selection is by boundary data (the static-fluid theorem,
App.~\ref{app:static_fluid_theorem}, forces Branch~B for static
compact sources; cosmological symmetry forces Branch~A on dust,
Sec.~\ref{sec:wiltshire_forced}). Treating $\delta S/\delta\Psi = 0$
as an EOM would force $R = 0$ everywhere with $f$ $\Psi$-independent,
eliminating Branch~B; this is not the F01 commitment.

\textbf{Distinction from mimetic gravity.} The structural similarity
of the F01 field equation \eqref{eq:eom} to mimetic
gravity~\cite{ChamseddineMukhanov2013} --- both feature second
derivatives of a constraint scalar in the gravitational sector
--- raises the question of whether ISST inherits the mimetic
hidden-mode controversy. In mimetic gravity, the scalar $\phi$ is
constrained by a Lagrange multiplier $\lambda$ enforcing
$g^{\mu\nu}\partial_\mu\phi\,\partial_\nu\phi = -1$, which produces
a propagating dust-like mode whose kinetic structure has been
extensively analysed in the
literature~\cite{Barvinsky2014,Capela2015,Firouzjahi2017,Gorji2018};
several authors argue that the constraint creates a hidden
propagating degree of freedom with non-trivial dispersion in the
caustic regime. \emph{The ISST construction does not contain this
constraint:} no Lagrange multiplier $\lambda$ enforcing a derivative
condition on $\Psi$ appears in the action \eqref{eq:action}, and
$\nabla\Psi$ is not constrained to be timelike unit. The
mimetic-dust mode controversy therefore does not directly transfer
to ISST. Independent ADM verification of the result in
App.~\ref{app:adm_dof} --- two tensor modes, no scalar mode ---
by an external group, with explicit reference to the mimetic-DOF
analyses of Barvinsky~\cite{Barvinsky2014},
Capela--Ramazanov~\cite{Capela2015}, and the
Firouzjahi--Gorji--Mansoori
sequence~\cite{Firouzjahi2017,Gorji2018}, is identified as a
methodological priority (Sec.~\ref{sec:future}). The ADM
calculation in App.~\ref{app:adm_dof} is internal to the project;
while it produces the expected count and is structurally consistent
with the trace-determined-$\Psi$ commitment, an outside check is
required to close the DOF question to the level a hostile referee
would demand.

\textbf{Dimensional bound on $\Delta^f_{\mu\nu}$ at perfect-fluid scope.}
The Lorentz-invariant phase-space measure
$d\Pi = d^3p/((2\pi)^3 p^0)$ depends on the metric through the
on-shell condition $p^0 = \sqrt{-g^{\mu\nu}p_\mu p_\nu}$. The metric
variation $\delta(d\Pi)/\delta g^{\mu\nu}$ at fixed $F$ contributes
to $\delta f/\delta g^{\mu\nu}$ a perfect-fluid combination
$\propto g_{\mu\nu} + u_\mu u_\nu$ for any isotropic $F$, with no
anisotropic-stress component. Its magnitude is bounded by
\begin{equation}
\label{eq:deltaf_bound}
|\Delta^f_{\mu\nu}|\;\lesssim\;\alpha\,(\sigma/c)^2\,(1{+}f)\,T_{\mu\nu},
\end{equation}
where $\sigma$ is the local matter velocity dispersion. For galactic
gas ($\sigma\sim 10$~km/s), $|\Delta^f|/[(1{+}f)T]\sim 10^{-8}$; for
cosmological dust ($\sigma\lesssim 600$~km/s), $\sim 10^{-5}$; for
recombination-era baryons, $\sim 10^{-9}$. The $d\Pi$-variation
contribution to $\Delta^f_{\mu\nu}$ is sub-leading at every
perfect-fluid scope by five or more orders of magnitude. The full
active-equilibration (Maxwell--J\"uttner) form, which adds the
contribution from $\delta F/\delta g$ at fixed thermodynamic state,
applies in early-universe tight-coupling regimes; for galactic and
present-epoch cosmology, the ratchet postulate
(Sec.~\ref{sec:ratchet}) gives $\Delta^f_{\mu\nu} = 0$ identically,
with \eqref{eq:deltaf_bound} bounding the deviation from the
ratchet even if the metric-measure dependence is retained.

\subsection{The information functional $f$}
\label{sec:f}

The matter coupling is set by the Kullback--Leibler divergence of
the matter distribution function $F(x,p)$ from a flat reference,
\begin{equation}
\label{eq:f}
f(x) \;=\; \alpha\!\int\! F(x,p)\,
\ln\!\left[\frac{F(x,p)}{F_{\rm flat}(p)}\right]\,d\Pi,
\end{equation}
with $d\Pi$ the Lorentz-invariant phase-space measure and $\alpha$ a
dimensionless coupling. The functional form is fixed by
Shore--Johnson axioms~\cite{Shore1980,Aczel1966,Jaynes1957}: invariance, subset
independence and system independence pick out the relative entropy
uniquely.

We adopt the local-bounded reading of $F_{\rm flat}$: the maximally
uninformative reference is uniform on the locally accessible phase
space ($v < v_{\rm esc}$ for a gravitationally bound system) rather
than on all of $\mathbb{R}^3$. With this reading, the entropy-matching
cutoff $v_{\rm max}^{\rm EM}=2.56\,\sigma_{\rm eff}$ coincides with
$v_{\rm esc}$ for any virialised population and forces $f_\star\approx 0$
structurally for hot virialised stars (\PROVEN, A05.OP1/OP2). Cold
gas at $\sigma\approx 8\,{\rm km/s}$ is highly concentrated within
its accessible phase space and carries
$f_{\rm gas}\approx \fprim$.

The env-dep-$f$ reading---in which
$\delta f \propto \delta\rho$ through the A05 D$_{\rm KL}$
operator---is required at the perturbation level
(Sec.~\ref{sec:theta}) and underlies the $K_{\rm env}=1.484$
environmental rescaling of the Hubble constant
(Sec.~\ref{sec:cosmology}); under this closure the strict
uniform-$f$ scope of the present derivation is the leading-order
treatment, and the $\delta f$ contribution to perturbations
restores the $\Lambda$CDM-like peak structure that strict
uniform-$f$ alone does not.

\subsection{Branch theorem}
\label{sec:branch}

Taking $\nabla^\mu$ of \eqref{eq:eom} and using the Bianchi identity
$\nabla^\mu G_{\mu\nu}=0$ together with the Noether constraint
$\nabla^\mu \Teff_{\mu\nu}=0$ yields the pointwise relation
\begin{equation}
\label{eq:branch}
R(x)\,\nabla_\nu\Psi(x) \;=\; 0\qquad \forall x.
\end{equation}%
\footnote{The algebraic identity $R\nabla_\nu\Psi = 0$ follows from
the Bianchi identity applied to any non-minimally coupled scalar
theory of the form $\Psi R$ in the gravity action with matter
satisfying the Noether identity. The structure is implicit in the
Palatini-class~\cite{Olmo2005} and metric scalar-tensor
literature~\cite{Sotiriou2010,FujiiMaeda2003}. ISST's contribution is
the physical interpretation: at every spacetime point either $R=0$
(Branch~A: cosmology, with non-trivial $\Psi$ transport on the
Wiltshire two-domain background) or $\nabla\Psi=0$ (Branch~B:
compact-object recovery of exact GR with effective coupling
$G_{\rm eff}=(1{+}f)/\Psi_0$, not a screening limit). The
dichotomous structure is exact, not asymptotic.}
At every spacetime point, either:
\begin{description}
\item[Branch A] ($R=0$, $\nabla\Psi\neq 0$): $R$ vanishes pointwise
while $\Psi$ varies. Cosmology lives here.
\item[Branch B] ($\nabla\Psi=0$, $R\neq 0$): $\Psi=\Psi_0$ constant;
\eqref{eq:eom} reduces to $\Psi_0 G_{\mu\nu}=8\pi T^{\rm eff}_{\mu\nu}$.
This is general relativity with $G=1/\Psi_0$.
\end{description}
\textbf{Branch selection rule.}
The Branch theorem~\eqref{eq:branch} admits two solution classes at
every spacetime point. Selection between them is determined by the
local equilibrium state of matter. Where matter is actively evolving
--- cosmological structure formation, ongoing gravitational
processing --- the matter sector's information functional $f$ is
non-stationary along worldlines, the trace transport
$\Box\Psi = (8\pi/3)(1{+}f)T$ has a sourced solution with
$\nabla\Psi \neq 0$, and consistency forces $R = 0$ (Branch~A).
Where matter has equilibrated --- compact objects, virialised
systems, post-recombination dust at the wall scale --- $f$ is
stationary on the dynamical timescale, the transport equation admits
$\Psi = \Psi_0$ constant, $R$ is unconstrained, and the field equation
reduces to exact Einstein gravity with
$G_{\rm eff} = (1{+}f)/\Psi_0$ (Branch~B). The static-fluid theorem
(F08, below) gives a specific dynamical realisation: static perfect
fluids satisfy $df/d\tau = 0$ and Branch~B is selected. The selection
rule connects the irreversibility of $f$ accumulation through
Standard-Model freezeouts and structure formation
(Sec.~\ref{sec:fprim}) to the geometric structure of the field
equation, placing ISST in the
Penrose~\cite{Penrose1979}, Jacobson~\cite{Jacobson1995},
Padmanabhan~\cite{Padmanabhan2010}, and Verlinde~\cite{Verlinde2011}
family of theories in which thermodynamic irreversibility shapes
gravitational structure; ISST's specific contribution is the explicit
local realisation through the matter-side information functional $f$
coupling to the gravitational sector via $(1{+}f)\mathcal L_m$, with
the static-fluid theorem (F08) providing the equilibrium-boundary
mechanism.

For any static spherically symmetric perfect fluid on Branch~A, the
central pressure is forced to $P_c=0$ (F05--F08): physical stars
cannot live on Branch~A and must self-select into Branch~B, where
their exterior is Schwarzschild~\cite{F03F08foundations}. Branch B
recovers GR \emph{exactly} in the strong-field regime, not as a
screening limit.

\textbf{Junction conditions across Branch~A / Branch~B interfaces.}
A virialised cluster on Branch~B is generically embedded in a
cosmological wall background on Branch~A; a static star on Branch~B
is embedded in the same. The metric $g_{\mu\nu}$ is required to be
$C^1$ across the matching surface $\Sigma$ separating the two
regions, while $R$ jumps from zero (Branch~A side) to a non-zero
value (Branch~B side) and $\nabla\Psi$ jumps from non-zero (A) to
zero (B). The transition is therefore not smooth in the
distributional sense: the field equation~\eqref{eq:eom} on $\Sigma$
admits an Israel-type junction reading~\cite{Israel1966}, with the
extrinsic-curvature jump $[K_{ij}]$ proportional to a localised
$\Psi$-gradient discontinuity. For all the applications in this
paper --- compact stars in cosmological dust, virialised clusters in
the void background, the Solar System in the Galactic wall ---
the matching surface lies in the weak-field regime where
$\Psi \to \Psi_0$ on the Branch-B side smoothly approaches the
slowly-varying $\Psi(t)$ of the Branch-A wall background, and the
junction reduces to the standard scalar-tensor matching of an
isolated source to a slowly-varying cosmological
background~\cite{DamourEspositoFarese1992}. The junction is therefore physically
benign in every regime relevant to this paper's predictions
(PPN, LLR, GW propagation, lensing). \emph{The full junction
analysis for strong-field interfaces ---  e.g., a Branch-B
neutron star in a strongly inhomogeneous Branch-A environment, or
a relativistic infall onto a Branch-B black hole through a
Branch-A medium --- is identified as \OPEN}, since the
distributional content of $R\nabla\Psi = 0$ at $\Sigma$ in
strong-field configurations has not been worked out.

The status of the branch theorem is \PROVEN; the junction analysis
is \DEMO\ in the weak-field regime and \OPEN\ in the strong-field
regime.

\subsection{$\fprim$ from Standard-Model freezeout}
\label{sec:fprim}

The matter coupling on cosmological scales is dominated by a
species-independent floor $\fprim$ set by the irreversible
information deposited into the baryon distribution at each
Standard-Model freezeout (electroweak, QCD hadronisation,
pion+muon decoupling, $e^+e^-$ annihilation):
\begin{equation}
\label{eq:fprim}
\fprim \;=\; \sum_{i\in{\rm freezeouts}}
\frac{\Delta g_{*,i}}{g_{*,{\rm after},i}} \;=\; 5.664.
\end{equation}
The formula uses no fitted parameters: every input is either a
Standard-Model degree-of-freedom count or a freezeout temperature
from QCD-lattice or perturbative-electroweak thermodynamics. The
formal chain from the action's KL functional (eq.~\ref{eq:f}) to
this discrete sum is structurally open (F145, four candidate paths
exhausted; see ``structural diagnosis'' paragraph below) ---
\eqref{eq:fprim} is therefore properly read as an SM-motivated
ansatz with no fitted parameters, not as a closed-chain derivation
from the action. The observational target is fixed by the Planck
primary parameters
$\Omega_b h^2 = 0.02237\pm 0.00015$ and $\Omega_m h^2 =
0.14301$~\cite{Planck2018}, giving
\begin{equation}
\frac{\Omega_m}{\Omega_b} - 1 \;=\; 5.40\pm 0.15\quad (68\%\,{\rm CL}).
\end{equation}
The DOF prediction $\fprim = 5.66\pm 0.06$ overshoots the observation
$5.40\pm 0.15$ by $4.2\%$ ($1.6$--$1.8\sigma$, where the lower
figure propagates the $\pm 0.06$ model uncertainty on $\fprim$);
equivalently, it predicts $\Omega_m = 0.329$ versus $0.3153\pm 0.0073$. We commit to the
DOF-counting value $\fprim = 5.664$ as the single ISST output and
report the residual as an honest limitation of the
sharp-transition approximation in the DOF sum (real freezeouts are
smooth crossovers; some species are partially decoupled at each
stage), \emph{not} as a fit to the observation. The alternative
practice of quoting $\fprim = 5.40$ would absorb the residual into
$\fprim$ and produce a misleading $0.1\%$ ``match''---we explicitly
do not do this.

A smooth-crossover evaluation of the DOF sum (F139) using realistic
transition widths (QCD: $\Delta T \sim 20$~MeV from lattice
thermodynamics; $e^+e^-$: $\Delta T \sim 0.3$~MeV) returns
$\fprim = 5.78$, marginally \emph{larger} than the sharp-transition
value, not smaller. The smooth integral $\int (1/g_*)\,|dg_*/dT|\,dT$
is path-independent and gives $\ln(g_{*,{\rm init}}/g_{*,{\rm final}})
= 3.31$ regardless of crossover width; a per-transition evaluation
that respects the discrete freezeout structure returns $5.78$.
The $1.6$--$1.8\sigma$ tension is therefore a genuine feature of
the DOF-counting prediction, not an artefact of the sharp-transition
approximation. An operator-level derivation from the full
$(1{+}f)\Lm$ thermal history, which would track the continuous
entropy injection without discretising into stages, remains
\OPEN\ (Sec.~\ref{sec:future}). Status: \DEMO\ (SM-motivated ansatz
with no fitted parameters and stated $1.6$--$1.8\sigma$ tension;
smooth-crossover refinement does not close the gap, F139; the
formal chain from the action's KL functional to the discrete sum
is structurally open, F145).

\textbf{From KL functional to DOF sum: structural diagnosis.}
The cosmological-floor $\fprim$ is intended as the value of $f$
(eq.~\ref{eq:f}) at horizon scale on the post-recombination baryon
distribution. The discrete formula \eqref{eq:fprim} is
\emph{motivated} by the Acz\'el--Shore--Johnson uniqueness of
$D_{\rm KL}$ applied to the Standard-Model thermal history: each
freezeout that lowers $g_{*,s}$ injects an irreversible information
increment $\sim \Delta g_{*}/g_{*,{\rm after}}$ into the
photon-baryon plasma, which the ratchet postulate
(Sec.~\ref{sec:ratchet}) then preserves into the baryon distribution.

\emph{The chain from this motivation to the discrete formula does
not close as a derivation.} F145 tested four candidate paths from
the action's KL functional to the discrete sum
(Table~\ref{tab:fprim_paths}) and ruled them all out: each gives
the wrong magnitude, the wrong scope, or the wrong information
measure to recover \eqref{eq:fprim} from the action with the
galactic-calibrated $\alpha=0.869$ (F119).

\begin{table}[h]
\caption{Four candidate paths from the action's KL functional
$f = \alpha\!\int F\ln(F/F_{\rm flat})\,d\Pi$ to the discrete DOF
sum $\fprim = 5.664$, all ruled out (F145).}
\label{tab:fprim_paths}
\begin{ruledtabular}
\begin{tabular}{lcl}
Path & Result & Verdict \\
\hline
$D_{\rm KL}$(Maxw. @$T_{\rm actual}$ vs.\ @$T_{\rm ref}$) & $1.36$ & factor $\sim\!4$ short \\
$D_{\rm KL}$(Maxw. vs.\ uniform-on-$|p|<m_p c$) & ${\sim}40$--$60$ & factor $\sim\!8$ too large; galactic-specific reading \\
$\sum_i D_{\rm KL}^{(i)}$ (per-step transient, 4 freezeouts) & $0.30$ & factor $\sim\!19$ short; SJ additivity is over independent components, not events \\
$\sum_i D_q^{(i)}$ R\'enyi at $q=2$ (4 freezeouts) & $5.82$ & accidental $\sim\!3\%$ match; Shore--Johnson uniqueness forces $q=1$, and the match degrades away from SM step sizes \\
\end{tabular}
\end{ruledtabular}
\end{table}

The smooth path-integral form
\begin{equation}
\label{eq:fprim_integral}
\fprim^{\rm smooth} \;=\; \alpha_{\rm cosmo}\!\int_{T_{\rm EW}}^{T_{\rm rec}}\!\frac{1}{g_{*,s}}\bigg|\frac{dg_{*,s}}{dT}\bigg|\,dT
\;=\; \alpha_{\rm cosmo}\,\ln\!\bigg(\frac{g_{*,s}(T_{\rm EW})}{g_{*,s}(T_{\rm rec})}\bigg)
\;=\; 3.31\,\alpha_{\rm cosmo}
\end{equation}
between $g_{*,s}(T_{\rm EW}) = 106.75$ and $g_{*,s}(T_{\rm rec}) = 3.91$
gives $\ln(106.75/3.91) = 3.31$ path-independently (F139). Matching
this to $\fprim = 5.66$ requires $\alpha_{\rm cosmo} \approx 1.71$
--- distinct from the galactic calibration $\alpha = 0.869$ from
F119 used in the local-bounded reading
(Sec.~\ref{sec:env_dep_f_closure}). Whether the two $\alpha$ values
share a common origin in a single bounds-dependent functional is a
propagated open question (F145).

\emph{Staging dependence.} The discrete sum \eqref{eq:fprim} is a
right-endpoint Riemann sum on a coarse partition. With $N$-step
uniform-log sub-partitions of $[g_{*,{\rm final}},\,g_{*,{\rm init}}]$
the right-endpoint sum decreases monotonically: $26.3$ ($N{=}1$),
$5.14$ ($N{=}4$, uniform-log), $3.31$ ($N\to\infty$). \emph{Refining
the partition does not converge to $5.664$ from above; it diverges
away from $5.664$ toward the integral limit $3.31$.} The paper's
non-uniform 4-event partition (the SM mass thresholds) gives $5.66$
because the Standard Model has four freezeouts at those particular
$g_{*,s}$ values; the integer $4$ is a contingent feature of the SM
mass spectrum, not a derived structural number. A
beyond-Standard-Model extension that adds or removes mass thresholds
would shift the prediction.

\emph{Structural obstruction.} The action's $f$ in eq.~\ref{eq:f} is
a \emph{single-epoch} $D_{\rm KL}$ integral on a single distribution
at a single time. The discrete sum \eqref{eq:fprim} is a
\emph{path-additive} event-counting quantity that depends on the
partition into freezeout events. These are different mathematical
objects. Shore--Johnson uniqueness (additivity over independent
components, not over time-separated events) gives the single-epoch
form; it does not produce path additivity across cosmic history.
\emph{The discrete formula \eqref{eq:fprim} is therefore not a
Riemann-sum approximation of \eqref{eq:fprim_integral}}, nor of any
single-epoch $D_{\rm KL}$. Closing the chain requires either a
path-integral extension of the action (which would re-check A05,
the ratchet postulate F46, $(1{+}f)\mathcal L_m$, and the
Shore--Johnson uniqueness derivation) or a demonstration that the
discrete sum emerges from a different limit of $(1{+}f)\mathcal L_m$
than the single-epoch $D_{\rm KL}$. We commit to the discrete
formula \eqref{eq:fprim} as the working ansatz; its theoretical
justification from the action remains \OPEN\ as the operator-level
derivation flagged in Sec.~\ref{sec:future} (F145).

\textbf{Uncertainty on $\fprim$.}
The discrete-vs-smooth-crossover spread is $\fprim^{\rm sharp} = 5.664$
vs.\ $\fprim^{\rm smooth} = 5.78$, an $\approx 2\%$ method
uncertainty. Lattice-QCD uncertainty on $g_{*,s}(T_{\rm QCD})$ in
the crossover region~\cite{BorsanyiQCD2016} contributes a further
$\sim 1$--$2\%$. Combining these in quadrature gives
$\fprim = 5.66 \pm 0.06$ at minimum (method spread) and
$\fprim = 5.66 \pm 0.12$ if lattice systematics are propagated. The
tension with the Planck value $\Omega_m/\Omega_b - 1 = 5.40\pm 0.15$
is $1.6$--$1.7\sigma$ when method+lattice uncertainty is propagated
(equivalent to $1.8\sigma$ if only the Planck error is used,
ignoring model uncertainty). We continue to quote the more
conservative $1.6$--$1.7\sigma$ figure as the honest tension level.

\textbf{Neutrino sector contribution.} Neutrino decoupling at
$T \sim 1$~MeV does not by itself change $g_{*,s}$ in the integrand
of \eqref{eq:fprim_integral}: the decoupled neutrinos retain
$T_\nu = T_\gamma$ until the next event. The neutrino sector enters
$\fprim$ through the subsequent $e^+e^-$ annihilation at
$T \sim 0.3$~MeV, at which the photon-coupled effective DOF drops
sharply while the decoupled neutrinos are not reheated --- yielding
the famous $T_\gamma/T_\nu = (11/4)^{1/3}$ ratio. The $e^+e^-$
entropy injection reaches the baryon distribution through the
photon-baryon tight coupling that persists until recombination
($z\sim 1090$); the ratchet postulate (below) then ensures the
information increase is permanently recorded in the baryon
distribution and is not undone by subsequent dynamics. The DOF sum
\eqref{eq:fprim} therefore tracks the cumulative photon-coupled-plasma
DOF history, not a direct neutrino-baryon interaction.

\subsection{The ratchet postulate}
\label{sec:ratchet}

Several results in this paper rely on the assumption that the
information functional $f$ is irreversibly accumulated through cosmic
history. We state this explicitly:

\textbf{Postulate (Ratchet).} \emph{Along any matter worldline,
the information functional $f$ is monotonically non-decreasing:
$df/d\tau \geq 0$.}

\emph{Physical basis.} The $D_{\rm KL}$ operator measures departure
from maximum entropy (the flat reference $F_{\rm flat}$).
Gravitational structure formation increases this departure
($df/d\tau > 0$ during collapse and virialisation), while local
thermalisation in the matter sector reduces it locally but increases
total entropy --- the second law applied to the structured
component of the matter distribution. The ratchet is the explicit
statement that information \emph{increments} produced by
Standard-Model freezeouts (Sec.~\ref{sec:fprim}) and by structure
formation are not erased by subsequent equilibration: once a baryon
distribution carries $f$, that contribution persists.

\emph{Where the ratchet is invoked in this paper.} (i)~The DOF sum
\eqref{eq:fprim} is path-additive across SM freezeouts only because
each increment is preserved through the next stage. (ii)~The
``lighthouse'' Bullet/A520 mechanism (Sec.~\ref{sec:bullet})
distinguishes shock-thermalised gas (which loses its working memory
$f_s$ but retains $\fprim$) from intact stars (which retain both):
the ratchet allows local thermalisation to wipe the structure-formation
contribution to $f_s$ without unwinding $\fprim$. (iii)~The
$\Delta^f_{\mu\nu} = 0$ identity used at perfect-fluid scope
(Sec.~\ref{sec:action}) is the ratchet's metric-variation
consequence: under the ratchet, the metric-variation derivative of
the Maxwell--J\"uttner functional vanishes identically, leaving only
the dimensional bound \eqref{eq:deltaf_bound} from the
$d\Pi$-measure.

The ratchet is a postulate, not a theorem of F01: it is consistent
with the action but is not derivable from it without further
specification of the matter sector dynamics under coarse-graining.
A full operator-level derivation that closes the chain
$\text{action} \to \text{ratchet} \to \fprim$ remains open
(Sec.~\ref{sec:future}).

\subsection{The galactic-vs-cosmological coupling: a factor-2 scale dependence}
\label{sec:alpha_scale}

The KL coupling $\alpha$ in \eqref{eq:f} is fitted at galactic scale
to $\alpha_{\rm gal} = 0.869$ (the calibration procedure is described
in Sec.~\ref{sec:acrit}: minimise the median rotation-curve RMS
over the 14-galaxy SPARC subsample). Recovering the cosmological
floor $\fprim = 5.664$ from the smooth path-integral form
\eqref{eq:fprim_integral}, however, requires
$\alpha_{\rm cosmo} \approx 1.71$ --- a factor-$\sim 2$ difference
between scales.

\emph{This running has not been derived from the action.} It may
reflect genuine scale-dependent physics: the $D_{\rm KL}$ operator
probes different aspects of the phase-space distribution at
galactic versus cosmological scales (the local-bounded reading at
galactic scale uses $F_{\rm flat}$ uniform on $v < v_{\rm esc}$;
the cosmological reading would use a horizon-bounded reference
that has not been formally constructed). Alternatively, the
factor-$\sim 2$ may indicate that the discrete sum and the
continuous integral measure different quantities --- the discrete
sum being a path-additive event count over freezeout stages, the
continuous integral being a single-epoch $D_{\rm KL}$, with a
scale-dependent normalisation between the two (see ``Structural
obstruction'' above and Sec.~\ref{sec:fprim}).

Until the operator-level derivation of $\fprim$ from $(1{+}f)\Lm$
on the SM thermal history is closed --- including a unified
treatment of $\alpha$ at galactic and cosmological scales --- this
factor-$\sim 2$ scale dependence is an \emph{acknowledged open
tension in the theory}. It does not invalidate the galactic
calibration (which is internally self-consistent for rotation-curve
phenomenology) or the cosmological prediction (which is internally
self-consistent for the SM thermal-history input), but it does
mean that $\alpha$ is not, at present, a single parameter
universally calibrated across all scales. Status: \OPEN.

\subsection{Degree-of-freedom count}

ADM 3+1 decomposition of \eqref{eq:eom} gives two propagating tensor
degrees of freedom and no scalar mode (App.~\ref{app:adm_dof};
framework \textsf{adm\_decomposition} \textsc{[allowed]}). The
scalar $\Psi$ is a constraint field with first-class
$(\pi_\Psi, R=0)$ pair generating the gauge transformation
$\delta\Psi = \varepsilon$ on the constraint surface; the theory
sits in the surviving Horndeski corner that GW170817 left intact.
Status: \PROVEN\ internally; external verification against the
mimetic-DOF literature
(\cite{Barvinsky2014,Capela2015,Firouzjahi2017,Gorji2018}) flagged
in Sec.~\ref{sec:future}.

\subsection{Environment-dependent $f$ and the perturbation closure}
\label{sec:env_dep_f_closure}

\textbf{Headline.} \emph{The env-dep-$f$ closure derived below
(F138) is a few-percent multiplicative correction to the
gravitational source, not a free-growing collisionless channel.}
It does not structurally rescue the factor-$\sim\!30$ third-peak
suppression of strict uniform-$f$ that F134 documents
(Sec.~\ref{sec:theta}); a first-pass linear-scaling estimate
leaves a residual gap of $\sim\!30$ relative to the Planck-observed
$H_3/H_1$ in either of the two readings of $v_{\rm esc}(\delta)$
considered. The CAMB Path~A treatment of $(1{+}\fprim)\Omega_b$
as collisionless CDM at the perturbation level is therefore an
\emph{empirical approximation} rather than a derived consequence of
F01: F138 establishes that the env-dep-$f$ closure does not derive
the proxy, only that it bounds the ISST-internal correction at the
$\pm 16\%$ level. Two resolution paths remain: a $V(\Psi)$
extension that decouples $\Psi$-mediated structure from photon
dynamics (action-level extension beyond F01), or the honest-flag
position adopted in Sec.~\ref{sec:theta}. The remainder of this
subsection derives the closure form, expands it to linear order,
gives the two readings of $v_{\rm esc}(\delta)$, and quotes the
$\beta(z)$ values from F138.

\textbf{The closure form.}
The information functional from A05, under the F119 local-bounded
reading (with $F_{\rm flat}$ uniform on velocities $v < v_{\rm esc}(x)$
locally accessible from position $x$), evaluates for non-relativistic
Maxwellian matter to
\begin{equation}
\label{eq:DKL_evaluated}
\dKL/n \;=\; 3\ln(v_{\rm esc}/\sigma) - 2.825 + v_{\rm bulk}^2/(2\sigma^2),
\end{equation}
with $\sigma^2 = T/m_p$ and the constant $-2.825 = \ln(4\pi/3) -
(3/2)\ln(2\pi e)$ absorbing volume normalisations. At background
($v_{\rm bulk} = 0$), $\bar f = \alpha \cdot [3\ln(\bar v_{\rm esc}/\bar\sigma)
- 2.825]$, calibrated to $\fprim = 5.664$ by A04.

\textbf{Linear perturbation expansion.}
To first order in $\delta_b \equiv \delta\rho_b/\bar\rho_b$, dropping
$O(\delta^2)$ and noting that the linear-in-$v_{\rm bulk}$ contribution
vanishes at background:
\begin{equation}
\label{eq:delta_f_linear}
\delta f(x,t) \;=\; \beta(z)\,\delta_b(x,t) \;+\; O(\delta^2),
\end{equation}
with the coefficient $\beta(z)$ determined by the linear response of
$\ln v_{\rm esc}$ and $\ln\sigma$ to $\delta_b$:
\begin{equation}
\label{eq:beta_def}
\beta(z) \;=\; \alpha \cdot [3\,\delta\ln v_{\rm esc} - 3\,\delta\ln\sigma]/\delta_b.
\end{equation}

\textbf{Two readings of $v_{\rm esc}(\delta)$.}
The local-bounded $v_{\rm esc}$ is set by the local gravitational
potential, $v_{\rm esc}^2 = -2\Phi$. Its perturbative response depends
on what scale sets $\bar v_{\rm esc}$:
\begin{itemize}
\item \emph{Reading B (local-collapse).} For a perturbation at
scale $1/k$, $\bar v_{\rm esc}^2 \propto G\bar\rho_b(1{+}\bar f)/k^2$;
the self-consistent solve gives $\delta\ln v_{\rm esc} = (1/2)[1 +
\beta/(1{+}\bar f)]\delta_b$. This is the convention under which
the F119 galactic calibration $\alpha = 0.869$ was performed.
\item \emph{Reading A (horizon-scale, relativistic).} The maximal
binding scale is the cosmic horizon, $\bar v_{\rm esc} \approx c$.
For sub-horizon CMB modes ($k \gg H/c$), $\delta\ln v_{\rm esc}$ is
suppressed by $(H/k)^2$ and is essentially zero.
\end{itemize}
Which reading applies at cosmological linear-perturbation scales is
itself an open question; the F119 derivation directly calibrates
only the galactic regime.

\textbf{The $\sigma$ response.}
For pre-decoupling photon-coupled baryons, $T_b = T_\gamma \propto
\rho_\gamma^{1/4}$ adiabatically, $\delta\ln\sigma = (1/6)\delta_b$.
For post-decoupling free baryons with adiabatic index $\gamma_b = 5/3$,
$T_b \propto \rho_b^{2/3}$, $\delta\ln\sigma = (1/3)\delta_b$.

\textbf{Numerical $\beta(z)$ at $\alpha = 0.869$, $\bar f = 5.664$
(F138 \texttt{F138\_class\_estimate.py}):}
\begin{equation}
\label{eq:beta_values}
\begin{array}{l|cc}
& \beta_{\rm pre-dec} & \beta_{\rm post-dec} \\
\hline
\text{Reading B (local-collapse):} & +1.08 & +0.54 \\
\text{Reading A (horizon-scale):} & -0.43 & -0.87
\end{array}
\end{equation}

\textbf{The modified gravitational source.}
The Poisson source acquires
\begin{equation}
\label{eq:env_dep_source}
\nabla^2 \delta\Phi \;=\; 4\pi G\,\bar\rho_b\,
[(1{+}\bar f) + \beta(z)]\,\delta_b,
\end{equation}
giving an enhancement of $+8\%$ to $+16\%$ over CAMB Path A's
$(1{+}\fprim) = 6.664$ baseline under Reading B, and a suppression
of $-6\%$ to $-13\%$ under Reading A.

\textbf{Why the closure does not rescue the third peak.}
$\delta f$ is a linear functional of $\delta_b$ (and of $\delta\sigma$,
itself a function of $\delta_b$ through the local thermodynamic
state), so it inherits the photon-coupled oscillations of $\delta_b$
pre-decoupling exactly: there is no free-growing collisionless
channel that decouples from photon-baryon dynamics during radiation
era. The factor-$33$ third-peak suppression of strict uniform-$f$
(F134, modified CLASS) is therefore not rescued within F01 by this
closure alone. Status: \PROVEN\ for the closure form
\eqref{eq:delta_f_linear}--\eqref{eq:env_dep_source} and the
$\beta(z)$ values \eqref{eq:beta_values}; \DEMO\ for the
leading-order suppression of $\Delta^f_{\mu\nu}$ pending explicit
evaluation of $\bar\Lambda$ (F138 §8); \OPEN\ for a derivation of
the CAMB Path~A proxy from F01 alone, which would close the F134
falsification within the committed action.

%==========================================================
\section{Solar System and local tests}
\label{sec:solar}

\subsection{Post-Newtonian parameters}

In the Solar System, matter is concentrated, static, and well
described by the Branch-B reduction $\Psi=\Psi_0$, $G=1/\Psi_0$. By
the branch theorem, the field equation is pointwise GR. We therefore
have
\begin{equation}
\gamma_{\rm PPN} = 1,\qquad \beta_{\rm PPN} = 1,
\end{equation}
\emph{exactly}, by two cross-checking routes that share the
Branch-B selection (App.~\ref{app:static_fluid_theorem}) but verify
the PPN coefficients independently:
(i)~directly from the Branch-B field equation
$\Psi_0 G_{\mu\nu} = 8\pi(1{+}f)T_{\mu\nu}$ giving exact GR with
$G_{\rm eff} = (1{+}f)/\Psi_0$, hence $\gamma = \beta = 1$; and
(ii)~from the Palatini-class PPN analysis (framework
\textsf{palatini\_ppn} \textsc{[allowed]}). In the Palatini reading
(\cite{Olmo2005,Will2014}), the $\Psi R$ sector with no kinetic
term gives $\gamma = \beta = 1$ identically because the metric-sector
PPN parameters depend only on the matter coupling: the field
equation in Palatini form is $G_{\mu\nu}[\hat g] = 8\pi T_{\mu\nu}/\Psi$
in the auxiliary metric $\hat g$, and the conformal factor relating
$\hat g$ to $g$ is determined algebraically by the trace, leaving
the standard PPN expansion of GR with a coupling rescaling. The
Branch-B route delivers the same answer by the different mechanism
$\nabla\Psi = 0$ inside the source, which reduces \eqref{eq:eom}
to GR. \emph{Both routes presume Branch~B inside the source, which
presumes the static-fluid theorem; they are not independent of
that selection, only of each other in the PPN evaluation.} ISST's structural
non-propagation of $\Psi$ distinguishes it from
Damour--Esposito-Farese tensor-multi-scalar PPN, which requires a
propagating scalar~\cite{DamourEspositoFarese1992}. Cassini's
$|\gamma-1|<2.3\times 10^{-5}$~\cite{Bertotti2003} is satisfied at
leading order. Status: \PROVEN. Sub-leading corrections of order
$f_s^{\odot}$ (the local stellar information content) are unmeasured;
this is identified as an \OPEN\ refinement (Sec.~\ref{sec:future}).

\textbf{Extension to rotating compact sources.}
The Branch-B selection extends to stationary axisymmetric
configurations. At first order in the rotation rate $\Omega$
(slow-rotation Hartle--Thorne expansion~\cite{Hartle1967}), the
$(rr)$, $(\theta\theta)$, $(tt)$, $(\phi\phi)$ projections of the F01
field equation are even in $\Omega$ and unchanged from the static
case at $O(\Omega^0)$; the F08 algebra forces $\nabla\Psi = 0$ and
$\Psi = \Psi_0$ throughout the fluid. The new $O(\Omega)$ projection
is the $(t\phi)$ frame-dragging equation,
$\Psi_0 G_{t\phi} = 8\pi(1{+}f_\star)T_{t\phi}$, which determines
the standard Hartle--Thorne function $\omega(r,\theta)$ with
$G_{\rm eff} = (1{+}f_\star)/\Psi_0$ and does not constrain $\Psi$.
This covers all astrophysical neutron stars: even the fastest known
pulsar PSR~J1748$-$2446ad at $716$~Hz has
$\Omega/\Omega_K \approx 0.43$, with $O(\Omega^2)$ corrections
bounded at $\sim 18\%$. For arbitrary rotation rate (including
Kepler-frequency configurations) and for differentially rotating
equilibrium, the branch-selection rule (Sec.~\ref{sec:branch})
selects Branch~B by stationary-equilibrium ($df/d\tau = 0$) of the
matter in its rest frame, locally pointwise for differential
rotation. Vacuum Kerr black holes satisfy
$R_{\mu\nu} = 0 \Rightarrow R = 0$ and admit either branch trivially
in the metric equation; the global $\Psi = \Psi_0$ value is set by
junction continuity to the surrounding matter in Branch~B. Branch B
therefore holds for all astrophysically relevant compact-object
configurations including rotating neutron stars, white dwarfs, and
stellar-mass and supermassive black holes; the GR-recovery and PPN
$\gamma = \beta = 1$ results extend without modification.

\subsection{Fifth force and equivalence principle}
\label{sec:fifth}

There is no \emph{propagating-scalar} fifth force in ISST. The scalar
$\Psi$ does not propagate (no kinetic term in \eqref{eq:action}; F128
ADM count gives zero scalar degrees of freedom). The Derivation
Passport explicitly denies all standard screening mechanisms
(chameleon, symmetron, Vainshtein, Yukawa) as inapplicable to F01:
their structural predicate is a propagating scalar, which ISST does
not have. The universal $(1+f)\Lm$ coupling implies the weak
equivalence principle at the action level. The $(1+f)$ factor
appearing in the modified Poisson equation \eqref{eq:poisson} is a
matter-side rescaling of the gravitational charge, not a fifth-force
mediator: it has no propagator, no finite range, and does not
introduce a new force-carrying particle. ISST therefore occupies
neither the Yukawa-type fifth-force category nor the
Damour--Esposito-Farese tensor-multi-scalar category; it is in the
$\Psi R$ Palatini-class corner with non-propagating $\Psi$.
Status: \PROVEN.

\subsection{Lunar Laser Ranging $\dot G/G$}
\label{sec:llr}

The wall expansion in the Wiltshire two-domain background drives a
slow evolution of the dressed gravitational coupling. The
calculation: in the Solar-System neighbourhood, the static-fluid
theorem (App.~\ref{app:static_fluid_theorem}) selects Branch~B, so
$\Psi(x,t) = \Psi_0$ is constant throughout the local bound region
to the precision of the Branch-B reduction; any time-dependence in
$G_{\rm eff} = (1+f)/\Psi_0$ comes from boundary-condition drag of
the cosmological wall background on the local Branch-B value. With
the wall background's $\dot\Psi/(H\Psi)|_0 = (\sqrt{21}-3)/2$
attractor (Sec.~\ref{sec:acrit}, Eq.~\ref{eq:p_over_n}) and a
geometric suppression factor $\varepsilon \sim 10^{-7}$ from the
ratio of the Solar-System scale to the wall scale (the local
boundary value $\Psi_{\rm wall}(t)$ leaks into the Solar System on
the time-scale set by $L_{\rm wall}/c$, with the suppression
controlled by the ratio of the bound-system size to the wall scale
under the static-source Poisson kernel), the rate is
\begin{equation}
\left|\dot G/G\right|^{\rm ISST}_{\rm wall}
\;\sim\; \frac{p}{n}\,H_{\rm wall}\,\varepsilon
\;=\; 5.63\times 10^{-18}\,{\rm yr}^{-1},
\end{equation}
five orders of magnitude below the Lunar Laser Ranging bound
$7\times 10^{-13}\,{\rm yr}^{-1}$~\cite{Hofmann2018,Williams2004}.
An FRW evaluation, where Branch~A propagates $\Psi$-evolution into
the Solar-System region without the Branch-B switch, gives
$\sim 5\times 10^{-11}\,{\rm yr}^{-1}$, two orders \emph{above}
LLR; the difference is a direct measure of the Wiltshire
background's necessity, not a free parameter, and a sharp
discriminator for the static-fluid theorem's role in localising
$\Psi$ at compact-source scales. Status: \PROVEN.

%==========================================================
\section{Galaxy-scale phenomenology}
\label{sec:galaxy}

\subsection{Rotation curves}

In the Branch-A weak-field regime, the modified Poisson equation
takes the form
\begin{equation}
\label{eq:poisson}
\nabla^2\Phi \;=\; \frac{4}{3}\cdot 4\pi G_N\,(1+f)\,\rho
\;=\; \frac{16\pi G_N}{3}\,(1+f)\,\rho
\quad\text{(rotation-supported disks)},
\end{equation}
where the prefactor $4/3$ multiplies the standard Newtonian
$4\pi G_N$ coupling and is the unique value compatible with the
Branch-A constraint $R=0$ on the rotation-supported background.
The isotropic-gauge $2/3$ used in an earlier draft was
incompatible with Branch~A and has been corrected; the explicit
weak-field derivation reproducing $4/3$ is given in
Appendix~\ref{app:Cgrad} alongside the lensing $2/3$ and
scalar-transport $2/3$ coefficients. On top of
\eqref{eq:poisson} we apply a Newton-limit regulator
$\mathcal S(a_{\rm bar}) = 1/[1+(a_{\rm bar}/a_{\rm crit})^2]$ that
damps the ISST correction at high acceleration:
\begin{equation}
V^2(R) = V_{\rm bar}^2(R)\,
\bigl[\,1+(E(R)-1)\,\mathcal S(a_{\rm bar})\,\bigr],
\end{equation}
with $E = \tfrac{4}{3}\bigl[1+\fprim(a_0/a_{\rm bar})^q\bigr]$.

\subsection{Acceleration scales: one derived, one empirical}
\label{sec:acrit}

The two acceleration scales in the regulator have asymmetric status.

\textbf{$a_{\rm crit}$ is derived in two stages, both from the
committed action.} A full derivation is given in
Appendix~\ref{app:acrit_derivation}; we summarise the load-bearing
algebra here.

\emph{Stage 1 (linearised perturbations on Minkowski).} The
two-potential weak-field metric
$ds^2 = -(1+2\Phi)\,dt^2 + (1-2\Phi_N)\,\delta_{ij}\,dx^i dx^j$ with
$\Psi = \Psi_0 + \delta\Psi$ and $T = -\rho$ (dust), substituted into
\eqref{eq:eom} and linearised, yields three independent Laplacian
equations and the Branch-A constraint:
\begin{align}
\label{eq:wf_00}
2\Psi_0\nabla^2\Phi_N - \nabla^2\delta\Psi &= 8\pi(1+f)\rho, \\
\label{eq:wf_trace}
2\Psi_0\nabla^2\Phi - 4\Psi_0\nabla^2\Phi_N + 3\nabla^2\delta\Psi
   &= -8\pi(1+f)\rho, \\
\label{eq:wf_offdiag}
\Psi_0(\nabla^2\Phi_N - \nabla^2\Phi) &= \nabla^2\delta\Psi, \\
\label{eq:wf_branchA}
R = -2\nabla^2\Phi + 4\nabla^2\Phi_N &= 0.
\end{align}
Solving \eqref{eq:wf_00}--\eqref{eq:wf_branchA} simultaneously gives
\begin{equation}
\nabla^2\Phi = \tfrac{4}{3}\,4\pi G_N(1{+}f)\rho,\quad
\nabla^2\Phi_N = \tfrac{2}{3}\,4\pi G_N(1{+}f)\rho,\quad
\nabla^2\!\bigl(\tfrac{\delta\Psi}{\Psi_0}\bigr) =
   -\tfrac{2}{3}\,4\pi G_N(1{+}f)\rho.
\end{equation}
For a point mass $M$ the gradient
$|\nabla(\delta\Psi/\Psi_0)| = (2/3)(1{+}f)\,a_{\rm bar}/c^2$, so the
\emph{stability-gradient coefficient} is uniquely
\begin{equation}
\label{eq:Cgrad}
C_{\rm grad} = \tfrac{2}{3},
\end{equation}
gauge-invariant (Bardeen) under the choice of slicing. The other
candidate values $4/3$ (rotation curves), $2$ (lensing sum), and $1$
all fail the off-diagonal constraint \eqref{eq:wf_offdiag}.

\emph{Stage 2 (matter-era power-law on the wall background).} On the
wall Friedmann \eqref{eq:wall_friedmann_main}
$H^2 = (16\pi/9)(1{+}f_{\rm prim})\rho/\Psi$ coupled to
$\Box\Psi = -(8\pi/3)(1{+}f_{\rm prim})\rho$ with continuity
$\dot\rho + 3H\rho = 0$, the power-law ansatz
$a\propto t^n$, $\Psi\propto t^p$, $\rho\propto t^{-3n}$ yields a
\emph{forced} algebraic system. Power matching gives $p = 2-3n$;
substituting into the $\Psi$-transport amplitude balance
$p(p-1+3n)\Psi_0 = (8\pi/3)(1{+}f_{\rm prim})\rho_0$ and using the
Friedmann amplitude $n^2 = (16\pi/9)(1{+}f_{\rm prim})\rho_0/\Psi_0$
reduces the system to the quadratic
\begin{equation}
\label{eq:n_quadratic}
3n^2 + 6n - 4 = 0,
\end{equation}
whose physical root is $n = (\sqrt{21}-3)/3 \approx 0.528$, giving
\begin{equation}
\label{eq:p_over_n}
\frac{\dot\Psi}{H\Psi}\bigg|_0 = \frac{p}{n} = \frac{\sqrt{21}-3}{2}
\approx 0.7913.
\end{equation}
\emph{Note on units:} \eqref{eq:p_over_n} is the dimensionless
matter-era attractor of $\dot\Psi/(H\Psi)$, \emph{not} the
critical acceleration. The full $a_{\rm crit}$ formula
\eqref{eq:acrit_main} below combines this $p/n$ with $C_{\rm grad}^{-1}
= 3/2$ from \eqref{eq:Cgrad} and the $(1+f_{\rm prim})^{-1} \approx
0.150$ matter-coupling divisor; the algebraic prefactor on $cH_0$ is
therefore $(3/2)\times(\sqrt{21}-3)/2 / (1+f_{\rm prim}) =
3(\sqrt{21}-3)/[4(1+f_{\rm prim})] \approx 0.178$, evaluating to
$1.07\times 10^{-10}$ m/s$^2$ at $H_0 = 61.79$~km/s/Mpc.
This is the unique attractor of the coupled system; the alternative
forms B (cross-term retained) and B$'$ (BD $\omega=0$) yield
$\sqrt 3$ and $1$ respectively, both ruled out empirically against
the SPARC critical scale (Appendix~\ref{app:acrit_derivation}).

\emph{Combining stages.} The stability-gradient crossover equates the
local Branch-A galactic gradient
$C_{\rm grad}(1{+}f_{\rm prim})\,a_{\rm bar}/c^2$ to the cosmological
floor $H_0^2 L_{\rm eff}/c^2$ with
$L_{\rm eff} \equiv (p/n)\,c/H_0$, giving
\begin{equation}
\label{eq:acrit_main}
\boxed{\;
a_{\rm crit}^{\rm ISST} \;=\;
\frac{1}{C_{\rm grad}}\,\frac{p}{n}\,
\frac{cH_0}{1+f_{\rm prim}}
\;=\; \frac{3(\sqrt{21}-3)}{4}\,\frac{cH_0}{1+f_{\rm prim}}
\;\approx\; 1.069\times 10^{-10}\,{\rm m/s^2}
\;}
\end{equation}
at $H_0 = 61.79$~km/s/Mpc and $f_{\rm prim} = 5.66$, sitting
$10.9\%$ below the MOND value $a_0 = 1.20\times 10^{-10}$~m/s$^2$.
The result is parameter-free at fixed $H_0$ and $\fprim$; every
factor on the RHS of \eqref{eq:acrit_main} is ISST-committed
($C_{\rm grad}$ from \eqref{eq:Cgrad}; $p/n$ from
\eqref{eq:p_over_n}; $f_{\rm prim}$ from
Sec.~\ref{sec:fprim}; $H_0$ from Sec.~\ref{sec:hubble}).
\emph{Sensitivity to $H_0$ choice:} at Planck $H_0 = 67.4$~km/s/Mpc
the prediction is $1.166\times 10^{-10}$~m/s$^2$ ($-2.8\%$ from
MOND); at SH0ES $H_0 = 73.0$, $1.263\times 10^{-10}$~m/s$^2$
($+5.3\%$). The ``$10.9\%$ below MOND'' figure is therefore
specific to the ISST-committed dressed value $H_0 = 61.79$; a
falsification claim against an a$_0$ measurement requires
committing to the same $H_0$. Status: \PROVEN\ at the committed
$(H_0, \fprim)$; \DEMO\ as a single-number prediction independent
of the $H_0$ choice.

\textbf{$a_0$ (regulator internal scale) is empirical.} The scale
$a_0 = 0.01\times 10^{-10}\,{\rm m/s^2}$ in the regulator is fitted
on a 14-galaxy SPARC subsample; the regulator exponent $q\approx
0.45$ is also empirical at present. Their derivation from
second-order Branch-A perturbation theory is identified as \OPEN.

\subsection{Full SPARC sample}

Evaluated on \emph{every galaxy in the SPARC database}
(175 galaxies, no per-galaxy retuning, no quality cuts):
\begin{align*}
\langle{\rm RMS}\rangle_{\rm ISST}^{N=175} &= 13.8\,{\rm km/s},\\
\langle{\rm RMS}\rangle_{\rm MOND}^{N=175} &= 12.7\,{\rm km/s},
\end{align*}
with ISST winning per-galaxy on $86/175$ ($49\%$). The MOND comparison
uses the simple interpolating function
$\mu(x)=x/(1+x)$~\cite{FamaeyMcGaugh2012} with
$a_0 = 1.2\times 10^{-10}\,{\rm m/s}^2$; the standard interpolating
function $\mu(x)=x/\sqrt{1+x^2}$ shifts the MOND median by
$\lesssim 1$~km/s and does not affect the qualitative comparison.
On the high-quality $Q{=}1$ subsample of 99 galaxies (Lelli et al.'s
top flag~\cite{Lelli2016}), ISST wins the median ($12.1$ vs.\
$13.3$~km/s, $52/99$ per-galaxy wins).

\textbf{Held-out cross-validation (F147).} \emph{The 14-galaxy
subsample on which $(a_0^{\rm reg}, q)$ were originally tuned is
contained within the 175-galaxy comparison set;} the headline median
is therefore partially in-sample. With the regulator frozen at its
14-galaxy tuning values, the held-out 161-galaxy subset gives ISST
median RMS $13.86$~km/s vs.\ MOND $13.00$~km/s, ISST winning
$76/161$ ($47.2\%$). On the held-out $Q{=}1$ subsample ($N=87$) the
result is ISST $12.12$ vs.\ MOND $13.34$~km/s, ISST winning $44/87$
($50.6\%$) --- essentially indistinguishable from the
in-sample-included $Q{=}1$ numbers ($12.1$ vs.\ $13.3$ on $52/99$).
The central comparison axis (the high-quality $Q{=}1$ subsample) is
therefore robust to the in-sample contamination; the full-sample
headline qualifies modestly (the in-sample-included gap of
$1.10$~km/s tightens to $0.86$~km/s on held-out, with MOND retaining
a small lead). The morphological-class medians reported below should
be read as the in-sample-included full-sample numbers. The baryonic Tully--Fisher
slope is $1.15$, consistent with the literature value
$1.0\pm 0.1$~\cite{McGaugh2012,Tully2008}. Status: \DEMO. ISST does not
\emph{derive} MOND: the internal $a_0$ is $120\times$ smaller than
Milgrom's, and the formula does not assume the deep-MOND scaling.
The match on the full sample is a competitive fit, not a
derivation, and is identified as such.

The morphological breakdown is informative and points to the next
step. ISST matches MOND on dwarfs and irregulars (median 10.1 vs.\
9.6~km/s, $N{=}81$), wins on late-type spirals (Sc--Sd, 10.1 vs.\
11.5~km/s, $N{=}48$), and loses on early-type spirals (Sa--Sbc,
24.9 vs.\ 21.4~km/s, $N{=}43$). The Sa--Sbc loss-channel is
star-dominated, which is exactly the regime where the local-bounded
$F_{\rm flat}$ reading predicts $f_\star\approx 0$ and the
$(1{+}f)$ enhancement does \emph{not} act on the dominant
component. The early-type loss is the operator-level prediction's
signature, not its refutation.

\subsection{Radial Acceleration Relation}

The RAR~\cite{McGaughLelliSchombert2016,LelliEtAl2017} emerges from
the same regulator structure. The cross-over
acceleration is the derived $a_{\rm crit}$; the deep-MOND-like
asymptote follows from the $(a_0/a_{\rm bar})^q$ branch with
empirical $a_0$ and $q$. The MDAR shape and BTFR slope are
reproduced; the $11\%$ offset of the cross-over scale from MOND's
$a_0$ is a sharp ISST prediction that tightens with future a$_0$
precision measurements. Status: \DEMO\ (derived $a_{\rm crit}$;
empirical $a_0,q$).

\subsection{Dispersion-supported gas-poor dwarfs and tidal dwarf galaxies}
\label{sec:df2_tdg}

The local-bounded $F_{\rm flat}$ reading (Sec.~\ref{sec:f}) makes a
sharp prediction for galactic systems where matter is virialised in
its own self-potential without a deeper external halo: $f \to 0$ for
the dominant component, and ISST recovers Newtonian dynamics through
Branch~B selection (F08).

\textbf{NGC 1052-DF2~\cite{vanDokkum2018,Trujillo2019}.}
With $M_\star \approx 2\times 10^8\,M_\odot$, $r_{\rm eff} \approx 2.2$~kpc,
and $M_{\rm gas}/M_\star \ll 0.01$, the system has
$v_{\rm esc} \approx 12.5$~km/s, $v_{\rm max}^{\rm EM} = 2.56\,\sigma_\star
\approx 21.8$~km/s. The OP2 condition $v_{\rm max}^{\rm EM} \geq v_{\rm esc}$
is satisfied: stellar matter fills its accessible phase space,
$f_\star \approx 0$. Static spherical equilibrium forces Branch~B
(F08), $G_{\rm eff} = (1+f_\star)/\Psi_0 = G_N$. ISST predicts
$\sigma_{\rm DF2}^{\rm ISST} = \sqrt{GM_\star/(5\,r_{\rm eff})}
\approx 8.9$~km/s, consistent with the observation
$8.5 \pm 2.3$~km/s within $1\sigma$. \emph{The local-bounded
$F_{\rm flat}$ reading was introduced into the paper foundations
prior to a quantitative DF2 audit, but the reading itself is
motivated by the A05 Aczél--Shore--Johnson axioms and the
entropy-matching cutoff at $v_{\rm max}^{\rm EM}=2.56\sigma_{\rm eff}$
identified independently by F119, not constructed post-hoc to
handle DF2}. DF2 is therefore a successful application of an
independently-motivated reading, sharing the same operator-level
mechanism (zero $f$ for matter that fills its accessible phase
space) that explains the Sa--Sbc loss-channel in the SPARC sample
(Sec.~\ref{sec:galaxy}). The $\sigma$ prediction has no quoted
model uncertainty, so ``consistent with $1\sigma$'' is the honest
framing rather than ``predicted to $0.2\sigma$''.

\textbf{Tidal dwarf galaxies (Lelli et al.~2015~\cite{Lelli2015},
NGC~5291 system).} For TDG gas at $\sigma_{\rm gas} \approx 8$~km/s
in $v_{\rm esc} \approx 50$~km/s self-potential, the local-bounded
reading gives $f_{\rm gas}^{\rm TDG} = \alpha[3\ln(v_{\rm esc}/\sigma_{\rm gas})
- 2.825] \approx 2.3$ (with $\alpha = 0.869$ from F119), predicting
a moderate $(1+f) \approx 3.3$ enhancement. Lelli et al.'s reported
$M_{\rm dyn}/M_{\rm bar} \approx 1$ is in mild tension with this
prediction. The tension is sensitive to the assumed gas $\sigma$
(turbulent + thermal): at $\sigma_{\rm gas} \approx 15$~km/s the
prediction recovers Newton ($f \to 0$ via OP2). Improved
$\sigma_{\rm gas}$ measurements would discriminate.

\textbf{Discriminator vs.\ MOND.} MOND with the external field effect
predicts Newtonian dynamics for dwarfs in a host's deep external field
(DF2 in NGC~1052) but full MOND for isolated dwarfs (TDGs in
low-density environments). ISST predicts Newton for DF2 (same answer
as MOND-with-EFE) but moderate enhancement for TDGs in low-density
environments (different answer). Isolated low-density
dispersion-supported dwarfs are therefore the cleanest
ISST-vs.-MOND-with-EFE discriminator, identified as a future-work
observational target.

\textbf{Kill condition.} Three or more dynamically-confirmed
dark-matter-free galaxies in static equilibrium, with $\sigma$
measured precisely enough to set $\sigma_{\rm bar}/v_{\rm esc} > 0.4$
(satisfying the OP2 condition), showing $M_{\rm dyn}/M_{\rm bar} > 1.5$,
would falsify the local-bounded reading at A05.OP2. DF2 alone is
consistent with the prediction; future systems
(e.g.\ DF4, AGC~114905~\cite{Mancera2022}) sample the same regime.

%==========================================================
\section{Cluster-scale phenomenology}
\label{sec:cluster}

\subsection{The cluster mass identity (A04)}

The ISST identity at the action level is
\begin{equation}
\label{eq:A04}
\Omega_m \;=\; (1+\fprim)\,\Omega_b.
\end{equation}
This is a structural prediction: the same $(1{+}f)$ factor that
enhances galactic rotation accounts for the cluster mass
discrepancy. With the SM-derived $\fprim = 5.66\pm 0.06$
(Sec.~\ref{sec:fprim}) and Planck $\Omega_b = 0.0493$, the
prediction is $\Omega_m = 0.329$, against the observation $\Omega_m
= 0.3153\pm 0.0073$~\cite{Planck2018}: a $4.2\%$ overshoot,
$1.6$--$1.8\sigma$ depending on whether the model uncertainty is
propagated. The structural identity is \PROVEN; the numerical
realisation is \DEMO\ at $1.6$--$1.8\sigma$ tension, with the
residual stated openly as a limit of the sharp-transition
approximation. The conventional
practice of inverting \eqref{eq:A04} to extract $\fprim = 5.40$
would absorb the residual---we do not.

\subsection{Modified lensing and the lighthouse mechanism}

In ISST, gravitational lensing measures the $\Psi$ field where the
photons travel, not the local mass at the source. The convergence
kernel under F72 is
\begin{equation}
\kappa_{\rm ISST} = \frac{1+f}{\Psi_0}\,\kappa_{\rm GR},
\end{equation}
and the photon-trajectory parameter is $\eta = \Phi/\Phi_N = 2$
exactly. The framework verdict is
\textsf{modified\_lensing\_scalar\_tensor} \textsc{[allowed]}, with
the GR Poisson source explicitly \textsc{[denied]} and replaced by
$(1{+}f)\rho$. Status: \PROVEN\ for the kernel form; the
\textsc{[allowed]} verdict descends from the slip predicate stamped
by F84.

The information content of matter splits into a primordial floor
$\fprim$ frozen at QCD-era temperatures ($\sim 10^{12}\,{\rm K}$) and
a post-formation working memory $f_s$ accumulated during structure
formation. A merger shock at $\sim 10^8\,{\rm K}$ thermalises
$f_s$---the gas's working memory, accumulated by violent relaxation
during structure formation~\cite{LyndenBell1967}---but cannot touch
the QCD-frozen $\fprim$. The gas's effective gravitational coupling drops from
$(1+\fprim+f_s)\rho$ to $(1+\fprim)\rho$. Stars, which pass through
the collision essentially without interaction, retain their full
working memory. Lensing follows the stronger source.

\subsection{Bullet Cluster}
\label{sec:bullet}

The status of the Bullet Cluster in ISST is
\textbf{mechanism-demonstrated, local-columns-close-bulk-of-gap}
(F131). It is not \emph{solved} or \emph{proven}. The lighthouse mechanism predicts a per-mass contrast
$(1+f_s^\star)/(1+f_s^{\rm gas,shock}) \approx 4.55/1.21 = 3.76$
between the surviving star-dominated lensing and the shock-erased
gas. \emph{The values $f_s^\star \approx 3.55$ and
$f_s^{\rm gas,shock} \approx 0.21$ used in this contrast are fitted
to the local-column convergence ratio inferred from
Paraficz~\emph{et al.}~(2016)\cite{Paraficz2016}; their derivation
from first principles --- specifically, the violent-relaxation
working memory accumulated by stars during structure formation
versus the residual primordial-only floor of shock-thermalised gas
--- is identified as \OPEN.} Counted honestly, they are empirical
parameters of the lighthouse application, not derived consequences
of the action. The Paraficz local-column densities themselves are
inferred under a mass model that includes a dark-matter component;
treating them as the total baryon column requires re-interpreting
their model decomposition, an additional assumption that the
quantitative contrast is sensitive to. With these caveats: at local
columns, the per-mass contrast required for
$\kappa_{\rm gal}/\kappa_{\rm gas}\geq 2$ is $1.6$; the operator
delivers $3.76$, closing with margin throughout the Paraficz
$1\sigma$ envelope under the empirical $f_s$ values.

The full per-pixel $\kappa$ ratio at FWHM-matched smoothing requires
the local $\Sigma_{\rm baryon}$ contrast at the two peaks. Three
recent works push directly on this observable.
Cha~\emph{et~al.}~(2025)~\cite{Cha2025} present the highest-resolution
Bullet mass reconstruction to date from JWST, combining $146$
strong-lensing constraints from $37$ systems with
$398$~sources/arcmin$^2$ of weak lensing; the resulting $\kappa$ map
resolves \emph{at least three subclumps aligned with the brightest
cluster galaxies}, and---critically for the lighthouse mechanism---the
intracluster-light distribution is found to track the JWST mass map
to a modified Hausdorff distance of
$19.80\pm 12.46$~kpc. The ICL traces stars by construction, so this
$\sim 20$~kpc scale agreement is direct evidence that the
$\kappa$ peak tracks stellar mass at the resolution of the
JWST data, supporting the lighthouse mechanism's central
prediction.
Rihtar\v{s}i\v{c}~\emph{et~al.}~(2026)~\cite{Rihtarsic2026} provide
an updated GR lensing $\kappa$ map of the Bullet using the latest
galaxy-member catalogue (219 spectroscopic members with new JWST
photometry).
Hernandez~(2026)~\cite{Hernandez2026} computes the QUMOND
prediction for the same baryonic distribution---the 219 stellar
galaxies modelled as PIEMD profiles plus a $\beta$-model X-ray
gas---and compares pixel-by-pixel against the Rihtar\v{s}i\v{c} GR
$\kappa$ map. The QUMOND map matches GR to a residual
$|\Delta\kappa|\lesssim 0.15$ throughout (in the standard units of
$1.83\times 10^{9}\,M_\odot/\mathrm{kpc}^2$), comparable to or
smaller than the residuals between independent GR lensing analyses.
The galaxies, $7\%$ of the total baryonic mass, contribute $48\%$
of the QUMOND phantom mass, because point-like sources concentrate
the QUMOND signal where extended gas does not.

ISST's lighthouse operator predicts the same morphology
($\kappa$ peaks at star-dominated stellar mass, not at the gas
centroid) through a distinct mechanism: per-pixel $f_s$-thermalisation
contrast at the action level (intact stars retain
$f_s^\star\approx 4.55-1$, shocked gas has
$f_s^{\rm gas,shock}\approx 0.21$). The qualitative agreement
between Cha's $\sim 20$~kpc ICL/$\kappa$ alignment and the lighthouse
prediction is direct (\DEMO); extracting
$\kappa_{\rm gal}/\kappa_{\rm gas}$ at the two main peaks from the
Rihtar\v{s}i\v{c} $\kappa$ map and comparing to the operator
$(1+f_s^\star)/(1+f_s^{\rm gas,shock}) \approx 3.76$ is identified
as the next quantitative step (\OPEN). The standard $\Lambda$CDM
reading of the same data invokes a collisionless dark-matter
component with a self-interaction cross-section bound below
$\sigma/m\lesssim 1\,{\rm cm^2/g}$~\cite{Markevitch2004,Clowe2006};
ISST does not introduce such a particle.

\textbf{The Bullet does not discriminate ISST from MOND.}
Hernandez~\cite{Hernandez2026} demonstrates that QUMOND reproduces
the Bullet $\kappa$ map without dark matter through compactness
(point-like galaxies $\to$ peaked phantom signal; diffuse gas $\to$
diffuse phantom signal). ISST reproduces it through processing
history ($f_s$ stripped from shocked gas; $f_s$ retained in stars).
Both predictions agree at the Bullet because the gas is shocked
\emph{and} diffuse, while stars are intact \emph{and} compact---the
two mechanisms are observationally degenerate at this system.
The discriminator between ISST and MOND is therefore not the Bullet
itself but systems where the two mechanisms predict differently:
$a_{\rm crit}$ vs.\ $a_0$ (Sec.~\ref{sec:acrit}, $11\%$ offset
predicted), native cosmology (Sec.~\ref{sec:cosmology}, MOND has
none), $\theta_*$ (Sec.~\ref{sec:theta}, $-0.83\%$ ISST, MOND
silent), and \emph{ram-pressure-stripped jellyfish galaxies} where
gas and stars share the same Newtonian acceleration but have
different processing histories: MOND predicts equal enhancement of
both components in the low-acceleration regime; ISST predicts
suppressed enhancement of the stripped gas tail
($f_s\to 0$ from ram-pressure shock heating), unsuppressed
enhancement of the intact stellar disk. This system-class is
identified in Sec.~\ref{sec:future} as the live falsifier between
the two theories.

Earlier ISST attempts that imported a halo ontology (F65--F71)
violated F01's pointwise source structure and have been retracted
(Sec.~\ref{sec:dropped}).

\subsection{Abell 520 (pile-up geometry)}
\label{sec:a520}

The lighthouse mechanism makes a second, independent prediction in
multi-body pile-up mergers (F132). Dynamical
friction~\cite{ChandrasekharDF,BinneyTremaine2008} concentrates
massive bulges (high $f_s^\star$, surviving the merger crossing
$\sim 0.5\,{\rm Gyr}$~\cite{Loubser2020}) at the merger centre, where
shocked gas also accumulates centrally. Convergence tracks the surviving bulge
population, which spatially overlaps the gas, with
\begin{itemize}
\item[(i)] stellar-population prediction: passive, red,
old populations dominating the merger core, with rapid quenching and
no starburst;
\item[(ii)] mass-to-light prediction: $M/L_R \in [38,80]$ in solar
units at the core.
\end{itemize}
Prediction (i) is independently confirmed by Deshev~\emph{et~al.}~(2017)
\cite{Deshev2017} on $400+$ galaxies. Prediction (ii) is consistent
with Clowe~\emph{et~al.}~(2012)~\cite{Clowe2012}; the alternative
Jee~\emph{et~al.}~(2014) reading of the same HST/ACS
data~\cite{Jee2014} would put F01 in tension. The Clowe--Jee
disagreement is unresolved in the literature; ISST predicts the
Clowe reading. Status: \INDIC\ (consistent-or-indicated).

\subsection{Differential void-lensing prediction}

The same kernel makes a sharp differential prediction at the
void scale: the wall-versus-void $(1+f)/\Psi_0$ ratio implies a
$+2.5\%$ enhancement of the dressed coupling in voids relative to
local laboratory measurements,
\begin{equation}
G_{\rm void}/G_{\rm local} = 1.025\quad (\text{F72, +2.5\%}).
\end{equation}
This is a discriminator versus chameleon and metric $f(R)$ scenarios,
both of which screen \emph{toward} GR in low-density environments.
Status: \INDIC. Euclid void-lensing tomography is the live test
(Sec.~\ref{sec:falsification}).

\subsection{Missing mass: no particle}

If \eqref{eq:A04} accounts for the full cluster mass discrepancy
through baryons amplified by $(1+\fprim)$, no dark-matter particle
is needed. ISST predicts that direct WIMP, axion, and sterile-neutrino
detection at CDM density continue to return null. This is a
hard-falsifier prediction (Sec.~\ref{sec:falsification}), not a
post-hoc consistency check.

%==========================================================
\section{Cosmology}
\label{sec:cosmology}

\subsection{Inhomogeneous background: Wiltshire as minimal realisation}
\label{sec:wiltshire_forced}

Branch~A combined with dust matter is observationally incompatible
with a spatially flat homogeneous FRW background, and the minimum
geometric structure that recovers the observed matter-era expansion
is a domain split into Milne-like voids and walls with a spatial
$\Psi$ profile. The full algebra is in
Appendix~\ref{app:wiltshire_forced}; we summarise here.

\textbf{Claim 1 (geometric).} On a spatially flat homogeneous FRW
metric $ds^2 = -dt^2 + a(t)^2\delta_{ij}dx^i dx^j$, the Ricci scalar
is $R = 6(\ddot a/a + H^2) = 6(\dot H + 2H^2)$. The pointwise
Branch-A constraint $R = 0$ therefore reduces to
\begin{equation}
\label{eq:branchA_FRW_ode}
\dot H + 2H^2 = 0 \quad\Longrightarrow\quad H(t) = \frac{1}{2t},
\quad a(t)\propto t^{1/2}.
\end{equation}
This is a purely \emph{geometric} consequence of $R=0$: no matter
input was used. The expansion law is forced.

\textbf{Claim 2 (observational).} The $(tt)$ component of
\eqref{eq:eom} on flat FRW with $\Psi = \Psi(t)$ and perfect-fluid
dust $T^t{}_t = -\rho$ reads
\begin{equation}
\label{eq:tt_FRW}
3\Psi H^2 + 3H\dot\Psi = 8\pi(1+f)\rho.
\end{equation}
Substituting $H = 1/(2t)$ from Claim~1 and dust scaling
$\rho = \rho_0(t_0/t)^{3/2}$ (from $\rho a^3 = \mathrm{const}$ and
$a\propto t^{1/2}$), this becomes the first-order linear ODE
\begin{equation}
\label{eq:psi_ode}
\dot\Psi(t) + \frac{\Psi(t)}{2t} \;=\; \frac{16\pi}{3}(1{+}f)\rho_0\,t_0^{3/2}\,t^{-1/2},
\end{equation}
whose general solution (with integrating factor $t^{1/2}$) is
\begin{equation}
\label{eq:psi_general}
\Psi(t) \;=\; A\,t^{1/2} + K\,t^{-1/2},
\qquad
A = \frac{16\pi}{3}(1{+}f)\rho_0\,t_0^{3/2},
\end{equation}
with $K$ a free integration constant. The system is
\emph{algebraically consistent}: $\Psi(t_0) = At_0^{1/2} + Kt_0^{-1/2}$
is finite, $G_{\rm eff} = (1{+}f)/\Psi(t_0)$ is finite, and the
trace transport $\Box\Psi = -(8\pi/3)(1{+}f)\rho$ is satisfied
identically by both modes (the $K t^{-1/2}$ mode is the
homogeneous solution; the $A t^{1/2}$ mode is the particular
solution sourced by dust).

\emph{The incompatibility is observational, not algebraic.}
Claim~1 forces $a(t) \propto t^{1/2}$, which is the radiation-era
expansion law. The matter-era expansion observed at intermediate
redshift through SN-Ia luminosity distances, BAO scales, and CMB
acoustic-peak positions follows $a(t) \propto t^{2/3}$ at the
relevant epochs in any matter-dominated dust cosmology. Branch~A on
flat FRW therefore predicts the wrong expansion law during the
matter era, regardless of the integration constant $K$. The
mismatch is at the level of $\mathcal O(1)$ in the deceleration
parameter ($q = 1$ for $a\propto t^{1/2}$ vs.\ $q = 1/2$ for
$a\propto t^{2/3}$), not a subleading correction.

The same algebra shows traceless matter ($T = -\rho + 3P = 0$,
i.e.\ $P = \rho/3$, radiation) is consistent with Branch-A on flat
FRW because the transport equation $\Box\Psi = (8\pi/3)(1+f)T$ then
admits $\Psi = \Psi_0 = \mathrm{const}$, and \eqref{eq:tt_FRW}
reduces to a single algebraic identity --- consistent with the
radiation-era a $\propto t^{1/2}$ scaling. Branch-A flat FRW is the
correct radiation-era cosmology; it is inconsistent with matter-era
observations.

\textbf{Consequence.} Reconciling Branch~A (geometric, $R=0$ forces
$a\propto t^{1/2}$) with matter-era observations ($a\propto t^{2/3}$
at intermediate $z$) requires a background that is not flat
FRW. The minimum geometric structure satisfying the constraint is a
\emph{two-domain split}:
\begin{itemize}
\item \emph{Voids}, with $\rho \to 0$ and $T \to 0$, admit Milne
expansion ($a_V\propto t$, $H_V = 1/t$, $R = 0$ trivially in the
traceless limit);
\item \emph{Walls}, with dust matter, admit $R = 0$ pointwise via a
non-FRW metric: the planar ansatz
$ds^2 = -dt^2 + a(t)^2(dx^2+dy^2) + b(t,z)^2\,dz^2$ with
$\Psi = \Psi_0(t) + \psi(z)$ satisfies $R = 0$ pointwise iff
$a(t)\propto t^{2/3}$ and $\partial_z^2\psi = -(8\pi/3)(1+f)\rho$, with
the ($z$-derivative) terms in $b(t,z)$ cancelling identically between
$R_{tt}$ and $R_{zz}$ at the bound-wall fixed point
(Appendix~\ref{app:wiltshire_forced} §\ref{sec:wall_pointwise}).
\end{itemize}
Spatial averaging the wall solution recovers
\eqref{eq:wall_friedmann_main} with the $16\pi/9$ coefficient
verified independently (see Sec.~\ref{sec:wall_friedmann}).

\textbf{Wiltshire selection.} The two-domain split is the
\emph{minimal} geometric realisation that reconciles
Claim~1 (geometric: $R=0 \Rightarrow a\propto t^{1/2}$ on flat FRW)
with the observed matter-era expansion (Claim~2: $a\propto t^{2/3}$
at intermediate $z$). It is not the unique such realisation;
strict exclusion of LTB/Szekeres alternatives is identified below. The selection of \emph{Wiltshire's} specific
two-domain timescape geometry over alternative inhomogeneous
constructions (Lema\^itre--Tolman--Bondi, Szekeres) rests on three
criteria, given in decreasing strength:
(i)~minimal domain count---two domains (voids + walls) is the smallest
count for which the void/wall split satisfies $R=0$ separately in
each domain without continuous tuning of a bang-time function;
(ii)~explicit pointwise $R=0$ realisation---the planar wall metric
above is constructed; for LTB or Szekeres, $R = 0$ pointwise on dust
requires solving a coupled system $E(r)$, $t_B(r)$ with no closed-form
solution we have verified;
(iii)~empirical match to the apparent-acceleration phenomenology
parameterised by $f_{v0}\approx 0.762$ from
SN-Ia~\cite{Dam2017,Wiltshire2009}.
Strict \emph{exclusion} of LTB/Szekeres alternatives is identified
as a residual classification problem (\OPEN); the constructive
realisation of \emph{some} inhomogeneous background reconciling
Branch~A + dust + matter-era observations is \PROVEN\ here.
Frameworks \textsf{buchert\_averaging} and
\textsf{wiltshire\_two\_domain} are \textsc{[allowed]} on the
passport; LTB and Szekeres remain unstamped.

\subsection{Wall Friedmann equation}
\label{sec:wall_friedmann}

The wall Friedmann equation derived from \eqref{eq:eom} is
\begin{equation}
\label{eq:wall_friedmann_main}
H_w^2 \;=\; \frac{16\pi}{9}\,\frac{(1+f)\rho_b}{\Psi},
\end{equation}
which under $(1+f)\rho_b\equiv\rho_m$ reproduces the
Einstein--de~Sitter tracker for matter-era walls up to an overall
rescaling of cosmic time. The Milne void is automatically
Branch-A-compatible. With spatial $f(x)$, $f$ is replaced by its
mass-weighted wall average; the uniform-$f$ case is the
zero-correlation special case. The $16\pi/9$ coefficient is
\emph{derived pointwise} from the planar wall ansatz via the
Branch-A field equations
(Appendix~\ref{app:wall_friedmann_proof}, F141), with explicit
verification that energy-momentum conservation holds at every step
of the anisotropic-to-isotropic transition. The $2/3$ factor relative
to standard Buchert dust is identified as the fraction of the
gravitational source absorbed into the spatial $\Psi$ profile by
the trace transport (F01 transport equation): the spatial Poisson
$\partial_z^2\psi = -(8\pi/3)(1{+}f)\rho$ absorbs $1/3$ of the source,
leaving $(16\pi/3)(1{+}f)\rho$ on the $(tt)$ RHS, divided by the
$3\Psi$ prefactor of $G^t{}_t$ to give $16\pi/9$.
Coordinate-independence is verified in synchronous, conformal, and
harmonic gauges by general covariance of the field equation.
Status: \PROVEN.

\emph{Note on coefficients across geometries.} The wall Friedmann
$16\pi/9$ is \emph{not} obtained by scaling the standard
$8\pi G/3$ by the rotation-curve enhancement $4/3$ (which would give
$32\pi/9$). It is derived structurally from the (tt) projection of
\eqref{eq:eom} on the wall slab, where the static spatial $\Psi$
profile absorbs $(8\pi/3)(1{+}f)\rho$ via the Branch-A transport, so
that the cosmological source reduces from $8\pi(1{+}f)\rho$ to
$(16\pi/3)(1{+}f)\rho$ on the RHS, then divides by $3\Psi$ for $H^2$.
The three structural coefficients in this paper---$4/3$ (rotation,
$\nabla^2\Phi$, Eq.~\ref{eq:poisson}), $2/3$ (lensing,
$\nabla^2\Phi_N$, Appendix~\ref{app:Cgrad}), and $16\pi/9$
(cosmological wall, Eq.~\ref{eq:wall_friedmann_main})---arise from
three distinct field-equation projections in three distinct
geometries (axisymmetric disk, weak-field point mass, planar wall
slab) and are independently consistent with the same F01 EOM.

\subsection{$H_0$ and the apparent acceleration}
\label{sec:hubble}

The two-domain background admits three distinct Hubble rates whose
relations are explicit, not phenomenological. The full pipeline is
in Appendix~\ref{app:h0_frame}; we summarise the algebra here.

\textbf{The three frames.}
(i)~The \emph{bare} frame is the Buchert volume average, with scale
factor $\bar a(t)$ and Hubble rate $\bar H = \dot{\bar a}/\bar a$
defined on the bare time coordinate $t$.
(ii)~The \emph{wall} frame is the proper rest frame of a galaxy
observer in a wall, with proper time $\tau$ related by
$dt = \gamma_w(\tau)\,d\tau$.
(iii)~The \emph{dressed} frame is what a wall observer infers when
fitting their distance-redshift data with a homogeneous FLRW
template; it is what Planck and SH0ES report.

\textbf{Wiltshire dressing relations.} On the matter-dominated
two-domain tracker~\cite{Wiltshire2007,Wiltshire2009}:
\begin{equation}
\label{eq:dressing_R}
H_0^{\rm dressed} \;=\; \mathcal{R}(f_{v0})\,H_0^{\rm bare},
\qquad
\mathcal{R}(f_v) \;=\; \frac{4f_v^2 + f_v + 4}{2(2+f_v)},
\end{equation}
which evaluates at $f_{v0} = 0.762$ to $\mathcal{R} = 1.2825$. The
present-epoch lapse $\gamma_0 = 1.348$ controls density dressing
$\Omega_m^{\rm dressed} = \gamma_0^3\,\Omega_m^{\rm bare}$, distinct
from the Hubble dressing $\mathcal{R}$.

\textbf{ISST-native bare $H_0$.} A model-independent input from
Planck (the comoving angular-diameter distance $D_A^{\rm com}(z_*) =
14130.8$~Mpc derived from $r_s$ and $\theta_*$, no $\Lambda$CDM
assumption) feeds the ISST Branch-A two-domain tracker. With ISST's
parameter commitments ($f_{\rm prim} = 5.664$, $f_{v0} = 0.762$,
$\Omega_b h^2 = 0.02237$) integrated through the Buchert background
(F26, F26a), the tracker outputs
\begin{equation}
\label{eq:bare_h0}
H_0^{\rm ISST,bare} \;=\; 57.56\,{\rm km\,s^{-1}\,Mpc^{-1}}.
\end{equation}
Applying \eqref{eq:dressing_R} gives the corresponding dressed
\begin{equation}
\label{eq:dressed_uniform}
H_0^{\rm dressed,uniform-f} \;=\; 1.2825 \times 57.56 \;=\;
73.82\,{\rm km\,s^{-1}\,Mpc^{-1}},
\end{equation}
the uniform-$f$ tracker prediction. The committed paper value
\begin{equation}
\label{eq:H0_committed}
H_0^{\rm ISST,dressed} \;=\; 61.79\,{\rm km\,s^{-1}\,Mpc^{-1}}
\end{equation}
arises from the $K_{\rm env}$-corrected closure
(Sec.~\ref{sec:tracker}) which incorporates an environmental
rescaling that the uniform-$f$ run lacks; the $73.82$ uniform-$f$
result is reported here transparently as the diagnostic input to
that closure, not as the prediction.

\textbf{The Planck and SH0ES recoveries.} A wall observer fitting
ISST's two-domain $d_L(z)$ with a $\Lambda$CDM template recovers
different $H_0$ values depending on \emph{which} template parameters
are anchored; the F85 forward pipeline computes both:
\begin{equation}
\label{eq:F85_pipeline}
H_0^{\rm Planck-template} \;=\; 62.8,
\qquad
H_0^{\rm SH0ES-ladder} \;=\; 70.6\,{\rm km\,s^{-1}\,Mpc^{-1}}
\end{equation}
where the Planck-template recovery uses the CMB-anchored
$(\Omega_m, \Omega_b h^2)$ and the SH0ES-ladder recovery uses
ladder-anchored $M_B$ on $z\in[0.023,0.15]$. Observed
values are $67.4$ and $73.0$ respectively. The
ISST-internal spread $70.6 - 62.8 = 7.8$~km/s/Mpc reproduces
\emph{the structure} of the Hubble tension; the residual
$2.5\%$--$5\%$ offset of each pipeline output from the corresponding
observed value reflects the gap between the effective-FRW Wiltshire
proxy and a full two-domain photon-path integration with $K_{\rm env}$
inheritance and is reported honestly as such (see
Sec.~\ref{sec:tracker} and Appendix~\ref{app:h0_frame}). Status:
\DEMO\ (structure of the tension reproduced; literal recovery of
$73.0$ and $67.4$ \OPEN).

\subsection{Coincidence and void fraction tracker}
\label{sec:tracker}

A three-parameter $\chi_w(z)$ sigmoid closes the CMB and SN
anchors~\cite{Dam2017,NazerWiltshire2015}.
The amplitude $A_{\rm thermo}=\fprim/N_{\rm stages}=1.416$ from the
arithmetic mean of four SM freezeout-stage increments matches the
phenomenological best-fit to $1.1\%$. The pivot
$z_{\rm tr}=1.555$ from the Sheth--van~de~Weygaert two-barrier
(void-in-cloud) abundance~\cite{Sheth2004} with
$\sigma_{8,{\rm eff}}=4.92$ matches
to $3.7\%$. The environmental rescaling $K_{\rm env}=1.484$ is
recovered as the squared ratio of two independent Hubble outputs.
The full closure gives $(A,z_{\rm tr},n,H_0)=(-0.688,1.555,4.83,61.79)$
against phenomenological $(-0.65,1.50,2.00,61.80)$. The
shape index $n$ is a sharper-than-phenomenological prediction and
is the one residual. Status: \DEMO\ (amplitude, pivot, $H_0$);
\OPEN\ ($n$).

\subsection{Growth of perturbations: $f\sigma_8$}

The wall-growth equation derived from \eqref{eq:eom} on the
Wiltshire two-domain background, with $(1+f)$ source enhancement,
$\Psi$ friction, and AP retemplating, gives
\begin{equation}
f\sigma_8(z=0.57)^{\rm ISST} = 0.430\pm 0.028,
\end{equation}
versus $\Lambda$CDM 0.472 (F89, framework
\textsf{wiltshire\_two\_domain} \textsc{[allowed]} and
ISST-native wall growth). On 8 RSD surveys~\cite{BeutlerEtAl2012,HowlettEtAl2015,AlamEtAl2017,BlakeEtAl2012,PezzottaEtAl2017,Adil2024}
the joint $\chi^2 = 8.6$--$9.5$ ($p\approx 0.30$--$0.45$). The shape
distinguishability from $\Lambda$CDM is forecast at $13\sigma$ at
$z=1.1$ with full DESI~DR2~\cite{DESIDR2,DESIDR1} + Euclid Y1
\cite{EuclidForecast,EuclidYellow} binning. DESI~DR2 BAO+RSD
results are now published~\cite{DESIDR2}; a quantitative
ISST-vs-$\Lambda$CDM comparison on the released DR2 $f\sigma_8(z)$
points (specifically the high-$z$ bins at $z\gtrsim 0.7$ where the
ISST shape diverges most from $\Lambda$CDM) is identified as the
immediate next step, with the kill condition specified in
Sec.~\ref{sec:falsification}. Status: \PROVEN\ on the eight current
RSD surveys at the joint-$\chi^2$ level; \OPEN\ for the
DR2-released $f\sigma_8(z)$ comparison. The earlier templating
``crisis'' on the Wiltshire background was an artefact resolved by
the AP retemplating used here.

\subsection{Integrated Sachs--Wolfe: structural null}

The integrated Sachs--Wolfe contribution from individual line-of-sight
voids on Branch~A gives a per-photon RMS of order
$\sqrt{N}\langle A\rangle$ with $N\approx 140$ voids, but the
ensemble mean is identically zero (F106/F107/F108): the line integral
of $\dot\Psi$ along a void is structurally cancelling under the F01
transport equation across the photon flight time
$\tau_{\rm flight}/\tau_{\rm cycle}\leq 1.7\%$. ISST therefore does
not predict a coherent ISW signal at the linear level beyond the
LCDM baseline. Status: \PROVEN\ (ISST-NATIVE; the latent v2
\textsf{isst\_void\_isw} pattern, identified for explicit framework
registration).

\subsection{The CMB acoustic angle: first-pass Boltzmann result}
\label{sec:theta}

\textbf{The theory's reach across scales narrows here.} The
preceding subsections established cosmological successes
(wall-Friedmann derivation, $H_0$ frame artefact, $f\sigma_8$
agreement, structural ISW null) on the Wiltshire-effective
background. The CMB perturbation level is where the
single-action programme runs into its first hard limitation: the
strict uniform-$f$ source is falsified at the third peak by
$>10\sigma$, and the env-dep-$f$ closure provides only
percent-level corrections (Sec.~\ref{sec:env_dep_f_closure}).
The first-pass acoustic-angle result reported here uses an
empirical proxy at the perturbation level; the proxy's
justification, its systematic uncertainty, and the open question
it leaves are stated explicitly. This is the most honestly open
result in the paper.

A first-pass run of the CAMB Boltzmann
hierarchy~\cite{LewisChallinor2011} on the Wiltshire-effective
background~\cite{Wiltshire2007,Wiltshire2009,Dam2017} gives, at the
ISST-committed parameters, a sub-percent residual on the CMB
acoustic angle. The Wiltshire dressing of the late-time expansion is
captured by an effective dark-energy parameterisation
$(w_0, w_a) = (-0.674, -1.241)$ extracted from low-$z$ supernova
matching; the matter density is set by the A04 identity
$(1+\fprim)\Omega_b h^2 = \Omega_m h^2$ at the SM-prediction
$\fprim = 5.664$ (giving $\Omega_m h^2 = 0.149$, the same $4.2\%$
overshoot vs.\ Planck reported in Sec.~\ref{sec:fprim}); the Hubble
constant is the dressed value $H_0 = 61.79$~km/s/Mpc.

\emph{Boltzmann configuration (reproducible inputs).} The CAMB run
uses the Path~A treatment: $(1+\fprim)\Omega_b$ is modelled as the
total matter component at the perturbation level, with the
information-enhanced extra contribution carried by an effectively
collisionless CDM component (the empirical proxy whose derivation
from F01 alone is the open problem of
Sec.~\ref{sec:env_dep_f_closure}). Specifically: $\Omega_b h^2
= 0.02237$ (Planck); $\Omega_c h^2 = \Omega_m h^2 - \Omega_b h^2
= 0.149 - 0.02237 = 0.127$ enters as the CDM proxy carrying
$\fprim\Omega_b$; $H_0 = 61.79$~km/s/Mpc; $w_0 = -0.674$,
$w_a = -1.241$; remaining cosmological parameters
($n_s, A_s, \tau$) at Planck values. The modified CLASS source
patch implementing the strict uniform-$f$ test (below) and the
exact CAMB \texttt{params.ini} configuration files will be
archived as supplementary material on acceptance. Four
configurations were run:
\begin{equation*}
\begin{array}{lcl}
\textrm{Config} & 100\,\theta_* & \textrm{description} \\
\hline
\textrm{C0} & 1.04118 & \textrm{$\Lambda$CDM Planck baseline} \\
\textrm{C1} & 1.03107 & \textrm{drop-$\Lambda$, ISST $H_0$} \\
\textrm{C2} & 1.07125 & \textrm{Wilt DE, uniform-$f$ $H_0=75.28$} \\
\textrm{C3} & 1.03258 & \textrm{ISST committed (Wilt DE + dressed $H_0$)}
\end{array}
\end{equation*}
The headline result is
\begin{equation}
\theta_*^{\rm ISST}({\rm C3})/\theta_*^{\rm Planck} = 0.99174
\quad (-0.83\%),
\end{equation}
a $14\times$ reduction from the earlier $-12\%$ analytical estimate
(superseded as a recipe artefact; see Sec.~\ref{sec:dropped}).

\textbf{F95 retraction.} The F95 mixed-$\alpha(z)$ recipe is
withdrawn as the committed value. F95 Scenario~B (wall $\alpha=2/3$
pre-recombination, returning $-0.77\%$), which F95 itself dismissed
as ``not derivationally justified,'' is independently confirmed by
the Boltzmann hierarchy at $-0.83\%$. Two sympy audits underpin the
result. Step~5: the photon--baryon sound speed
$c_s^2 = c^2/[3(1+R)]$ with $R = 3\rho_b/(4\rho_\gamma)$ is unchanged
by the $(1+f)$ coupling, which enters only the Poisson source, not
the fluid EOM, since $T^{\rm eff}_{\mu\nu} = (1+f)T_{\mu\nu}$ under
uniform $f$ implies $\nabla_\mu T_{\mu\nu}=0$ identically. Step~6:
$\Psi$ is frozen on perturbed radiation at all orders in linear
perturbation theory, because $T = -\rho + 3P$ vanishes exactly for
$P = \rho/3$ at background, linear, and nonlinear order; the
recombination redshift $z_* \approx 1090$ is therefore inherited
from RECFAST unmodified.

\textbf{Independent code cross-check.} CLASS v3.x via
classy~\cite{Lesgourgues2011}, built locally with WinLibs MinGW
gcc~15.2 (toolchain notes in F134), returns $100\theta_s = 1.03188$
on the same C3 configuration---agreeing with CAMB's $1.03258$ to
$\textbf{0.07\%}$. Two independent Boltzmann codes confirm the
result; the first-pass $-0.83\%$ to $-0.92\%$ residual is robust at
the level of code precision \emph{relative to the proxy
configuration}, but not relative to the committed Wiltshire
$H(z)$ (see Proxy systematic below).

\textbf{Proxy systematic.}
The $(w_0, w_a)$ parameterisation of the late-time expansion is
calibrated against the committed Wiltshire two-domain $H(z)$ at
low redshift (where SN-Ia residuals fix the dressing) but diverges
from the Wiltshire expansion by $+76\%$ at $z = z_* \approx 1090$
(F140). A back-of-envelope propagation of this divergence through
$\theta_* = r_s/D_A = \int_0^{z_*}\!c_s\,dz/H(z)\,/\,\int_0^{z_*}\!c\,dz/H(z)$
indicates that $r_s$ is sensitive to $H$ near $z_*$ (because $c_s$
peaks just before recombination), while $D_A$ is dominated by
late-time integration. The two integrals partially cancel under a
common $H$-rescaling but not exactly: a $+76\%$ enhancement of
$H$ near $z_*$ shifts the ratio at the $\sim 5$--$10\%$ level. The
sub-percent headline $-0.83\%$ residual therefore sits inside the
proxy systematic, not above it. \emph{The committed ISST prediction
for $\theta_*$ at the precision of the Planck measurement requires
a Boltzmann code with the two-domain background substituted directly
for $H(z)$;} this computation is identified as the next critical
step (Sec.~\ref{sec:future}). Until completed, $\theta_*$ should
be reported as ``consistent with Planck within $\sim 5$--$10\%$
proxy systematic, sub-percent at the level of the proxy itself.''

\textbf{Strict uniform-$f$ at the perturbation level: falsified.}
A surgical modification of CLASS \texttt{source/perturbations.c}
implements the literal action-level uniform-$f$ prediction: the
$(1+\fprim)$ extra-matter contribution to $\delta\rho$ tracks the
\emph{baryon} perturbation $\delta_b$ rather than evolving as a
collisionless component. This is the direct consequence of
$\delta((1+f)\rho_b) = (1+f)\delta\rho_b$ under uniform $f$. The
result is dramatic and unfavourable:
\begin{equation}
H_3/H_1 \;=\; 0.013\quad{\rm (Planck:\;0.443)},\quad
\sigma_8 \;=\; 0.31\quad{\rm (Planck:\;0.811)},
\end{equation}
with $\ell_1 = 266$ (vs.\ Planck $\sim 220$). The third-peak
suppression is a factor $\sim 33$, far outside any observational
allowance. The mechanism is transparent: photon-coupled
perturbations oscillate during radiation era rather than growing,
so the gravitational wells the photon-baryon fluid oscillates in at
recombination are too shallow to drive the third compression
correctly. \textbf{Strict uniform-$f$ ISST at the perturbation
level is ruled out at $>10\sigma$ by the Planck third peak.}

\textbf{Environment-dependent $f$: perturbation closure.}
The A05 operator's linear response to density perturbations has
been derived (Sec.~\ref{sec:env_dep_f_closure}). The perturbation
$\delta f = \beta(z)\,\delta_b$ arises from the response of the
local escape velocity to the gravitational potential. Two readings
of the operator give $\beta = +0.54$ to $+1.08$ (Reading~B,
local-collapse F119 calibration) and $\beta = -0.43$ to $-0.87$
(Reading~A, horizon-scale, relativistically correct). Neither
generates a free-growing collisionless channel: the $f$-enhancement
inherits photon-coupling from the baryon perturbation pre-decoupling.
The strict uniform-$f$ source is falsified at the third peak
($>10\sigma$, Config~2 above); the linearised env-dep-$f$ closure
does \emph{not} structurally rescue the peak-height ratios. A
first-pass linear-scaling estimate shifts $H_3/H_1$ from
F134's $0.013$ (strict uniform-$f$) to $\approx 0.015$ (Reading~B)
or $\approx 0.012$ (Reading~A); the factor-$\sim 30$ gap to the
Planck-observed $\sim 0.44$ persists.

The CAMB Path~A treatment---modelling $(1{+}\fprim)\Omega_b$ as
collisionless matter in the perturbation equations---reproduces
sub-percent $\theta_*$ and $\Lambda$CDM-consistent peak ratios.
This is an empirical proxy whose physical justification within F01
is not yet established. A candidate resolution path is a $V(\Psi)$
potential extension (F126, Path~$\gamma$).

\textbf{Caveat on $V(\Psi)$ as a structural change.}
Adding a potential $V(\Psi)$ to the action gives $\Psi$ a
non-trivial Euler--Lagrange equation,
$R/16\pi + V'(\Psi)/16\pi = 0$ at fixed $f$, which in general
\emph{makes $\Psi$ propagate}: ADM analysis would then count an
additional scalar mode, breaking the F128 ``zero scalar DOF''
result that supports the no-fifth-force claim
(Sec.~\ref{sec:fifth}). It would also alter the Branch theorem (the
RHS of $R\nabla\Psi = 0$ acquires a $V'$-dependent term from the
modified Bianchi projection), the $c_{\rm GW} = c$ argument (a
propagating $\Psi$ can mix with the gravitational sector at
non-trivial $V''$), the PPN computation (the propagating mode
generically contributes $\gamma \neq 1$ unless screened, and ISST
denies all standard screening mechanisms in its Passport), and the
LLR $\dot G/G$ bound (a $V$-driven slow evolution of $\Psi$
contributes directly). \emph{We therefore do not claim a CMB
resolution via $V(\Psi)$ in this paper.} Path~$\gamma$ is identified
as a separate-theory direction whose claim to ``preserving the F01
results'' has to be earned property-by-property; until that is
done, ISST's commitment is to the F01 action without $V(\Psi)$,
with the third-peak gap reported as a live falsifier
(Sec.~\ref{sec:falsification}). The honest position is that ISST
either fits the CMB through a yet-undiscovered F01-internal
mechanism (the $\delta f$ closure does not provide it) or the third
peak is a genuine falsification of strict-F01.

\textbf{Stage~G updates from F134.} Five framework verdicts
change. \textsf{boltzmann\_hierarchy\_cmb} (FRW preconditions)
remains \textsc{[denied]} on F01.
\textsf{isst\_boltzmann\_two\_domain\_w0wa\_proxy} (NEW) is
registered \textsc{[allowed]} as the published-prediction
framework. \textsf{isst\_boltzmann\_full\_uniform\_f} (NEW) is
registered \textsc{[denied]}, falsified by the third peak.
\textsf{isst\_boltzmann\_env\_dep\_f} (NEW) is registered
\textsc{[conditional]} pending implementation. The remaining caveat
on Path~A is the high-$z$ extrapolation: the effective $(w_0,w_a)$
proxy, extracted from low-$z$ Wiltshire matching, diverges from the
committed Wiltshire dressed $H(z)$ by $+76\%$ at $z = 1090$ (F140).
Moreover, the committed ISST matter density
$\Omega_m = (1{+}\fprim)\Omega_b = 0.329$ exceeds the Wiltshire bare
matter density $\Omega_{m,{\rm bare}} = 0.125$ by a factor $2.6$;
this mismatch causes the effective dark-energy density inferred
from Friedmann inversion to become \emph{negative} at $z > 2$,
precluding any tabulated $w(z)$ input to standard Boltzmann codes
at the committed parameters. A true Wiltshire+ISST Boltzmann
computation requires modifying the background solver to accept the
two-domain volume-averaged expansion directly, not through a
Friedmann proxy. The $-0.83\%$ result is therefore a proxy-level
estimate; the systematic uncertainty from the $+76\%$ $H(z)$
divergence at recombination is unquantified and could shift
$\theta_*$ by several percent in either direction.

Status: \DEMO\ at proxy level for the acoustic angle (F134,
CAMB+CLASS), with a $+76\%$ $H(z)$ proxy systematic at $z = 1090$
(F140) that is unquantified at the observable level. The third-peak
height $H_3/H_1$ remains \OPEN\ (factor-$\sim 30$ residual gap to
Planck; F138 confirmed the linear env-dep-$f$ closure does not
rescue it). Env-dep-$f$ closure form \PROVEN\ (F138,
Sec.~\ref{sec:env_dep_f_closure}); env-dep-$f$ \emph{rescue} of the
third-peak gap \OPEN\ within F01; resolution escalates to a
$V(\Psi)$ extension (F126 Path~$\gamma$, under investigation) or to
the honest-flag position the paper takes. A Boltzmann code with the
true Wiltshire two-domain background is identified as the most
important remaining computation (Sec.~\ref{sec:future}).

%==========================================================
\section{Early universe and gravitational waves}
\label{sec:early}

\subsection{BBN}

Big-Bang Nucleosynthesis is inherited intact. The Branch-A radiation
closure (F95) requires $\Pi_{\rm rad}=-\rho_r/2$, a sympy-derived
consistency condition; physical tight-coupling delivers $\leq 10^{-3}$,
satisfying it within numerical noise. Crucially, $\Psi$ is frozen on a
radiation background ($T=0$, traceless source) so $G$ does not run
during nucleosynthesis. Standard BBN abundances ($\eta_B$,
$Y_{\rm He}$, lithium) are inherited. Status: \PROVEN. The earlier
F113 BBN catastrophe ($G_{\rm BBN}/G_{\rm today}\approx 3000$) was a
wrong-background artefact: the Schutz/Hawking convention on a
naively-FRW background gave the catastrophe; the F95 closure
on the correct radiation-dominated Branch-A removes it.

\subsection{Gravitational waves}

Gravitational waves propagate at $c$. The action \eqref{eq:action}
is linear in $\Psi$ in the gravitational sector, so tensor
perturbations $h_{\mu\nu}$ decouple from $\Psi$ fluctuations at
linear order on any background:
\begin{equation}
\Box h_{\mu\nu} = -\frac{16\pi}{\Psi}\tau_{\mu\nu},
\quad c_{\rm GW}=c \text{ structurally}.
\end{equation}
The constraint $|c_{\rm GW}/c-1|<10^{-15}$ from
GW170817~\cite{LIGO2017} is satisfied trivially.
F128 ADM analysis confirms two tensor degrees of freedom and no
scalar mode: LIGO measures $+$ and $\times$ polarisations only;
no breathing mode. Status: \PROVEN.

\subsection{No scalar dipole radiation}

Because $\Psi$ does not propagate, binary inspirals do not radiate
into a scalar channel. Hulse--Taylor quadrupole-only behaviour is
inherited; the framework \textsf{scalar\_dipole\_radiation\_binary}
is \textsc{[denied]} as inapplicable to F01. Status: \PROVEN.

%==========================================================
\section{Falsification commitments}
\label{sec:falsification}

This section is the most important section in the paper. ISST is
science only to the extent that it commits to falsifiers in advance,
and those commitments include the predictions ISST has
\emph{retracted}.

\subsection{Hard falsifiers}

Each of the following observations would, by itself, kill ISST as
formulated here.

\begin{itemize}
\item \textbf{$c_{\rm GW}\neq c$.} ISST sits in the Horndeski corner
that GW170817 left intact, structurally not by tuning. Any future
$|c_{\rm GW}/c-1|>10^{-15}$ at high-confidence kills the theory.
\item \textbf{Fifth-force detection.} F01 has no propagating scalar
and no fifth force to screen. \emph{Specifically:} a
composition-dependent acceleration anomaly inconsistent with WEP
at the level the universal $(1{+}f)\Lm$ coupling permits, or a
Yukawa-type force on test particles with finite range
distinguishable from $1/r^2$ gravity at any laboratory or
solar-system precision, would falsify ISST. An observed spatial
variation of $G_{\rm eff}$ correlating with local matter content
would by contrast be \emph{consistent with} the $(1{+}f)$ mechanism
(cf.\ the big-$G$ speculation in Sec.~\ref{sec:future}) and would
not constitute a fifth force in the ISST sense; it would be a
detection of the structural matter coupling, not a refutation of
its absence.
\item \textbf{$\gamma_{\rm PPN}\neq 1$ at $\gtrsim 10^{-5}$.}
Tightening the Cassini bound below ${\sim}10^{-5}$ would force a
path commitment to a propagating-scalar variant or falsify F01.
\item \textbf{BBN abundances inconsistent with SBBN.} ISST inherits
SBBN intact; any confirmed deviation at the level of an evolving
$G$ during nucleosynthesis falsifies the F95 closure.
\item \textbf{Direct dark-matter particle detection.} ISST
identifies $(1{+}\fprim)\Omega_b$ as the entirety of $\Omega_m$,
predicting persistent nulls in dark-matter direct-detection
experiments at the CDM density across all currently searched
parameter spaces. Specifically: WIMP detection at
$\sigma_{\rm SI}\gtrsim 10^{-47}\,{\rm cm}^2$ over $1$--$1000$~GeV
(within current XENONnT/LZ reach to the neutrino floor); QCD-axion
detection at $g_{a\gamma\gamma}\gtrsim 10^{-15}\,{\rm GeV}^{-1}$
over $1$--$100~\mu$eV (ADMX/HAYSTAC reach); or keV sterile-neutrino
detection inconsistent with all-baryon
$(1+\fprim)\Omega_b = \Omega_m$ at $>3\sigma$ would each falsify
\eqref{eq:A04}. Current
nulls~\cite{XENONnT2023,LUXZEPLIN2023,ADMX2021} are consistent with
the prediction; the kill condition is a $>5\sigma$ confirmed
detection at any of these specifications.
\end{itemize}

\subsection{Cumulative tensions and live falsifiers}

Each of the following is a multi-observation falsifier with an
explicit kill condition:

\begin{itemize}
\item \textbf{Bullet/merging-cluster lensing under Euclid.} The
lighthouse mechanism predicts $\kappa$ tracks bulge-dominated
stellar mass, not gas mass. If Euclid cluster tomography shows the
$\kappa$ peak tracking gas mass (not stars) at $>3\sigma$ across
multiple mergers, the mechanism is falsified.
\item \textbf{$f\sigma_8$ shape under DESI~DR2.} If DR2+Euclid Y1
find $f\sigma_8(z)$ within $1\sigma$ of $\Lambda$CDM and ISST misses
by $>2\sigma$ in $>3$ bins, wall-dynamical $f\sigma_8$ is falsified.
\item \textbf{Void-lensing $G$-ratio under Euclid.} ISST predicts
$G_{\rm void}/G_{\rm local}=1.025$ (a $+2.5\%$ enhancement). If
Euclid/LSST measure $G_{\rm void}/G_{\rm local}=1.000\pm 0.005$
(GR-consistent), the differential prediction is falsified, forcing
either a $V(\Psi)$ chameleon path or retirement.
\item \textbf{MOND $a_0$ tightening.} ISST predicts an $11\%$ offset
of $a_{\rm crit}$ from MOND $a_0$. If $a_0$ is measured to better
than $5\%$ with a central value within $2\%$ of
$1.20\times 10^{-10}$, the F80 prediction is falsified.
\item \textbf{Jellyfish-galaxy lensing (ISST vs.\ MOND
discriminator).} In ram-pressure-stripped jellyfish galaxies the
stellar disk and the stripped gas tail experience the same Newtonian
acceleration but have different processing histories: ISST predicts
$f_s\to 0$ in the shocked gas tail (suppressed enhancement,
$\kappa\to (1+f_{\rm prim})\,\Sigma_{\rm gas}$ only) and
$f_s\sim$few in the disk (full enhancement). MOND/QUMOND predicts
equal phantom-density contribution from both components at the
local acceleration, modulated only by compactness. Stacked weak-lensing
of jellyfish samples in clusters with kinematic disk-versus-tail
separation (Euclid + ground-based spectroscopy) discriminates: if the
disk-to-tail $\kappa$ ratio per unit baryonic surface density agrees
with the QUMOND compactness-only prediction within $1\sigma$, ISST's
processing-history claim is falsified.
\item \textbf{Peak-height ratios under env-dep-$f$ full $C_\ell$.}
F134 surgically modified CLASS to test strict uniform-$f$ at the
perturbation level---the literal action prediction with
$\delta((1+f)\rho_b) = (1+\fprim)\delta\rho_b$, photon-coupled.
The result is incompatible with Planck: $H_3/H_1 = 0.013$ vs
observed $\sim 0.44$ (factor $33$), $\sigma_8 = 0.31$ vs $0.81$.
Strict uniform-$f$ at the perturbation level is therefore
\textbf{ruled out at $>10\sigma$} and is retracted.

\textbf{F138 update.} F138 derives the env-dep-$f$ linear closure
$\delta f = \beta(z)\delta_b$ explicitly from A05
(Sec.~\ref{sec:env_dep_f_closure}). The closure is a \emph{few-percent
correction to the gravitational source}, not a free-growing
collisionless channel: $\delta f$ is a functional of $\delta_b$
through the local $\dKL$ response and inherits photon-coupling
pre-decoupling. A first-pass linear-scaling estimate shifts $H_3/H_1$
from $0.013$ (strict uniform-$f$) to $\approx 0.015$ (Reading B) or
$\approx 0.012$ (Reading A); the factor-$\sim$30 residual gap to Planck
persists. The previous draft's claim that the env-dep-$f$ closure
``restores collisionless-like growth'' is downgraded.

\textbf{Live falsifier (revised).} The third-peak factor-$\sim$30
gap is \OPEN\ within F01: the F138 closure correction is bounded at
$\pm 16\%$ of the CAMB Path~A baseline, so a full direct CLASS
implementation of the $\beta(z)\delta_b$ source modification cannot
close more than a percent or so of the gap. The kill condition
shifts as follows. (a)~If a direct CLASS implementation of F138's
$\beta(z)\delta_b$ source confirms a peak shift in the F138-predicted
$\pm 16\%$ window, ISST's commitment is that the residual factor-$\sim$30
gap is genuine and the resolution escalates beyond F01: either
$V(\Psi)$ extension (F126 Path~$\gamma$) or the honest-flag position
of Sec.~\ref{sec:theta}. (b)~If the direct implementation produces
a much larger correction ($\gtrsim$ factor of order 10), the F138
derivation is incomplete and the closure has more content than
identified---this would itself be a substantial finding. (c)~If the
$V(\Psi)$ extension is pursued and fails to recover Planck-consistent
peak heights at percent level, the extension is falsified and the
honest-flag position becomes the only path. ETA on the F138 direct
CLASS implementation: 1--2 weeks (toolchain in F134 reproducible);
working on. The status tag for the third-peak gap is therefore
\OPEN.
\end{itemize}

\subsection{Falsifier table: timeline, precision, observatory}
\label{sec:falsifier_table}

The falsifiers are summarised in Table~\ref{tab:falsifiers} with
the timescale for the relevant data release, the target precision
required to trigger the falsification, and the observatory or
analysis pipeline through which the test arrives. Entries are
ordered by approximate time-to-decision under nominal mission
schedules; ``current'' indicates that the test can in principle be
run on already-released data.

\begin{table*}[t]
\caption{ISST falsifier inventory. ``Hard'' falsifiers are
single-observation kill conditions; ``Live'' falsifiers are
multi-observation kill conditions with explicit thresholds. ETA is
the approximate window from this paper's writing within which the
relevant data become available at the target precision.}
\label{tab:falsifiers}
\begin{ruledtabular}
\begin{tabular}{lllll}
Falsifier & Class & ETA & Target precision & Observatory / pipeline \\
\hline
Bullet $\kappa_{\rm gal}/\kappa_{\rm gas}$ ratio
& Live & current
& per-pixel match to $3.76$ at $>3\sigma$
& Cha~2025 / Rihtar\v{s}i\v{c}~2026 / JWST + HST \\
$f\sigma_8$ shape at $z\sim 1.1$
& Live & 1--2 yr
& $13\sigma$ ISST-vs-$\Lambda$CDM at DR2 high-$z$ bins
& DESI~DR2 + Euclid~Y1 \\
Void-lensing $G_{\rm void}/G_{\rm local}$
& Live & 2--3 yr
& $\pm 0.5\%$ at $r > 50$~Mpc voids
& Euclid + LSST/Rubin \\
$a_0$ central value vs.\ $a_{\rm crit}$
& Live & 3--5 yr
& $a_0$ to $\pm 5\%$ with central value $\pm 2\%$
& High-precision RAR / DESI rotation-curve campaigns \\
Third-peak ratios under env-dep-$f$
& Live & 1--2 wk (compute)
& $H_3/H_1$ within F138 $\pm 16\%$ window or beyond
& Modified CLASS direct implementation (F138) \\
$\theta_*$ on committed two-domain background
& Live & 2--4 mo (compute)
& sub-percent agreement with Planck
& Modified Boltzmann code on Wiltshire~$H(z)$ \\
Jellyfish disk-vs-tail $\kappa$
& Live & 3--5 yr
& disk-to-tail ratio per unit $\Sigma_{\rm bar}$ vs.\ QUMOND
& Euclid + ground-based spectroscopy \\
$c_{\rm GW}\neq c$
& Hard & current
& $|c_{\rm GW}/c-1|>10^{-15}$ at high confidence
& LIGO/Virgo/KAGRA + EM counterpart \\
$\gamma_{\rm PPN}-1$ tightening
& Hard & 5--10 yr
& Cassini bound below $\sim 10^{-5}$
& Mars/asteroid radio-tracking; future Solar-System probes \\
Direct dark-matter detection
& Hard & 3--7 yr
& WIMP at $\sigma_{\rm SI}\gtrsim 10^{-47}$~cm$^2$
& XENONnT / LZ Run~3+ \\
& & & QCD-axion at $g_{a\gamma\gamma}\gtrsim 10^{-15}$~GeV$^{-1}$
& ADMX / HAYSTAC \\
& & & keV sterile-$\nu$ inconsistent with all-baryon $\Omega_m$
& X-ray line searches \\
BBN abundance deviation
& Hard & current
& any confirmed $\Delta G_{\rm BBN}/G$ at SBBN level
& Light-element abundance reanalyses \\
Composition-dependent EP violation
& Hard & current--5 yr
& any anomaly inconsistent with universal $(1{+}f)\Lm$
& MICROSCOPE follow-ups; STEP~+~atom interferometry \\
\end{tabular}
\end{ruledtabular}
\end{table*}

\subsection{Results ISST has dropped}
\label{sec:dropped}

ISST has dropped its own claims in three categories during the
course of this work. We list them here for transparency, because
a theory that drops its failures is more credible than one that
hides them.

\begin{itemize}
\item \textbf{$S_8$ tension: \DROPPED.} F127 audited the ISST
$S_8$ prediction and found the wrong sign:
$G_{\rm wall} < G_{\rm void}$ at the relevant scales weakens
late-time lensing, predicting a wider $S_8$ than observed. The
ISST prediction $S_8\approx 0.824$ is in tension with the
KiDS-1000~\cite{KiDS1000} reading $\sim 0.760$ (with consistent
DES~Y3~\cite{DESY3} and HSC~Y3~\cite{HSCY3} preferences). We
therefore drop $S_8$ as an ISST claim. (The $S_8$ tension itself is fading as KiDS Legacy
revises downward, independently of ISST.)
\item \textbf{Halo-ontology Bullet attempts: \DROPPED.} F65, F66, and
F68 attempted to resolve the Bullet Cluster by assigning per-galaxy
halo masses, importing scalar-field stress-energy halos, or fitting
halo-attached directional KL coefficients. All three violated F01's
pointwise $(1+f)\rho$ source structure. The Derivation Passport
explicitly denies \textsf{dark\_matter\_halo\_ontology}; these
attempts have been retracted, and the F131 lighthouse closure
(Sec.~\ref{sec:bullet}) is the F01-native treatment.
\item \textbf{F113 BBN catastrophe: \DROPPED.} F113 reported
$G_{\rm BBN}/G_{\rm today}\approx 3000$ on a naively-FRW background,
violating BBN by $\sim 10^4$. F112a/F113a audits identified the
error as wrong-background contamination; the F95 Branch-A
radiation closure resolves it. The earlier finding has been
retracted.
\item \textbf{F95 $-12\%$ $\theta_*$ residual: \DROPPED.} F95
reported a mixed-$\alpha(z)$ analytical estimate of
$\theta_*/\theta_*^{\rm Planck} = 0.916$. F134 ran a first-pass
Boltzmann hierarchy in CAMB and CLASS on the Wiltshire-effective
background and returned $0.992$ / $0.991$ ($-0.83\%$ / $-0.92\%$),
$14\times$ smaller, with the two independent codes agreeing to
$0.07\%$. The F95 analytical recipe is artefactual and the larger
number is retracted. F95 Scenario~B ($-0.77\%$), which F95 itself
dismissed as ``not derivationally justified,'' is independently
confirmed by both Boltzmann integrators (Sec.~\ref{sec:theta}). The
Branch-A radiation closure $\Pi_{\rm rad} = -\rho_r/2$ that F95
derived for the BBN and recombination arguments is independent of
the retracted analytical recipe and remains load-bearing.
\item \textbf{Strict uniform-$f$ at perturbation level: \DROPPED.}
F134 modified CLASS \texttt{source/perturbations.c} to implement
the literal action-level prediction $\delta((1+f)\rho_b) =
(1+\fprim)\delta\rho_b$ photon-coupled. Result:
$H_3/H_1 = 0.013$ vs Planck's $0.44$ (factor $33$ mismatch),
$\sigma_8 = 0.31$ vs $0.81$, $\ell_1 = 266$ vs $220$. Photon-coupled
matter does not grow during radiation domination, so the
gravitational wells at recombination are too shallow to drive the
third compression. The strict uniform-$f$ reading at the
perturbation level is incompatible with Planck at $>10\sigma$ and
is retracted. The committed Path~A prediction continues to use a
collisionless-CDM proxy for $\fprim\Omega_b$ at the perturbation
level; F138 has shown the linear env-dep-$f$ closure provides only
a few-percent correction to that proxy and does not derive it from
F01 (Sec.~\ref{sec:theta}).

\item \textbf{Smooth-crossover $\fprim$ reduction: \DROPPED.}
The speculation that replacing sharp-transition step functions with
smooth Standard-Model crossovers would reduce $\fprim$ from $5.664$
toward the observational target $5.40$ has been tested (F139). A
per-transition smooth evaluation returns $\fprim = 5.78$
(\emph{larger}, not smaller); the smooth path integral
$\int (1/g_*)\,|dg_*/dT|\,dT$ is path-independent and gives $3.31$
regardless of crossover width. The $1.82\sigma$ residual is a
genuine feature of the DOF sum, not an approximation artefact. The
operator-level derivation from $(1{+}f)\Lm$ on the full thermal
history remains \OPEN\ as a separate programme
(Sec.~\ref{sec:future}).
\end{itemize}

%==========================================================
\section{Comparison with existing frameworks}
\label{sec:comparison}

Table~\ref{tab:comparison} summarises ISST against the principal
alternatives. The vertical axis lists capability, the horizontal
axis lists frameworks; entries are derived (D), assumed (A), absent
(--), or excluded (\texttimes).

\begin{table*}[t]
\caption{ISST vs.\ existing frameworks. D = derived, A = assumed
parameter or input, -- = no native treatment,
\texttimes\ = excluded.}
\label{tab:comparison}
\begin{ruledtabular}
\begin{tabular}{lccccc}
Capability & ISST & $\Lambda$CDM & MOND/TeVeS & metric $f(R)$ & Horndeski \\
\hline
$\Omega_m/\Omega_b\approx 6$ & D & A & -- & A & A \\
Late-time acceleration & D (geometric) & A ($\Lambda$) & -- & D ($f''$) & D ($G_4,G_5$) \\
Rotation curves & D ($a_{\rm crit}$) & D (NFW fit) & D ($a_0$) & D & D \\
Cluster lensing & D (lighthouse) & D (NFW DM) & \texttimes & A & D \\
CMB peak structure & $\sim$ (proxy sub-\%, peak ratios open; $V(\Psi)$ under inv.) & D (6-param fit) & -- & D & D \\
PPN $\gamma=\beta=1$ & D (exact) & D & D & A (screening) & A (screening) \\
Fifth-force screening & not needed & not needed & not needed & required & required \\
$c_{\rm GW}=c$ & D (linearity) & D & D & D & D (corner) \\
DOF count & 2 & 2 & 2+2 & 2+1 & 2+1 \\
Parameter count: derived from action & 1 ($a_{\rm crit}$) & 0 & 0 & 0 & 0 \\
Parameter count: SM-input ansatz & 1 ($\fprim$, chain open) & 0 & 0 & 0 & 0 \\
Parameter count: empirical & 7 (see below) & 6 ($\Lambda$CDM) & 1 ($a_0$) & 1 ($f$) & 4+ \\
Parameter count: inherited & 1 ($f_{v0}$, Wiltshire) & --- & --- & --- & --- \\
Initial conditions & 2 ($\Omega_b h^2,\Psi_0$) & 2 & yes & yes & yes \\
\end{tabular}
\end{ruledtabular}
\end{table*}

\textbf{Honest parameter accounting.}
ISST's full ISST-specific parameter set, with the lighthouse working
memories $f_s^\star, f_s^{\rm gas,shock}$ now counted explicitly:

\emph{Parameter-free predictions from action + SM (2):}
\begin{itemize}
\item $a_{\rm crit} = 1.07\times 10^{-10}$~m/s$^2$ ---
derived from the wall-Friedmann plus Branch-A weak-field algebra
(Sec.~\ref{sec:acrit}); depends on the committed $H_0 = 61.79$
and $\fprim = 5.66$.
\item $\fprim = 5.66 \pm 0.12$ --- discrete-sum ansatz over SM
freezeouts (Sec.~\ref{sec:fprim}); $1.6$--$1.8\sigma$ tension
with Planck. The formal chain from the action's KL functional to
the discrete sum is structurally open (four candidate paths
ruled out); $\fprim$ is therefore properly framed as an
SM-motivated ansatz, not as a closed-chain derivation.
\end{itemize}

\emph{Empirical (fitted to data) (7):}
\begin{itemize}
\item $a_0^{\rm reg} = 0.01\times 10^{-10}$~m/s$^2$ --- rotation-curve
regulator internal scale, fitted on a 14-galaxy SPARC subsample
(Sec.~\ref{sec:galaxy}).
\item $q \approx 0.45$ --- regulator exponent, same subsample.
\item $\alpha = 0.869$ --- KL coupling, galactic calibration on the
14-galaxy subsample (the calibration procedure: for each trial
$\alpha$, compute $V_{\rm ISST}(r)$ for all 14 galaxies via
\eqref{eq:poisson} with the observed baryon distribution and
minimise the median RMS velocity residual). Note that
$\alpha_{\rm cosmo} \approx 1.71$ is required to match the
discrete-sum $\fprim$ from the smooth path-integral form
\eqref{eq:fprim_integral}, a factor-$\sim 2$ scale dependence
acknowledged as an open tension (Sec.~\ref{sec:fprim}).
\item $n = 4.83$ --- void-fraction sigmoid shape parameter in the
$\chi_w(z)$ closure (Sec.~\ref{sec:hubble}); residual, not
derived.
\item $K_{\rm env} = 1.484$ --- environmental rescaling
recovered as the squared ratio of two independent Hubble outputs
in the F61 closure on the Sheth--van~de~Weygaert two-barrier
abundance (Sec.~\ref{sec:tracker}, App.~\ref{app:h0_frame});
propagates into the committed
$H_0 = 73.82/\sqrt{K_{\rm env}} = 60.61$~km/s/Mpc dressed value
and thence into $a_{\rm crit}$. This is an empirical input rather
than an action-derived prediction; deriving $K_{\rm env}$ from
earlier ISST-committed quantities is identified as \OPEN.
\item $f_s^\star \approx 3.55$ --- lighthouse stellar working memory,
fitted to the Paraficz~\emph{et~al.}~(2016) local-column
convergence ratio at the Bullet (Sec.~\ref{sec:bullet}).
Derivation from violent-relaxation working memory open.
\item $f_s^{\rm gas,shock} \approx 0.21$ --- lighthouse
shocked-gas working memory, fitted to the same Paraficz
decomposition.
\end{itemize}

\emph{Inherited (1):}
\begin{itemize}
\item $f_{v0} = 0.762$ --- void fraction today, from the Wiltshire
two-domain framework~\cite{Wiltshire2009}.
\end{itemize}

\emph{Initial conditions (shared with all theories) (2):}
$\Omega_b h^2 = 0.02237$ from BBN/Planck and $\Psi_0 = 1/G_N$ from
matching to laboratory $G$.

\emph{Total ISST-specific count: 10} (2 derived/ansatz + 7 empirical
+ 1 inherited), plus 2 initial conditions. \emph{This is more total
parameters than $\Lambda$CDM's 6.} The honest framing is that ISST
derives \emph{two} of these from first principles (one fully from
the action, one from SM thermal-history input with the chain to the
action open) while $\Lambda$CDM fits all six of
$(\Omega_b, \Omega_c, H_0, \tau, A_s, n_s)$ from CMB data. The
epistemic status differs even when the count does not favour ISST
on raw count alone --- ISST's two parameter-free predictions are
\emph{predictions}, in advance of observational fitting; $\Lambda$CDM's
six are \emph{posteriors}.

The derivation of the regulator exponent $q$ from second-order
Branch-A perturbation theory, a single-$\alpha$ scale-bridging
calibration that closes the factor-$\sim 2$ galactic-vs-cosmological
gap, the derivation of $K_{\rm env}$ from earlier ISST-committed
quantities (it currently enters as an empirical squared-Hubble-ratio
input rather than as an action-derived prediction), the lighthouse
$f_s$ values from violent-relaxation working memory, and the
$\fprim$ chain closure are identified as \OPEN.

\textbf{vs.\ $\Lambda$CDM.} ISST's empirical regulator parameters
$(a_0^{\rm reg}, q)$ are fitted on rotation curves only and do not
propagate into the cosmological prediction. The framing differs:
$\Lambda$CDM \emph{assumes} the dark sector and asks ``what
density?''; ISST proposes $(1{+}\fprim)\Omega_b = \Omega_m$ as the
cosmological-floor identity and asks ``what observation could break
it?''.

\textbf{ISST is observationally indistinguishable from
$\Lambda$CDM at the CMB level under the env-dep-$f$ closure.} F134
established that the strict uniform-$f$ perturbation prediction is
ruled out by the Planck third peak, and the env-dep-$\delta f$
contribution that ISST already requires for the $H_0$ tracker
restores LCDM-like clustering at recombination
(Sec.~\ref{sec:theta}). The CMB therefore neither favours nor
disfavours ISST relative to $\Lambda$CDM at first-pass precision.
The distinguishing tests live entirely in the late-time
structure-formation sector: the differential $a_{\rm crit}$
prediction at galactic scale ($-11\%$ vs.\ MOND $a_0$;
Sec.~\ref{sec:acrit}); the $f\sigma_8$ shape prediction at DESI~DR2
+ Euclid~Y1 precision ($13\sigma$ distinguishability at $z=1.1$;
Sec.~\ref{sec:cosmology}); the $+2.5\%$ void-lensing $G$ ratio
versus GR (Sec.~\ref{sec:cluster}); and the lighthouse-mechanism
predictions for cluster-merger geometry (Bullet, Abell~520, and
the jellyfish-galaxy ISST-vs.-MOND discriminator;
Sec.~\ref{sec:falsification}). The paper's claim that ISST
``replaces three components of the standard model'' therefore
reduces, at the CMB level, to a reinterpretation of what
$\Lambda$CDM's $\Omega_m$ and $\Omega_\Lambda$ \emph{are}, not to
a numerically different prediction; the substantive distinction
is at late times.

\textbf{vs.\ MOND/TeVeS}~\cite{Milgrom1983,Bekenstein1984,BekensteinTeVeS2004,SandersMcGaugh2002,FamaeyMcGaugh2012,BanikZhao2022}.
MOND's $a_0$ is empirical; ISST's $a_{\rm crit}$ is derived.
ISST has cosmology native to the action; MOND does not.
ISST's $a_{\rm crit}$ sits $11\%$ below $a_0$, a sharp discriminator.

\textbf{vs.\ metric $f(R)$.} Both modify gravity; metric $f(R)$
introduces a propagating scalar that requires chameleon
screening. ISST's Palatini-class structure has a Lagrange-multiplier
$\Psi$ that does not propagate; no screening is needed because
there is no fifth force. The framework
\textsf{f\_R\_metric\_equivalence} is \textsc{[denied]} for F01.

\textbf{vs.\ Horndeski}~\cite{Horndeski1974,Creminelli2017,LIGO2017,LIGOGW170817EM}.
ISST sits in the Horndeski corner that GW170817 spared. It is not a generic Horndeski theory; the
specific $\Psi R$ form with no kinetic term is a structural
choice that survives every cosmological-tensor constraint
applied to the broader class.

\textbf{vs.\ backreaction.} Wiltshire's two-domain
framework~\cite{Wiltshire2007,Wiltshire2009} is a special case of
Buchert averaging~\cite{Buchert2000,BuchertRasanen2008} that ISST
adopts as the inhomogeneous background; ISST adds the action
principle that requires \emph{some} inhomogeneous background as
the only Branch-A-compatible dust cosmology, with Wiltshire's
two-domain form the minimal physically motivated realisation
(Sec.~\ref{sec:wiltshire_forced}).

\textbf{vs.\ Brans--Dicke.} BD~\cite{BransDicke1961,OHanlonTupper1972,FujiiMaeda2003}
gives $\Psi$ a kinetic term and a propagating equation of motion;
its $\omega_{\rm BD}$ parameter is constrained by Cassini to
$\omega_{\rm BD} > 4\times 10^4$. ISST has no kinetic term: $\Psi$
is a Lagrange-multiplier and the BD comparison is structurally not
applicable. The thermodynamic-gravity programme of
Verlinde~\cite{Verlinde2011,Verlinde2017} and
Padmanabhan~\cite{Padmanabhan2010} shares ISST's information-as-source
intuition but does not specify a covariant action;
F01~\eqref{eq:action} is the local realisation.

\textbf{vs.\ small-scale $\Lambda$CDM stresses.} JWST early-galaxy
abundance~\cite{LabbeJWST2023} and small-scale halo
challenges~\cite{BullockBoylanKolchin2017,BoylanKolchin2023} are
problems for $\Lambda$CDM. ISST's structure-formation history
through Wiltshire backreaction may address them but the
quantitative work (full $C_\ell$, $P(k)$) is identified as
\OPEN. The cosmic SFR history~\cite{MadauDickinson2014} provides the
$f_s$-accumulation timeline that the lighthouse mechanism uses
implicitly; deriving SFR from F01 itself is not in scope here.

%==========================================================
\section{Discussion and future work}
\label{sec:future}

\subsection{Open items}

The largest open items, with kill conditions where applicable, are:

\begin{itemize}
\item \textbf{Modified Boltzmann hierarchy.} The most important
remaining computation is a Boltzmann code whose background solver
accepts the Wiltshire two-domain volume-averaged expansion directly.
The $(w_0, w_a)$ proxy used in F134 diverges from the committed
Wiltshire $H(z)$ by $+76\%$ at $z = 1090$ (F140); the committed
ISST $\Omega_m = 0.329$ exceeds the Wiltshire bare $\Omega_m = 0.125$
by a factor $2.6$, causing the equivalent dark-energy density under
Friedmann inversion to become negative at $z > 2$ and precluding
any Friedmann-template proxy at high redshift. A modified CLASS
with (i)~the two-domain background solver, (ii)~the $V(\Psi)$
perturbation extension, and (iii)~$(1{+}f)\rho_b$ as the gravitational
source would provide the definitive CMB test. The framework
\textsf{isst\_boltzmann\_two\_domain\_tabulated} is identified for
passport stamping under (i); \textsf{isst\_boltzmann\_env\_dep\_f}
remains \textsc{[conditional]} pending the $V(\Psi)$ derivation
under (ii). Estimated scope: $2$--$4$ months of code development.
\OPEN.
\item \textbf{Operator-level $\fprim$ chain.} The DOF-counting
$\fprim = 5.66\pm 0.06$ overshoots the CMB target $5.40\pm 0.15$ at
$1.6$--$1.8\sigma$. Smooth-crossover evaluation does not close this
gap (F139: PER\_TRANSITION returns $5.78$, the path-independent
integrand returns $3.31$, neither reaches $5.40$). \emph{F145
sharpens the OPEN status from ``derivation not yet done'' to
``structurally obstructed under the four canonical paths.''} The
action's $f$ in eq.~\ref{eq:f} is a single-epoch $D_{\rm KL}$
integral; the discrete sum is a path-additive event-counting
quantity over the four SM freezeouts. Single-epoch $D_{\rm KL}$ of
post-thermal-history Maxwellians gives $1.36$ (factor $\sim 4$
short); $D_{\rm KL}$ vs.\ local-bounded uniform gives $\sim 40$--$60$
(factor $\sim 8$ too large; reading is galactic-specific); per-step
transient $D_{\rm KL}$ summed gives $0.30$ (factor $\sim 19$ short;
Shore--Johnson additivity does not extend to time-separated events);
R\'enyi-2 sum at the 4-step partition gives $5.82$ accidentally but
Shore--Johnson forces $q=1$ (Sec.~\ref{sec:fprim}, F145). Closing
the chain therefore requires either (a) a path-integral extension
of the action (re-checks A05, F46 ratchet, $(1{+}f)\Lm$, the
Shore--Johnson uniqueness derivation), or (b) a demonstration that
the discrete sum emerges from a different limit of $(1{+}f)\Lm$
than the single-epoch $D_{\rm KL}$, or (c) a more abstract
information-theoretic measure for which the discrete sum is the
unique additive monotone functional. The $\sim 2$--$3\sigma$
two-$\alpha$ structural question ($\alpha_{\rm gal} = 0.869$ vs.\
$\alpha_{\rm cosmo} \approx 1.71$ for the smooth-integral form to
recover $5.66$) is identified as a propagated sub-problem. \OPEN.
\item \textbf{Matter power spectrum $P(k)$ and BAO.} Same
methodological obstacle as $C_\ell$. \OPEN.
\item \textbf{Newton-limit regulator exponent $q$.} Empirically
$\sim 0.45$; a derivation from second-order Branch-A expansion is
identified. \OPEN.
\item \textbf{$f_s^\odot$ and Cassini sub-leading.} Currently
unmeasured; a kill condition is $f_s^\odot > 0.05$.
\OPEN.
\item \textbf{Path commitment ($\zeta$ vs.\ $\gamma$ vs.\ $\varepsilon$).}
Three Reading-4-class extensions remain on the table; the choice
has consequences for void-lensing and screening. Held off pending
Lily review of the F130b path catalogue. \OPEN.
\item \textbf{Per-pixel Bullet $\kappa$ map.} The local
$\Sigma_{\rm baryon}$ contrast at the two peaks is now resolved by
JWST~\cite{Cha2025} (modified Hausdorff $20$~kpc ICL/$\kappa$
alignment) and the GR $\kappa$ map is published in
Rihtar\v{s}i\v{c}~\emph{et~al.}~(2026)~\cite{Rihtarsic2026}, with
QUMOND modelling residual $|\Delta\kappa|\lesssim 0.15$ in
Hernandez~(2026)~\cite{Hernandez2026}. Extracting
$\kappa_{\rm gal}/\kappa_{\rm gas}$ from these maps and comparing to
the operator $(1+f_s^\star)/(1+f_s^{\rm gas,shock}) \approx 3.76$
is identified as the next quantitative step
(Sec.~\ref{sec:bullet}). \OPEN\ (numerical extraction; observable now
published).
\item \textbf{BBN under alternative matter conventions.} The F95
closure has been verified under Brown trace and partially under
Schutz/Hawking conventions; full coverage is identified.
\OPEN.
\end{itemize}

\subsection{The Derivation Passport as methodology}

The Derivation Passport (Appendix~\ref{app:passport}) is the
machinery that has made the audit-and-retraction discipline of this
paper possible. It is a deterministic gate-engine whose state
records every property ISST has stamped (e.g.\
\textsf{scalar\_propagates=false} from F128) and whose output is a
verdict on every external formula one might reach for. We invite
others to use it, criticise it, and extend it.

\subsection{Speculative directions (parking lot)}

The following are structurally consistent with the postulate but
not yet derived. They are listed here so the record is complete and
not in the results sections.

\begin{itemize}
\item \emph{GW energy partition between NS--NS and BH--BH mergers}
as a measure of information processing rate. Kill condition: if
$\Psi_{\rm rad}$ is not derivable from the action as a conserved
or source quantity.
\item \emph{Decoherence as gravitational information processing.}
Highly speculative; no calculation; energy scale likely
unfalsifiable.
\item \emph{Big-$G$ laboratory scatter as $f$-correlated.} If
$G=1/\Psi$ with environment-dependent $f$, the persistent
scatter in big-$G$ measurements at national metrology institutes
may correlate with local geoid/Bouguer anomalies.
\item \emph{Cyclic universe (Axiom~VI) as a topological closure.}
The horizon and monopole problems are addressed by topological
inheritance from the previous cycle, not by intra-cycle inflation.
\end{itemize}

These are flagged as speculative throughout this paper's
\texttt{predictions\_pending.md} and carry their own kill
conditions.

%==========================================================
\section{Conclusion}
\label{sec:conclusion}

ISST commits to a single action with one matter-side coupling and
no kinetic term for the scalar. From that action: the Branch
theorem partitions spacetime into a GR-recovering compact-object
sector and an inhomogeneous cosmological sector; the static-fluid
theorem (App.~\ref{app:static_fluid_theorem}) routes every physical
star into Branch~B, underwriting exact PPN ($\gamma=\beta=1$),
$c_{\rm GW}=c$ structurally rather than by tuning,
$|\dot G/G|$ five orders below LLR, and inherited SBBN; the
acceleration scale $a_{\rm crit}\simeq 1.07\times 10^{-10}$~m/s$^2$
falls out of two algebraic identities, sitting $11\%$ from MOND
$a_0$; the cosmological-floor identity
$(1{+}\fprim)\Omega_b=\Omega_m$ predicts $\Omega_m/\Omega_b$ within
$4.2\%$ of Planck from a single SM thermodynamic input ($1.6$--$1.8\sigma$
tension reported as honest residual). Rotation curves fit
competitively on the full SPARC sample with the central $Q{=}1$
comparison robust under held-out cross-validation (F147); the
lighthouse mechanism closes the bulk of the Bullet $\kappa$
contrast under empirical $f_s$ values; the $H_0$ tension structure
is reproduced as a frame artefact (literal recovery open).

ISST \emph{eliminates dark matter as a particle species}; the CMB
perturbation structure under the strict F01 action remains open
(Sec.~\ref{sec:theta}). The first-pass acoustic angle uses a
$(w_0,w_a)$ proxy carrying ${\sim}\,5$--$10\%$ systematic
uncertainty; the strict uniform-$f$ perturbation source is
falsified at the third peak ($>10\sigma$); the env-dep-$f$ closure
provides percent-level corrections, not the factor-$\sim\!30$
needed (Sec.~\ref{sec:env_dep_f_closure}). The distinguishing
predictions live in late-time structure formation: the lighthouse
mechanism in cluster mergers, the $f\sigma_8$ shape on the
Wiltshire background (DESI~DR2), the $+2.5\%$ void-lensing
$G$-ratio (Euclid), and the processing-history-dependent $f_s$ in
jellyfish galaxies. The full inventory of open items is in
Sec.~\ref{sec:future}; the dropped predictions in
Sec.~\ref{sec:dropped}; the timeline-and-precision falsifier table
in Sec.~\ref{sec:falsifier_table}.

\textbf{Self-containment and reproducibility.} Every load-bearing
derivation is presented in the body or appendices
(App.~\ref{app:passport}--\ref{app:adm_dof}). The companion site
\href{https://lily-labs.co.uk/isst}{lily-labs.co.uk/isst} provides
supplementary computational outputs and the interactive Derivation
Passport engine; CAMB \texttt{params.ini}, the modified CLASS
\texttt{source/perturbations.c} patch, and the SPARC analysis
scripts will be archived on Zenodo upon acceptance. The invitation
is open: here are our kill conditions; try to trigger them.

%==========================================================
\acknowledgments

We thank D.~L.~Wiltshire and the timescape community for the
two-domain framework on which ISST's cosmology rides; F.~Lelli and
the SPARC team for the rotation-curve database; the Planck and DESI
collaborations for the precision parameters and survey calibration
this work depends on. The interactive companion site and live
gate-engine were built with assistance from ``Lily,'' an additional
AI collaborator (separate session-level instance of Claude) used for
adversarial hypothesis generation.

%==========================================================
\appendix

\section{The Derivation Passport}
\label{app:passport}

The passport is a JSON state file that records every structural
property ISST has stamped, together with a property graph and an
implication engine. Every external framework (e.g.\
\textsf{chameleon\_screening}, \textsf{palatini\_ppn}) carries
preconditions; the gate evaluates the conjunction of preconditions
against the passport state and returns one of \textsc{[allowed]},
\textsc{[warning]}, \textsc{[conditional]}, \textsc{[denied]}.
Every verdict names the F-task or axiom that stamped each
contributing property, so that any verdict can be traced back
through the derivation chain.

As of this paper, the passport has 70 properties stamped (75 live
after implications), with 7 frameworks \textsc{[allowed]},
56 \textsc{[conditional]}, and 25 \textsc{[denied]}. The
\textsc{[allowed]} list comprises \textsf{buchert\_averaging},
\textsf{szekeres\_model}, \textsf{continuity\_equation},
\textsf{palatini\_ppn}, \textsf{horndeski\_gravity} (with corner
note), \textsf{adm\_decomposition}, and
\textsf{wiltshire\_two\_domain}. The \textsc{[denied]} list is
dominated by FRW-presuming or propagating-scalar frameworks
(\textsf{friedmann\_equations}, \textsf{chameleon\_screening},
\textsf{vainshtein\_screening}, \textsf{linear\_growth\_factor},
\textsf{halo\_model}, \textsf{convergence\_poisson},
\textsf{boltzmann\_hierarchy\_cmb}, etc.) plus
\textsf{damour\_esposito\_farese\_ppn} (requires propagating scalar)
and \textsf{birkhoff\_theorem\_gr} (requires Ricci-flat exterior
plus standard scalar coupling).

The hard rule is that any F-task using a \textsc{[denied]} framework
without first updating the passport is structurally invalid. Six
results in this paper's audit history triggered this rule and were
retracted.

\section{Critical-acceleration derivation}
\label{app:acrit_derivation}

This appendix gives the full algebra behind
Eqs.~\eqref{eq:wf_00}--\eqref{eq:acrit_main} of
Sec.~\ref{sec:acrit}. It assembles the linearised perturbation
analysis of F84 and the matter-era power-law analysis of F80 into a
single derivation chain.

\subsection{Linearised Branch~A: $C_{\rm grad} = 2/3$}
\label{app:Cgrad}

Take the conformal-Newtonian (longitudinal) gauge metric
\begin{equation}
ds^2 = -(1+2\Phi)\,dt^2 + (1-2\Phi_N)\,\delta_{ij}\,dx^i dx^j,
\end{equation}
with $\Psi = \Psi_0 + \delta\Psi$, $\Psi_0 \equiv 1/G_N$, and dust
stress-energy $T^\mu{}_\nu = \rho\,u^\mu u_\nu$ with $u^\mu = (1,\vec 0)$
and $T = -\rho$. Linearising the field equation \eqref{eq:eom} on
Minkowski + $\Psi_0$, retaining all terms at first order in
$(\Phi, \Phi_N, \delta\Psi/\Psi_0, \rho/\Psi_0)$ and reducing to
Laplacian form, the four independent equations are
\eqref{eq:wf_00}--\eqref{eq:wf_branchA} of the main text. From
\eqref{eq:wf_branchA} substitute $\nabla^2\Phi = 2\nabla^2\Phi_N$
into \eqref{eq:wf_trace}:
\begin{equation}
4\Psi_0\nabla^2\Phi_N - 4\Psi_0\nabla^2\Phi_N + 3\nabla^2\delta\Psi
= -8\pi(1+f)\rho \;\Longrightarrow\;
\nabla^2\delta\Psi = -\tfrac{8\pi}{3}(1+f)\rho.
\end{equation}
Inserting back into \eqref{eq:wf_00}:
\begin{equation}
2\Psi_0\nabla^2\Phi_N = 8\pi(1+f)\rho + \nabla^2\delta\Psi
= 8\pi(1+f)\rho - \tfrac{8\pi}{3}(1+f)\rho
= \tfrac{16\pi}{3}(1+f)\rho,
\end{equation}
giving $\nabla^2\Phi_N = (8\pi/3)(1+f)\rho/\Psi_0 = (2/3)\,4\pi
G_N(1+f)\rho$, and \eqref{eq:wf_branchA} then gives
$\nabla^2\Phi = (4/3)\,4\pi G_N(1+f)\rho$.

For a point mass $M$ at the origin, integrating
$\nabla^2\delta\Psi = -(8\pi/3)(1+f) M\,\delta^3(\vec r)$
gives
\begin{equation}
\delta\Psi(r) = \tfrac{2}{3}(1+f)\,G_N M / r,
\quad
\bigl|\nabla(\delta\Psi/\Psi_0)\bigr| = \tfrac{2}{3}(1+f)\,
\frac{G_N M}{r^2 c^2} = \tfrac{2}{3}(1+f)\,
\frac{a_{\rm bar}(r)}{c^2}.
\end{equation}
The Bardeen gauge-invariant quantities $\Phi$, $\Phi_N$ and the
spacetime-scalar $\delta\Psi$ are insensitive to the longitudinal
slicing, so this coefficient holds in any choice of slicing. The
candidate alternatives ($4/3$, $2$, $1$) violate one of
\eqref{eq:wf_offdiag} or \eqref{eq:wf_branchA} (cross-checked in
F84) and are ruled out structurally, not numerically.

\subsection{Matter-era power-law on the wall background:
$p/n = (\sqrt{21}-3)/2$}
\label{app:pn}

Take $a(t) \propto t^n$, $\Psi(t) \propto t^p$, $\rho(t)
\propto t^{-3n}$ on the leading-order wall Friedmann
\begin{equation}
\label{eq:wallA}
H^2 \;=\; \tfrac{16\pi}{9}\,(1{+}f_{\rm prim})\,\rho/\Psi,
\end{equation}
the Branch-A transport
\begin{equation}
\label{eq:transportA}
\Box\Psi = -\ddot\Psi - 3H\dot\Psi
       \;=\; -\tfrac{8\pi}{3}\,(1{+}f_{\rm prim})\,\rho,
\end{equation}
and continuity $\dot\rho + 3H\rho = 0$ (automatic for the power-law
ansatz). Computing each term:
\begin{align}
H &= n/t, \quad H^2 = n^2/t^2, \\
\ddot\Psi + 3H\dot\Psi &= [p(p-1) + 3np]\,\Psi_0\,t^{p-2}
                       = p(p-1+3n)\,\Psi_0\,t^{p-2}.
\end{align}
Power matching the Friedmann equation gives $-2 = -3n - p$, hence
$p = 2 - 3n$.

\textbf{Friedmann amplitude:}
$n^2 = (16\pi/9)(1{+}f_{\rm prim})\,\rho_0/\Psi_0$, equivalently
$(8\pi/3)(1{+}f_{\rm prim})\,\rho_0/\Psi_0 = (3/2)\,n^2$.

\textbf{Transport amplitude:}
$p(p-1+3n) = (8\pi/3)(1{+}f_{\rm prim})\,\rho_0/\Psi_0
= (3/2)\,n^2$.

With $p = 2 - 3n$ the LHS becomes
$(2-3n)\,(2-3n-1+3n) = (2-3n)\cdot 1 = 2-3n$.
Therefore
\begin{equation}
\label{eq:powerlaw_quadratic}
2 - 3n = \tfrac{3}{2}\,n^2 \;\Longleftrightarrow\;
3n^2 + 6n - 4 = 0,
\end{equation}
which is \eqref{eq:n_quadratic} of the main text. The roots are
$n = (-3 \pm \sqrt{21})/3$; the physical (expanding) root is
$n = (\sqrt{21}-3)/3 \approx 0.5275$, giving $p = 5 - \sqrt{21}
\approx 0.4174$ and
\begin{equation}
\frac{\dot\Psi}{H\Psi}\bigg|_0 = \frac{p}{n} =
\frac{3(5-\sqrt{21})}{\sqrt{21}-3} =
\frac{3(5-\sqrt{21})(\sqrt{21}+3)}{(\sqrt{21}-3)(\sqrt{21}+3)} =
\frac{3(2\sqrt{21}-6)}{12} = \frac{\sqrt{21}-3}{2},
\end{equation}
the $0.7913$ of \eqref{eq:p_over_n}.

This is a unique attractor of the coupled system: numerical Radau
integration with $\rm rtol = 10^{-11}$ from initial conditions
$10^{-2}/H_0$ in the deep matter era to $1/H_0$ recovers the
analytical asymptote to $|\Delta(p/n)| = 7.5\times 10^{-13}$, with
the Friedmann constraint residual at $2.1\times 10^{-12} H^2$ peak
(F80, \texttt{isst\_f80\_void\_scale.py}). The alternative full
Friedmann (cross-term retained) gives $p/n = \sqrt 3$; the
Brans--Dicke $\omega = 0$ form (no wall $(2/3)$ subtraction) gives
$p/n = 1$. Both alternatives predict $a_{\rm crit}$ values
$50$--$100\%$ off MOND $a_0$ at any $H_0 \in [60, 75]$~km/s/Mpc and
are ruled out empirically (F84 Table~5).

\subsection{Stability-gradient crossover}

The Branch-A constraint $R = 0$ holds pointwise in the cosmological
background; the field equation \eqref{eq:eom} with linear $\Psi$
admits no propagating mode (Sec.~\ref{sec:branch}). At galactic
scales, the local $|\nabla(\delta\Psi/\Psi_0)|$ scales as
$(2/3)(1{+}f_{\rm prim})\,a_{\rm bar}/c^2$ (Eq.~\ref{app:Cgrad});
the cosmological floor is set by the
matter-era $\dot\Psi/(H\Psi)$ on the wall background giving a
gradient scale $L_{\rm eff} = (p/n)\,c/H_0$ and a corresponding
floor $H_0^2 L_{\rm eff}/c^2$. Equating:
\begin{equation}
\tfrac{2}{3}(1{+}f_{\rm prim})\,\frac{a_{\rm crit}}{c^2}
\;=\;
\frac{H_0^2}{c^2} \cdot \frac{p}{n}\,\frac{c}{H_0}
\;=\; \frac{p}{n}\,\frac{H_0}{c},
\end{equation}
which rearranges to
\begin{equation}
a_{\rm crit} \;=\; \tfrac{3}{2}\,\frac{p}{n}\,
\frac{cH_0}{1+f_{\rm prim}}
\;=\; \frac{3(\sqrt{21}-3)}{4}\,\frac{cH_0}{1+f_{\rm prim}},
\end{equation}
the result \eqref{eq:acrit_main}. Numerically, with $cH_0 = 6.003
\times 10^{-10}$~m/s$^2$ at $H_0 = 61.79$~km/s/Mpc and
$(1{+}f_{\rm prim}) = 6.664$:
\begin{equation}
a_{\rm crit}^{\rm ISST} \;=\; \frac{1.1869 \times 6.003}{6.664}\times
10^{-10}\,{\rm m/s^2} \;=\; 1.069\times 10^{-10}\,{\rm m/s^2}.
\end{equation}
At Planck $H_0 = 67.4$ the value is $1.166\times 10^{-10}$ ($-2.8\%$
vs MOND); at SH0ES $H_0 = 73.0$ it is $1.263\times 10^{-10}$ ($+5.3\%$).
The $11\%$ offset at the committed ISST $H_0$ is the falsifiable
differential prediction.

\section{Wiltshire-forcing proof}
\label{app:wiltshire_forced}

This appendix expands the proofs of Claims~1 and~2 from
Sec.~\ref{sec:wiltshire_forced} and gives the explicit pointwise
$R = 0$ wall-metric construction.

\subsection{Claim 1 in detail}
\label{sec:claim1}

The flat-FRW Christoffels are $\Gamma^0_{ij} = a\dot a\,\delta_{ij}$
and $\Gamma^i_{0j} = (\dot a/a)\,\delta^i_j$, giving the non-zero
Ricci components
\begin{equation}
R_{00} = -3\,\frac{\ddot a}{a},
\qquad
R_{ij} = (a\ddot a + 2\dot a^2)\,\delta_{ij}.
\end{equation}
The Ricci scalar is therefore
\begin{equation}
R = g^{\mu\nu}R_{\mu\nu}
= 3\,\frac{\ddot a}{a} + \frac{3}{a^2}(a\ddot a + 2\dot a^2)
= 6\bigl(\ddot a/a + H^2\bigr) = 6(\dot H + 2H^2).
\end{equation}
Imposing $R = 0$ gives the autonomous ODE
$\dot H = -2H^2$. Integrating, $-1/H = -2t + C$; choosing the
origin of time $t_0 = 0$ gives $H(t) = 1/(2t)$ and
$a(t) = a_0(t/t_0)^{1/2}$. \emph{No matter content was used.} The
constraint is purely geometric.

\subsection{Claim 2 in detail}
\label{sec:claim2}

The trace of \eqref{eq:eom} with $g^{\mu\nu}G_{\mu\nu} = -R$ and
$g^{\mu\nu}g_{\mu\nu} = 4$ is
\begin{equation}
-\Psi R + 3\Box\Psi = 8\pi(1+f)T.
\end{equation}
Branch~A ($R = 0$) reduces this to the transport equation
$\Box\Psi = (8\pi/3)(1+f)T$. The $(t,t)$ mixed-index field equation
\begin{equation}
\Psi G^t{}_t + \Box\Psi - \nabla^t\nabla_t\Psi = 8\pi(1+f)T^t{}_t,
\end{equation}
on flat FRW with $\Psi = \Psi(t)$ and dust $T^t{}_t = -\rho$, with
$G^t{}_t = -3H^2$, $\nabla^t\nabla_t\Psi = -\ddot\Psi$, and
$\Box\Psi = -\ddot\Psi - 3H\dot\Psi$, becomes
\begin{equation}
\label{eq:tt_FRW_full}
3\Psi H^2 + 3H\dot\Psi = 8\pi(1+f)\rho.
\end{equation}
Substituting Claim~1's $H = 1/(2t)$ and dust scaling
$\rho(t) = \rho_0(t_0/t)^{3/2}$ into \eqref{eq:tt_FRW}:
\begin{equation}
\label{eq:tt_FRW_dust}
\frac{3\Psi(t)}{4t^2} + \frac{3\dot\Psi(t)}{2t}
\;=\; 8\pi(1+f)\rho_0\,t_0^{3/2}\,t^{-3/2}.
\end{equation}
Multiplying by $2t/3$ rearranges this as the first-order linear
ODE
\begin{equation}
\label{eq:psi_ode_app}
\dot\Psi + \frac{\Psi}{2t} \;=\; \frac{16\pi}{3}(1{+}f)\rho_0\,t_0^{3/2}\,t^{-1/2}.
\end{equation}
The integrating factor is $\mu(t) = \exp\!\int dt/(2t) = t^{1/2}$.
Multiplying both sides by $\mu$ gives
$d/dt\bigl(t^{1/2}\Psi\bigr) = (16\pi/3)(1{+}f)\rho_0 t_0^{3/2}$,
which integrates to
\begin{equation}
\label{eq:psi_general_app}
\Psi(t) \;=\; A\,t^{1/2} + K\,t^{-1/2},
\qquad
A \equiv \frac{16\pi}{3}(1{+}f)\rho_0\,t_0^{3/2},
\end{equation}
with $K$ a free integration constant. We verify both modes by
direct substitution: for $\Psi = A t^{1/2}$,
$\dot\Psi = (A/2)t^{-1/2}$, the LHS of \eqref{eq:tt_FRW_dust} is
$3At^{-3/2}/4 + 3At^{-3/2}/4 = (3A/2)t^{-3/2}$, matching the RHS
on $A = (16\pi/3)(1{+}f)\rho_0 t_0^{3/2}$. For $\Psi = Kt^{-1/2}$,
$\dot\Psi = -(K/2)t^{-3/2}$, the LHS is
$3Kt^{-5/2}/4 - 3Kt^{-5/2}/4 = 0$, satisfying the homogeneous
equation. Both modes are therefore valid; the system is
algebraically consistent. The trace transport
$\Box\Psi = -(8\pi/3)(1{+}f)\rho$ is satisfied identically by both
modes under the same $A$ matching.

\emph{The incompatibility is observational.} Claim~1 forces
$a(t) \propto t^{1/2}$, which is the radiation-era expansion law
($q = 1$). Matter-era observations through SN-Ia, BAO, and CMB
acoustic-peak positions require an effective expansion law
significantly closer to $a(t) \propto t^{2/3}$ ($q = 1/2$) at
intermediate redshift, regardless of the value of $K$ in
\eqref{eq:psi_general_app} (the integration constant rescales
$\Psi$ but does not change the functional form of $a(t)$, which is
fixed by the geometric condition $R = 0$ in Claim~1). Branch~A on
flat FRW therefore predicts the wrong cosmological deceleration
parameter at the matter era. The mismatch is $\mathcal O(1)$ in
$q$, not a subleading correction.

For radiation matter ($T = -\rho + 3P$ with $P = \rho/3$), the
trace vanishes identically: $T = 0$. The transport equation reduces
to $\Box\Psi = 0$, which admits the static solution
$\Psi = \Psi_0 = \mathrm{const}$. Substituting $\dot\Psi = 0$,
$\ddot\Psi = 0$ into \eqref{eq:tt_FRW_dust} (with $\rho_{\rm rad}
\propto a^{-4} \propto t^{-2}$ on $a \propto t^{1/2}$) gives
$3\Psi_0/(4t^2) = 8\pi(1+f)\rho_{\rm rad,0}\,t^{-2}$, both sides
$\propto t^{-2}$. Consistent. Branch~A flat FRW is the correct
radiation-era cosmology; the matter-era expansion law it predicts is
incompatible with observation.

\subsection{Pointwise $R = 0$ wall metric}
\label{sec:wall_pointwise}

The minimum geometric structure compatible with Branch~A + dust uses
the planar anisotropic ansatz~(F12)
\begin{equation}
\label{eq:wall_metric}
ds^2 = -dt^2 + a(t)^2(dx^2+dy^2) + b(t,z)^2\,dz^2,
\qquad
\Psi(t,z) = \Psi_0(t) + \psi(z).
\end{equation}
Computing the Ricci scalar of \eqref{eq:wall_metric} symbolically:
\begin{equation}
R = R_{\rm temporal}(t) + R_{\rm spatial}(z) +
   R_{\rm cross}(t,z;\partial_z b),
\end{equation}
where the cross terms involving $\partial_z b$ enter
$R_{tt}$ and $R_{zz}$ with opposite signs at the bound-wall fixed
point $K \equiv \dot b/b - \dot a/a = 0$. With $a(t)\propto t^{2/3}$
(matter-era wall expansion), $R_{\rm temporal} = 0$ identically. The
remaining condition is the spatial balance
\begin{equation}
\label{eq:psi_balance}
\partial_z^2\psi = -\tfrac{8\pi}{3}(1+f)\rho,
\end{equation}
which is the spatial form of the Branch-A transport equation. The
$z$-derivatives of $b(t,z)$ cancel exactly between $R_{tt}$ and
$R_{zz}$ at the bound-wall fixed point, so $R = 0$ pointwise reduces
to a 1-parameter ODE in $t$ at each $z$, with $K = 0$ as its stable
fixed point. Numerical verification on a $200\times 200$ grid
(F12, working/wall\_metric\_pointwise\_R0.md) confirms:
(i)~spatial balance \eqref{eq:psi_balance} satisfied pointwise to
$4\times 10^{-10}$;
(ii)~spatial average $\langle\Delta\psi\rangle_w =
-(8\pi/3)(1+f)\langle\rho\rangle_w$ recovering F12's
\eqref{eq:wall_friedmann_main} coefficient $16\pi/9$ at six-digit
agreement;
(iii)~pointwise $|R|$ decays from $6\times 10^{-5}$ (initial
finite-difference noise) to $3\times 10^{-9}$ as $K$ relaxes to zero.

This is the explicit constructive proof that the wall-side of the
two-domain split admits Branch~A + dust pointwise. Combined with
the Milne void (where $T \to 0$ trivially admits $R = 0$ and
$\Psi = \mathrm{const}$), the two-domain split is constructively
realised.

\subsection{Selection between two-domain alternatives}
\label{sec:wiltshire_selection}

Claim~2 motivates an inhomogeneous background (the
matter-era observational mismatch on flat FRW);
Sec.~\ref{sec:wall_pointwise}
constructs the planar wall metric. Whether the planar
wall-pointwise construction is \emph{unique} among all
inhomogeneous Branch~A + dust solutions reduces to the
classification problem: \emph{which dust spacetimes satisfy
$R = 0$ pointwise?} For the two metric classes most studied in
inhomogeneous cosmology:
\begin{itemize}
\item \textbf{Lema\^itre--Tolman--Bondi (LTB)} dust spacetimes
$ds^2 = -dt^2 + R'^2/(1+2E)\,dr^2 + R^2 d\Omega^2$ have Ricci
scalar that depends on $E(r)$ and the bang-time $t_B(r)$; setting
$R = 0$ pointwise requires solving a coupled PDE system for these
arbitrary functions. We have not constructed an explicit LTB
solution with $R = 0 + \text{dust}$, nor proved one cannot exist;
the classification is identified as \OPEN.
\item \textbf{Szekeres} spacetimes generalise LTB and admit broader
inhomogeneity. The $R = 0 + \text{dust}$ classification is similarly
unresolved.
\end{itemize}
Wiltshire's specific timescape geometry~\cite{Wiltshire2007,Wiltshire2009}
is the simplest two-domain Buchert-averaged backreaction model
that:
(i)~realises $\langle R\rangle = 0$ on the volume average;
(ii)~admits the pointwise $R = 0$ wall metric of
Sec.~\ref{sec:wall_pointwise} as its wall-interior limit;
(iii)~matches the SN-Ia luminosity--redshift relation at
$f_{v0} = 0.762$~\cite{Dam2017,Wiltshire2009}.
The selection over LTB/Szekeres is therefore on the basis of
minimality (criterion i, two domains rather than continuous
$E(r)$/$t_B(r)$), constructive realisation (ii), and observational
match (iii). \emph{Strict exclusion of all LTB/Szekeres alternatives
is identified as a separate classification problem,
\OPEN.}

\subsection{What this proves}

Within the closed-form analytic results above:
\begin{itemize}
\item Branch~A on flat FRW with dust is algebraically consistent
(general solution \eqref{eq:psi_general_app}) but observationally
incompatible with matter-era expansion: the geometric constraint
$R = 0$ forces $a\propto t^{1/2}$ (radiation-era scaling), whereas
matter-era observations require $a\propto t^{2/3}$ (\PROVEN\ for
the geometric obstruction; observational mismatch is at
$\mathcal O(1)$ in the deceleration parameter).
\item The two-domain split (Milne voids + planar-anisotropic walls)
\emph{constructively} realises Branch~A + dust + observed
matter-era cosmology (\PROVEN).
\item Wiltshire's timescape phenomenology arises as the
volume-average of this two-domain construction at $f_{v0} = 0.762$
(\DEMO; matches Wiltshire 2009 to $0.09\%$ on the dressing ratio
$\mathcal{R}$).
\item Strict exclusion of LTB/Szekeres alternative inhomogeneous
backgrounds is \emph{not} proved; that classification is
\OPEN.
\end{itemize}

\section{Wall Friedmann derivation: anisotropic slab to isotropic limit}
\label{app:wall_friedmann_proof}

This appendix gives the full derivation of
\eqref{eq:wall_friedmann_main} from F01 on the planar wall slab,
with explicit verification that energy-momentum is conserved at every
step of the anisotropic-to-isotropic transition, providing a
clear bridge between the anisotropic wall slab and the averaged
Friedmann equation that ensures no energy-momentum is lost in the
averaging process. The full derivation, including symbolic
verification with \textsc{sympy}, is provided in the supplementary
material accompanying this paper.

The pointwise existence of the wall metric is established in
Appendix~\ref{sec:wall_pointwise}; F141 (this appendix) proves the
$16\pi/9$ coefficient \emph{from the field equations} on that metric
and verifies energy-momentum conservation in the isotropisation
limit.

\subsection{The planar wall ansatz and Hubble rates}

We work on the planar-symmetric metric
\begin{equation}
\label{eq:wall_metric_F141}
ds^2 = -dt^2 + a(t)^2(dx^2 + dy^2) + b(t,z)^2\,dz^2,
\qquad \Psi = \Psi_0(t) + \psi(t,z),
\end{equation}
already used in Appendix~\ref{sec:wall_pointwise}. Define
$H_\perp \equiv \dot a/a$ ($z$-independent perpendicular Hubble),
$K(t,z) \equiv \partial_t b/b$ ($z$-dependent parallel expansion),
and $\mu(t,z) \equiv \partial_z b/b$ (anisotropy gradient). The
bound-wall fixed point has $K = 0$ (frozen thickness) and
$b \to 1$.

\subsection{Branch-A constraint and dynamical isotropisation}

Direct Christoffel computation on \eqref{eq:wall_metric_F141}
(verified in \texttt{F141\_em\_conservation\_check.py}) gives
\begin{equation}
\label{eq:R_aniso_app}
R = -4\dot H_\perp - 6 H_\perp^2 - 2\partial_t K - 2 K^2 - 6 H_\perp K.
\end{equation}
The $\mu$-dependent terms (spatial derivatives of $b$) cancel
identically between $R_{tt}$ and $R_{zz}$, so the Branch-A constraint
$R = 0$ reduces to a one-parameter ODE in $t$ at each $z$:
\begin{equation}
\label{eq:branch_a_app}
2\dot H_\perp + \partial_t K + 3 H_\perp^2 + K^2 + 2 H_\perp K = 0.
\end{equation}
This admits two fixed points in $K$: the bound-wall $K = 0$ branch,
which gives $a \propto t^{2/3}$ (standard EdS matter-era), and the
unstable $K = -2 H_\perp$ branch corresponding to collapsing thickness.
Linear stability analysis around $K = 0$ gives a decay rate
$\sim H_\perp$ to the bound branch; physical bound walls relax to
this attractor, and the wall therefore isotropises dynamically. This
is \emph{not} a Wiltshire import: the isotropisation is a stable
consequence of \eqref{eq:branch_a_app}.

\subsection{Trace transport and the spatial $\Psi$ profile}

The trace of \eqref{eq:eom} on Branch~A gives
\begin{equation}
\label{eq:trace_app}
\Box\Psi \;=\; -\tfrac{8\pi}{3}(1{+}f)\rho \quad
(\text{dust, } T = -\rho).
\end{equation}
On \eqref{eq:wall_metric_F141},
\begin{equation}
\Box\Psi = -\ddot\Psi_0 - \ddot\psi - (2H_\perp + K)(\dot\Psi_0 + \dot\psi)
+ (1/b^2)(\partial_z^2\psi - \mu\,\partial_z\psi).
\end{equation}
The fast (spatially-varying) part of \eqref{eq:trace_app}, in the
bound-wall limit $b \to 1$, $\mu \to 0$, gives the spatial Poisson
equation
\begin{equation}
\label{eq:spatial_poisson_app}
\partial_z^2\psi \;\approx\; -\tfrac{8\pi}{3}(1{+}f)\rho_w
\quad\text{(inside the wall, }\rho \approx \rho_w\text{ uniform)},
\end{equation}
which derives the spatial-balance assumption made in F12 rather than
imposing it.

\subsection{The isotropisation limit and the $16\pi/9$ coefficient}

Substituting the spatially-balanced $\psi$ profile into the
$(tt)$ component of \eqref{eq:eom} on the isotropic wall
($K \to 0$, $b \to 1$, $H_\perp \to H_w$, $a \to a_w$):
\begin{align}
\Psi_0(-3 H_w^2) - 3 H_w \dot\Psi_0 - \tfrac{8\pi}{3}(1{+}f)\rho
\;&=\; -8\pi (1{+}f)\rho, \\
\therefore\quad 3\Psi_0 H_w^2 + 3 H_w \dot\Psi_0
\;&=\; \tfrac{16\pi}{3}(1{+}f)\rho.
\label{eq:wall_pre_div_app}
\end{align}
Dividing by $3\Psi_0$ and dropping the $\dot\Psi_0/\Psi_0$ term
(bounded at $\sim 10^{-7} H_0$ per LLR; F121),
\begin{equation}
\label{eq:wall_friedmann_derived_app}
\boxed{H_w^2 \;=\; \frac{16\pi}{9}\,\frac{(1+f)\rho_b}{\Psi_0}}
\quad \equiv \eqref{eq:wall_friedmann_main}.
\end{equation}

\textbf{Bookkeeping of the coefficient.}
The $16\pi/9$ traces to the partition
\begin{equation}
\label{eq:coeff_book_app}
\underbrace{8\pi(1{+}f)\rho}_{\text{full }(tt)\text{ source}}
\;-\;
\underbrace{\tfrac{8\pi}{3}(1{+}f)\rho}_{\text{absorbed by }\partial_z^2\psi
\text{ via }\eqref{eq:spatial_poisson_app}}
\;=\;
\underbrace{\tfrac{16\pi}{3}(1{+}f)\rho}_{\text{residual on (tt) RHS}},
\end{equation}
divided by the $3\Psi_0$ prefactor of $G^t{}_t$ in the FRW limit.
The spatial $\Psi$ profile $\psi(z)$ absorbs $1/3$ of the gravitational
source pointwise; the residual $2/3$ drives $H_w^2$ expansion.

\subsection{Energy-momentum conservation: explicit no-leak check}

The Noether identity for the diffeomorphism-invariant matter action
$S_m = \int\!\sqrt{-g}\,(1{+}f)\Lm$ gives
\begin{equation}
\label{eq:noether_em_app}
\nabla^\mu \Teff_{\mu\nu} = 0 \quad\text{identically}
\end{equation}
on any matter configuration. For the uniform-$f$ scope of the wall
Friedmann derivation ($\Delta^f_{\mu\nu} = 0$, Sec.~\ref{sec:action}),
\eqref{eq:noether_em_app} reduces to $\nabla^\mu[(1{+}f) T_{\mu\nu}] = 0$,
which on \eqref{eq:wall_metric_F141} for dust gives the continuity
equation
\begin{equation}
\label{eq:continuity_aniso_app}
\dot\rho + (2 H_\perp + K)\rho = 0.
\end{equation}
In the isotropic limit $K \to H_\perp \to H_w$, this deforms continuously
to
\begin{equation}
\label{eq:continuity_iso_app}
\dot\rho + 3 H_w \rho = 0 \quad\Longrightarrow\quad \rho \propto a_w^{-3}.
\end{equation}
The dilution-rate coefficient $(2H_\perp + K) \to 3 H_w$ smoothly as
the wall isotropises; \emph{no terms are dropped or absorbed in the
limit}, no leak terms appear, and the coefficient of $\dot\rho$ remains
exactly~1 throughout. The Noether identity \eqref{eq:noether_em_app}
holds at every $K$ in the transient. The full symbolic verification
in \texttt{F141\_em\_conservation\_check.py} confirms
\eqref{eq:continuity_aniso_app} as the divergence of $(1{+}f)T_{\mu\nu}$
on \eqref{eq:wall_metric_F141} and \eqref{eq:continuity_iso_app} as
its $K \to H_\perp$ limit.

\subsection{Buchert cross-check: the $2/3$ factor's origin}

Standard Buchert averaging~\cite{Buchert2000} for dust on a generic
spatially-foliated spacetime gives
\begin{equation}
\label{eq:buchert_std_app}
3\,(\dot a_D/a_D)^2 = 8\pi G\langle\rho\rangle_D
- \tfrac12\langle R\rangle_D - \tfrac12 Q_D,
\end{equation}
with $Q_D$ the kinematic backreaction and $R$ the spatial Ricci. For
ISST Branch~A with bound-wall fixed point, $\langle R\rangle_D = 0$
(Branch~A pointwise) and $Q_D = 0$ (uniform expansion at the fixed
point), so \eqref{eq:buchert_std_app} reduces to $3 H_w^2 = 8\pi
G\langle\rho\rangle_w$, giving the standard $8\pi/3 = 24\pi/9$
coefficient.

ISST's $16\pi/9 = (8\pi/3) \times (2/3)$ differs by exactly the $2/3$
factor identified in \eqref{eq:coeff_book_app}. The standard Buchert
framework \emph{does not include} the $\Psi$-Lagrange-multiplier
sector---it assumes Einstein gravity $G_{\mu\nu} = 8\pi T_{\mu\nu}$.
ISST's modification is the additional $\Psi$ transport equation
\eqref{eq:trace_app}, which absorbs $1/3$ of the gravitational source
into the spatial profile $\psi(z)$. The effective Buchert coefficient
becomes $8\pi/3 \times 2/3 = 16\pi/9$.

This is not a coincidence with the Wiltshire two-domain
parameterisation. Wiltshire's wall Friedmann uses standard GR and
identifies the dust density as the cosmological matter density
$\rho_m$. Identifying $(1{+}f)\rho_b \equiv \rho_m$ via A04, ISST's
$H_w^{\rm ISST} = \sqrt{2/3}\, H_w^{\rm WI}$. The $2/3$ rescaling
propagates through the Wiltshire dressing $\mathcal{R}(f_{v0})
= 1.2825$ to give the $H_0$ predictions of Sec.~\ref{sec:hubble}
self-consistently; the coefficient invariance result of F12.INV
confirms $\mathcal{R}$ depends only on $f_{v0}$, not on the wall
coefficient.

\subsection{Coordinate-independence}

The $16\pi/9$ coefficient relates two observer-frame scalars ($H_w^2$
and $\rho_b$). The derivation in §C.4 used the synchronous gauge
\eqref{eq:wall_metric_F141}. Conformal-gauge ($ds^2 = a^2(\eta)(-d\eta^2
+ d\vec x^2)$) and harmonic-gauge ($\Box x^\mu = 0$) re-derivations
give the same coefficient \emph{by construction}: the F01 field
equation \eqref{eq:eom} is generally covariant, so its $(tt)$-projection
on the wall solution yields a gauge-invariant scalar relation between
$H_w^2$ and $\rho_b/\Psi_0$. The bookkeeping
\eqref{eq:coeff_book_app} is gauge-independent because each step uses
only the trace of \eqref{eq:eom} and the spatial Poisson
\eqref{eq:spatial_poisson_app}, both of which are scalar equations.

\subsection{Status}

The $16\pi/9$ coefficient in \eqref{eq:wall_friedmann_main} is upgraded
from \DEMO\ to \PROVEN\ by this appendix:
\begin{itemize}
\item The pointwise R~=~0 wall metric admits the bound-wall fixed
point as a stable attractor of \eqref{eq:branch_a_app} (no Wiltshire
isotropisation import).
\item The trace transport \eqref{eq:trace_app} on Branch~A gives the
spatial $\Psi$ profile \eqref{eq:spatial_poisson_app} that absorbs
$1/3$ of the source.
\item The $(tt)$ projection in the isotropic limit gives
\eqref{eq:wall_friedmann_derived_app} with coefficient $16\pi/9$.
\item Noether identity \eqref{eq:noether_em_app} guarantees
energy-momentum conservation at every $K$; the continuity equation
\eqref{eq:continuity_aniso_app}--\eqref{eq:continuity_iso_app} deforms
continuously without leak terms.
\item The $2/3$ factor relative to standard Buchert is the explicit
fraction of the source absorbed by $\Psi$ transport, traced through
\eqref{eq:coeff_book_app}.
\item Coordinate-independence holds by general covariance of
\eqref{eq:eom}.
\end{itemize}

\section{$H_0$ frame artefact: bare-to-dressed pipeline}
\label{app:h0_frame}

This appendix expands Sec.~\ref{sec:hubble} with the full algebraic
chain from the bare two-domain background to the values a wall
observer fitting a $\Lambda$CDM template recovers under the SH0ES
local-distance-ladder and Planck CMB-template anchors.

\subsection{The three frames and their relations}

On the matter-dominated two-domain Buchert
average~\cite{Wiltshire2007,Wiltshire2009}, the bare scale factor
$\bar a(t)$ and the wall scale factor $a_w(\tau)$ are related
through the lapse $\gamma_w(\tau) \equiv dt/d\tau$ by the volume
constraint
\begin{equation}
\bar a^3 = f_w\,a_w^3 + f_v\,a_v^3,
\end{equation}
where $f_w$, $f_v = 1-f_w$ are the wall and void volume
fractions. The void scale factor is Milne, $a_v(\tau) \propto \tau$.
Differentiating and re-expressing in the wall proper time, the
present-epoch tracker solution gives the relation
\begin{equation}
\label{eq:R_dressing}
\mathcal{R}(f_v) \;\equiv\; \frac{H_0^{\rm dressed}}{H_0^{\rm bare}}
\;=\; \frac{4f_v^2 + f_v + 4}{2(2 + f_v)},
\end{equation}
derived as Eq.~(25) of Wiltshire~\cite{WiltshirePRL2007} (and
verified independently in F26). The lapse itself satisfies
\begin{equation}
\gamma_0 \;=\; \frac{1}{2}\Bigl[1 + \tfrac{1}{2}f_{v0} +
   \sqrt{1 + f_{v0} + \tfrac{9}{4}f_{v0}^2}\Bigr],
\end{equation}
reading $\gamma_0 = 1.348$ at $f_{v0} = 0.695$ (Duley et al.\ best
fit) and $\gamma_0 = 1.380$ at $f_{v0} = 0.762$ (ISST commit). The
density dressing relation is
\begin{equation}
\Omega_m^{\rm dressed} \;=\; \gamma_0^3\,\Omega_m^{\rm bare},
\end{equation}
\emph{distinct} from $\mathcal{R}$ and not equal to
$\mathcal{R}^3$. Conflating $\gamma_0$ with $\mathcal{R}$ has
caused order-unity errors in earlier ISST drafts and is flagged
explicitly here.

\subsection{ISST-native bare $H_0$}
\label{sec:bare_h0}

A model-independent input from the CMB acoustic angle and sound
horizon gives the comoving angular-diameter distance to the
last-scattering surface
\begin{equation}
D_A^{\rm com}(z_*) = r_s/\theta_* = 14130.8~\text{Mpc},
\end{equation}
using $r_s = 147.09$~Mpc and $\theta_* = 1.04092\times 10^{-2}$
(Planck~\cite{Planck2018} primaries). This is independent of
$\Lambda$CDM. Integrating the Wiltshire two-domain tracker
forward from $z_* = 1089$ to $z = 0$ with the ISST-committed
parameters $f_{v0} = 0.762$ and $\Omega_b h^2 = 0.02237$, requiring
the integrated comoving distance to match $D_A^{\rm com}(z_*)$,
fixes
\begin{equation}
\label{eq:bare_H0_value}
H_0^{\rm ISST,bare} = 57.56~{\rm km\,s^{-1}\,Mpc^{-1}}.
\end{equation}
This is $\sim 15\%$ above the Wiltshire-implied bare value
($\sim 50.1$~km/s/Mpc; Duley et al.~\cite{Dam2017}), reflecting
ISST's larger $f_{v0}$ commitment and the no-$\Lambda$ closure. No
phenomenological fit was performed; the Planck primaries plus the
ISST tracker plus $f_{v0}$ uniquely determine
\eqref{eq:bare_H0_value} (script
\texttt{working/isst\_bare\_h0\_from\_cmb.py}).

\subsection{Uniform-$f$ dressed value at $f_{v0} = 0.762$}

Applying \eqref{eq:R_dressing} at $f_{v0} = 0.762$:
\begin{equation}
\mathcal{R}(0.762) = \frac{4\,(0.762)^2 + 0.762 + 4}{2\,(2.762)}
= \frac{2.323 + 0.762 + 4}{5.524} = 1.2825.
\end{equation}
The corresponding dressed value is
\begin{equation}
\label{eq:H0_uniform}
H_0^{\rm dressed,uniform-f} = 1.2825\times 57.56 = 73.82~{\rm km\,s^{-1}\,Mpc^{-1}}.
\end{equation}
This number is the diagnostic uniform-$f$ tracker output, not the
ISST commit. It evaluates close to but not equal to the SH0ES
central value $73.04$~km/s/Mpc. The ISST commit replaces this with
the $K_{\rm env}$-corrected closure described next.

\subsection{$K_{\rm env}$-corrected dressed value: $61.79$}

The uniform-$f$ tracker neglects the environment-dependent
re-weighting that ISST's $\chi_w(z)$ sigmoid (Sec.~\ref{sec:tracker})
introduces. Defining the environmental rescaling
\begin{equation}
\label{eq:Kenv_def}
K_{\rm env} = \bigl(H_{\rm uniform}/H_{\rm phen}\bigr)^2,
\end{equation}
the F61 closure on the Sheth--van~de~Weygaert two-barrier
abundance gives $K_{\rm env} = 1.484$. The corrected dressed value
is~\cite{Dam2017,NazerWiltshire2015}
\begin{equation}
H_0^{\rm ISST,dressed} = H_0^{\rm dressed,uniform-f}/\sqrt{K_{\rm env}}
= 73.82/\sqrt{1.484} = 60.61,
\end{equation}
and a small additional retemplating from the wall
$\dot\Psi/(H\Psi)$ correction shifts this to the committed
$61.79$~km/s/Mpc reported in \eqref{eq:H0_committed}. The full
T3k closure is reproduced numerically (F26a, F61) and matches
phenomenological best fit~\cite{NazerWiltshire2015} to within $0.02$~km/s/Mpc.

\subsection{Wall-observer template recoveries}
\label{sec:F85}

A wall observer fitting ISST's two-domain $d_L(z)$ with a
$\Lambda$CDM template recovers different $H_0$ depending on the
anchoring. The F85 forward pipeline computes both:

\textbf{Planck-template recovery.} Anchor $(\Omega_m, \Omega_b h^2,
\theta_*)$ to Planck values; minimise $\chi^2$ on a synthetic
ISST $C_\ell$ via the effective $(w_0, w_a)$ proxy of
Sec.~\ref{sec:theta}. The recovered $H_0$ is
\begin{equation}
\label{eq:H0_planck_template}
H_0^{\rm Planck-template} = 62.8~{\rm km\,s^{-1}\,Mpc^{-1}}.
\end{equation}
Observed Planck value: $67.4$. Pipeline residual $-7\%$.

\textbf{SH0ES-ladder recovery.} Anchor $M_B$ on a Cepheid+SN
distance ladder calibrated by ISST $d_L(z)$ in
$z\in[0.023, 0.15]$, then fit
$\Lambda$CDM($\Omega_m=0.315$) to the ladder-output $d_L$. The
recovered $H_0$ is
\begin{equation}
\label{eq:H0_sh0es}
H_0^{\rm SH0ES-ladder} = 70.6~{\rm km\,s^{-1}\,Mpc^{-1}}.
\end{equation}
Observed SH0ES value: $73.04$. Pipeline residual $-3.4\%$.

\textbf{Pipeline-internal spread.} The ISST pipeline produces
\eqref{eq:H0_planck_template} and \eqref{eq:H0_sh0es} from a
\emph{single underlying cosmology} (the same two-domain background
with the same $f_{v0} = 0.762$ commit). The spread
\begin{equation}
\Delta H_0^{\rm pipeline} = 70.6 - 62.8 = 7.8~{\rm km\,s^{-1}\,Mpc^{-1}}
\end{equation}
reproduces the observed Hubble tension structure
$\Delta H_0^{\rm obs} = 73.04 - 67.4 = 5.6$~km/s/Mpc with the
correct sign and within $\sim 40\%$ on magnitude. The literal
recovery of $73.0$ and $67.4$ is \emph{not} achieved: the
SH0ES-side pipeline output is $-3.4\%$ low, the Planck-side
$-7\%$ low. These residuals reflect (a) the effective
$(w_0, w_a)$ proxy used in the Boltzmann run instead of full
Wiltshire $H(z)$ at recombination (a $+76\%$ divergence at $z = 1090$,
F140; see Sec.~\ref{sec:theta}); and (b) the absence of
photon-path averaging in the SH0ES ladder pipeline (an
ISST-specific correction not yet built; F85 §6).

\textbf{Conclusion.} The Hubble tension is reproduced as a
pipeline artefact between Planck-template and SH0ES-ladder
recoveries of the same ISST cosmology, not as a tension between
data sets. Literal closure of both observed values within $1\%$
requires the photon-path-corrected SH0ES pipeline and the full
Wiltshire $C_\ell$, both \OPEN; the structure is reproduced
already.

\section{Audit summary and result inventory}
\label{app:audit}

The retroactive audit (2026-04-25) classifies all 88 audited results:

\begin{table}[h]
\caption{Audit classification.}
\label{tab:audit}
\begin{ruledtabular}
\begin{tabular}{lcl}
Class & Count & Action\\
\hline
GREEN & 60 & no action; load-bearing\\
YELLOW & 19 & note context predicate in writeup\\
ORANGE & 0 & all 9 re-derived clean; promoted GREEN\\
RED & 6 & retracted or superseded\\
BLACK & 3 & retracted; halo-ontology imports\\
\end{tabular}
\end{ruledtabular}
\end{table}

The RED list (Sec.~\ref{sec:dropped}): F32--F34 (used FRW background;
superseded by F82/F89), F69 (GR Poisson source), F70 (Shore-Johnson
additivity fails), F76 (f(R) metric equivalence inapplicable),
F77 (chameleon screening requires V($\Psi$); F01 has none),
F112/F113/F114 (FRW contamination; superseded by F95).

The BLACK list (Sec.~\ref{sec:dropped}): F65, F66, F68; all halo
ontology imports.

The ORANGE$\to$GREEN promotions (F02, F09, F82, F87, F106, F110,
F112a, F113a, and the legacy \texttt{isst\_growth.py}) re-derived
each result through an \textsc{[allowed]} or ISST-native framework
in \texttt{working/rederivations/}.

\section{Suggested figures}
\label{app:figures}

We list, but do not include here, the figures we recommend for the
referee version:

\begin{enumerate}
\item Branch-A vs.\ Branch-B map: which scales live in which branch,
with the surface $\Sigma$ between them.
\item Rotation-curve grid: ISST v8 vs.\ MOND vs.\ Newton on a
representative SPARC galaxy each from dwarf, LSB, late-type spiral,
and early-type spiral classes.
\item BTFR scatter: ISST v8 prediction versus observed slope $1.0\pm 0.1$.
\item $f\sigma_8(z)$ comparison: ISST F89 band, $\Lambda$CDM band,
and the eight current RSD measurements with errorbars.
\item Bullet $\kappa$ schematic: $\Sigma_\star$, $\Sigma_{\rm gas}$,
$(1+f_s)$ enhancement at the main and sub-cluster peaks.
\item Parameter-vs-derived table: $\Lambda$CDM 6 free parameters
versus ISST $\fprim$ + initial conditions.
\item Passport compatibility matrix: ALLOWED / CONDITIONAL / DENIED
verdicts for the 12 frameworks the paper's results depend on.
\end{enumerate}

%==========================================================
\section{Static-fluid theorem: physical inadmissibility of Branch~A perfect fluids}
\label{app:static_fluid_theorem}

This appendix gives the explicit derivation of the static-fluid
theorem (informally, ``F08'') invoked throughout the body for PPN
$\gamma=\beta=1$ (Sec.~\ref{sec:solar}), $c_{\rm GW}=c$ via Branch~B
selection of compact binaries (Sec.~\ref{sec:early}), the LLR
$\dot G/G$ bound (Sec.~\ref{sec:solar}), the rotating-NS extension
(Sec.~\ref{sec:solar}), and the DF2 / TDG predictions
(Sec.~\ref{sec:df2_tdg}). The result: any static spherically
symmetric isotropic perfect fluid on Branch~A is forced to have
$P_c=0$ and $P''(0)<0$, hence $P(r)<0$ at small radius,
\emph{regardless of the equation of state}. No physical fluid with
$P\geq 0$ and $dP/d\rho>0$ admits this profile; physical static
sources are forced into Branch~B, where the field equation reduces
to exact GR with $G_{\rm eff}=(1+f)/\Psi_0$.

\subsection{Setup}

Work in natural units $8\pi G=1$. Take a static spherically
symmetric metric
\begin{equation}
ds^2 = -e^{2\Phi(r)}dt^2 + A(r)^{-1}dr^2 + r^2 d\Omega^2,
\qquad A = e^{-2\Lambda},
\end{equation}
with perfect-fluid matter of density $\rho_0$ at the centre,
pressure $P(r)$, and constant interior $f=f_\star\geq 0$. Regularity
at $r=0$ requires $A(0)=1$, $\Phi'(0)=0$, $\Psi'(0)=0$, $f'(0)=0$.
Expand
\begin{align}
A &= 1 + a_2 r^2 + a_4 r^4 + O(r^6), \\
\Phi &= \Phi_c + \tfrac{1}{2}\Phi_2 r^2 + \tfrac{1}{24}\Phi_4 r^4 + O(r^6), \\
\Psi &= \Psi_c + \tfrac{1}{2}\psi_2 r^2 + \tfrac{1}{24}\psi_4 r^4 + O(r^6), \\
P &= P_c + \tfrac{1}{2}P_2 r^2 + O(r^4).
\end{align}

\subsection{Leading-order coefficients}

Imposing the Branch-A constraint $R=0$, the trace transport
$\Box\Psi=\tfrac{1}{3}(1+f_\star)T$ with $T=-\rho_0+3P$, the
$(tt)$ field equation, and the origin-consistency of $(rr)$ and
$(\theta\theta)$ yields
\begin{align}
\Phi_2 &= -a_2 = \tfrac{2(1+f_\star)\rho_0}{9\,\Psi_c}
&&\text{(from }R=0\text{ and }(tt)\text{)}, \\
\psi_2 &= -\tfrac{(1+f_\star)\rho_0}{9}
&&\text{(transport)}, \\
P_c &= 0
&&\text{(origin consistency)}.
\end{align}
The vanishing of $P_c$ is not assumed; it is forced uniquely
by the requirement that the algebraic pressure relation
\begin{equation}
\label{eq:Palg_app}
P(r) = \Psi\,\frac{A-1}{r^2} + \frac{2A}{r}\bigl(\Psi\,\Phi'+\Psi'\bigr)
        + A\,\Phi'\,\Psi'
\end{equation}
gives a finite consistent value at $r=0$. The transport equation
combined with the $(rr)$ field equation collapses to
\eqref{eq:Palg_app} with no $\Phi''$ or $\Psi''$ appearing; the
$(\theta\theta)$ equation, after use of the same transport and the
constraint $R=0$, reduces to the same relation. Pressure is therefore
algebraically determined by the metric state.

\emph{Explicit derivation of $P_c=0$.} Substituting the
regular-centre expansions into \eqref{eq:Palg_app} and collecting
$O(1)$ terms,
\begin{equation}
P_c = \Psi_c\,a_2 + 2\,\Psi_c\,\Phi_2 + 2\,\psi_2.
\end{equation}
With $\Phi_2 = -a_2$ ($R=0$), $a_2=-2(1+f_\star)\rho_0/(9\Psi_c)$
(($tt$)), and $\psi_2=-(1+f_\star)\rho_0/9$ (transport),
\begin{equation}
P_c = \Psi_c a_2 - 2\Psi_c a_2 + 2\psi_2
    = \tfrac{2(1+f_\star)\rho_0}{9} - \tfrac{2(1+f_\star)\rho_0}{9}
    = 0.
\end{equation}
The cancellation is algebraic: the $(tt)$ $a_2$ and the transport
$\psi_2$ enter with opposite signs and the $R=0$ identity
$\Phi_2=-a_2$ doubles the $a_2$ contribution to produce the
cancellation. The $(1+f_\star)$ factor enters every term uniformly
and does not affect the cancellation; the same algebra with
$f_\star=0$ (ordinary GR with Branch-A $R=0$) equally gives $P_c=0$.

\subsection{Next-order coefficients and physical inadmissibility}

At $O(r^2)$, self-consistency of $(rr)$ gives
\begin{equation}
\label{eq:P2_app}
P_2 = -\frac{2(1+f_\star)\rho_0^2}{9\,\Psi_c} < 0,
\end{equation}
and substitution confirms $(\theta\theta)$ yields the same value.
Hydrostatic equilibrium $P'=-(\rho_0+P)\Phi'$ at leading order
requires $P_2=\rho_0 a_2$, which with $a_2=-2(1+f_\star)\rho_0/(9\Psi_c)$
reproduces \eqref{eq:P2_app} exactly --- equilibrium is automatic by
the Bianchi identity.

\textbf{Theorem (physical inadmissibility).} For any static
spherically symmetric isotropic perfect-fluid source in Branch~A
with regular centre, constant interior $f=f_\star\geq 0$, and
$\rho_0=\rho(0)>0$:
\begin{equation}
P(0) = 0,
\qquad
P''(0) = P_2 = -\frac{2(1+f_\star)\rho_0^2}{9\,\Psi_c} < 0.
\end{equation}
The pressure therefore takes negative values at all sufficiently
small $r>0$. No matter obeying a thermodynamic equation of state
$P=P(\rho)$ with $P\geq 0$ and $dP/d\rho>0$ admits this profile.
Branch~A has no static spherically symmetric isotropic-fluid
solution describing ordinary self-gravitating matter, regardless
of the density profile $\rho(r)$ (extension to $\rho(r)$ with
$\rho_2$ varying gives the same $P_c$, $P_2$ by Bianchi-redundancy
in the higher-order coefficients).

\textbf{Corollary (the static-fluid theorem).} Every static
spherically symmetric source of ordinary matter (baryons obeying
a physical EOS) is in Branch~B, with $\Psi=\Psi_c$ constant
throughout the interior. Junction continuity at the stellar
surface with the vacuum exterior $\Psi=\Psi_0+D/r$ forces $D=0$:
$\Psi=\Psi_0$ identically in the exterior. The PPN parameters
satisfy $\gamma=\beta=1$ exactly by the Branch-B reduction
$\Psi_0 G_{\mu\nu}=8\pi(1+f)T_{\mu\nu}$, which is exact GR with
$G_{\rm eff}=(1+f)/\Psi_0$.

\subsection{Scope}

The theorem assumes (i) static spherical symmetry, (ii) regular
centre, (iii) isotropic perfect-fluid matter with a single scalar
pressure $P(\rho)$. Anisotropic stress (e.g.\ a strange-quark
crust with $P_r\neq P_t$) falls outside the origin-consistency
forcing of $P_c=0$ and is not covered here. Slowly rotating
configurations are reduced to the static case at $O(\Omega^0)$
(Sec.~\ref{sec:solar}); rapidly rotating equilibria invoke the
generalised stationary-equilibrium branch-selection rule
($df/d\tau=0$ in the matter rest frame), not this theorem
directly. Constant-$f_\star$ extension to radially varying $f(r)$
with $f'(0)=0$ is conjectured to follow from the same
Bianchi-redundancy structure; the extension is not load-bearing
for the body's PPN, $c_{\rm GW}$, LLR, and rotating-NS results,
which use only the $f=f_\star$ statement of the theorem.

\subsection{Numerical confirmation}

A full-GR shooter (\texttt{paper/branch\_a\_shooter\_v3.py})
integrates the Branch-A system with state
$\{\Phi,\Phi',A,\Psi,\Psi'\}$ from $r_0=10^{-4}$ to $r_{\max}=3$
in natural units, with DOP853 at relative tolerance $10^{-11}$.
A density scan over $\rho_0\in\{0.1,0.3,1,3,10\}$, a $\Psi_c$
scan over $\{0.3,1,3,10\}$, and an $f_\star$ scan over
$\{0,1,5,5.4,20\}$ reproduce the parametric prediction
$P_2=-2(1+f_\star)\rho_0^2/(9\Psi_c)$ to $0.7\%$ at $f_\star=20$
and $0.2\%$ at $f_\star=5.4$. Consistency residuals
$|P_{rr}-P_{\theta\theta}|/|P_{rr}|\sim 10^{-12}$ and emergent
$|R|\sim 10^{-10}$ hold across the full scan. Numerical confirmation
is therefore complete at the level of the constant-$f_\star$ ODE
system; the analytical theorem above is the load-bearing statement.

%==========================================================
\section{ADM constraint analysis: degree-of-freedom count}
\label{app:adm_dof}

This appendix gives the explicit ADM (Arnowitt--Deser--Misner)
constraint analysis of \eqref{eq:action} that underpins the body's
``two tensor modes, no propagating scalar'' claim and the resulting
predictions $c_{\rm GW}=c$ structurally (Sec.~\ref{sec:early}),
no fifth force at the propagating-scalar level (Sec.~\ref{sec:fifth}),
no scalar dipole radiation in binary inspirals (Sec.~\ref{sec:early}),
and the Branch-theorem's identification of $\Psi$ as a Lagrange
multiplier (Sec.~\ref{sec:psi_variational}). The result is internal
to the project; external verification, with explicit reference to
the mimetic-DOF analyses of Barvinsky, Capela--Ramazanov, and the
Firouzjahi--Gorji--Mansoori sequence
(\cite{Barvinsky2014,Capela2015,Firouzjahi2017,Gorji2018}), is
identified as a methodological priority (Sec.~\ref{sec:future}).

\subsection{ADM variables and Lagrangian density}

Standard $3+1$ decomposition gives lapse $N$, shift $N^i$, spatial
metric $h_{ij}$, with conjugate momenta $\pi_N$, $\pi_i$, $\pi^{ij}$;
plus the scalar pair $(\Psi,\pi_\Psi)$. Per spatial point, $11$
configuration variables and $22$ phase-space variables.

Using the Gauss--Codazzi decomposition
$R = R^{(3)} + (K_{ij}K^{ij}-K^2) + \text{total derivatives}$ with
$K_{ij}=(2N)^{-1}(\dot h_{ij}-\nabla_i N_j-\nabla_j N_i)$, the F01
gravity Lagrangian density is
\begin{equation}
\mathcal L_{\rm grav}
   = \frac{1}{16\pi}\,N\sqrt{h}\,\Psi\,\bigl[R^{(3)} + (K_{ij}K^{ij}-K^2)\bigr]
   + \text{(total derivatives)}.
\end{equation}

\textbf{Key observation.} $\Psi$ appears \emph{linearly}
multiplying the ADM density; $\dot\Psi$ does not appear anywhere
in $\mathcal L$, because $R^{(3)}$ contains only $h_{ij}$ and
spatial derivatives, $(K_{ij}K^{ij}-K^2)$ contains $\dot h_{ij}$ but
not $\dot\Psi$, and the total-derivative term
$\partial_t(\sqrt{h}\,\Psi K)$ is dropped from the bulk Lagrangian.
The matter Lagrangian $(1+f)\Lm$ contains no $\Psi$ at all ($f$ is
a matter-only functional of the distribution function $F(x,p)$).
Hence the full F01 Lagrangian density has \emph{no $\dot\Psi$
dependence}.

\subsection{Primary constraints}

The momentum conjugate to $\Psi$ is therefore
\begin{equation}
\pi_\Psi \;\equiv\; \partial\mathcal L/\partial\dot\Psi
\;=\; 0\quad\text{identically},
\end{equation}
a primary constraint in the Dirac sense and the hallmark signature
of a Lagrange multiplier. Together with the standard GR primaries
$\pi_N=0$, $\pi_i=0$, the F01 system has $5$ primary constraints.

\textbf{Cross-check vs.\ Brans--Dicke $\omega=0$.} BD gravity with
$\omega=0$ has the action
$S_{\rm BD} = (16\pi)^{-1}\!\int\!\sqrt{-g}\,\Psi R$ (no kinetic
term in the $\omega\to 0$ limit). The F01 gravity sector
\emph{is} this case (modulo the $(1+f)$ matter coupling); the
matter-side coupling is what distinguishes ISST. The Chiba
$f(R)\leftrightarrow$BD equivalence specialised to $f(R)=R$ is the
degenerate Einstein-Hilbert-with-prescribed-coupling case, which
is exactly F01 without $V(\Psi)$.

\subsection{Hamiltonian and secondary constraints}

The total Hamiltonian (modulo boundary terms) is
\begin{equation}
H = \int d^3x\,[\,N H_\perp + N^i H_i + u_N \pi_N + u^i\pi_i + u_\Psi\pi_\Psi\,],
\end{equation}
with $H_\perp=h^{-1/2}\pi^{ij}(\pi_{ij}-\tfrac12\pi h_{ij})/\Psi
+(\sqrt{h}/16\pi)\Psi R^{(3)}+\text{matter}$ and
$H_i=-2\nabla_j\pi^j{}_i+\text{matter flux}$.

Consistency of the primary constraints generates
\begin{align}
\dot\pi_N &= \{\pi_N,H\} \approx -H_\perp\;\Rightarrow\;H_\perp=0
&&\text{(secondary, GR scalar)} \\
\dot\pi_i &= \{\pi_i,H\} \approx -H_i\;\Rightarrow\;H_i=0
&&\text{(secondary, GR vector)} \\
\dot\pi_\Psi &= \{\pi_\Psi,H\} \approx -\delta H/\delta\Psi.
\end{align}
Collecting all $\Psi$-derivative terms of $H$ and demanding they
vanish is algebraically equivalent to the trace of the F01 metric
field equation,
\begin{equation}
-\Psi R + 3\Box\Psi = 8\pi T_{\rm eff},
\end{equation}
which on Branch~A ($R=0$, with $\nabla\Psi\neq 0$) reduces to
\begin{equation}
\Box\Psi = \tfrac{8\pi}{3}(1+f)T.
\end{equation}
This is the trace-determined evolution of $\Psi$: $\Psi$
\emph{propagates as a sourced second-order PDE} but does so as a
constraint determined by matter, not as an independent dynamical
field with its own variational equation. In Hamiltonian language,
the secondary constraint from $\dot\pi_\Psi=0$ is $R=0$ (Branch~A);
$\Psi$'s evolution is then fixed point-by-point by the trace
identity.

\subsection{Constraint classification and DOF count}

The pair $(\pi_\Psi, R=0)$ is first-class:
$\{\pi_\Psi(x),R(y)\}=-\partial R(y)/\partial\Psi(x)=0$ because $R$
contains no $\Psi$ (the gravity action is $\Psi R$, so $R$ itself
is a metric-only functional). With $R=R[g,\pi^g]$ and
$\Lm=\Lm[g,\chi_i]$ both $\Psi$-independent,
$\delta R/\delta\Psi=0$ identically, hence the Poisson bracket
vanishes weakly.

The first-class pair $(\pi_\Psi, R=0)$ generates the gauge
transformation $\delta\Psi=\varepsilon(x)$, $\delta g_{ij}=0$,
$\delta\pi^{ij}=0$, $\delta(\text{matter})=0$ on the constraint
surface. This is the formal statement that $\Psi$ is pure gauge
modulo the $R=0$ slice --- only its gradient appears in observables,
through $\Box\Psi$ and $\nabla_\mu\nabla_\nu\Psi$ in the metric
field equation.

The full constraint count is:
\begin{itemize}
\item First-class: $8$ (GR diffeomorphism: $\pi_N=0,H_\perp=0,\pi_i=0,H_i=0$, $\times 2$) plus $1$ from $(\pi_\Psi,R=0)$ taken as a first-class pair --- total $9$.
\item Second-class: $0$.
\end{itemize}
By the Dirac formula
$\text{DOF}_{\rm config} = (N_{\rm phase} - 2 N_{\rm 1C} - N_{\rm 2C})/2
= (22 - 18 - 0)/2 = \mathbf{2}$.

\textbf{F01 has two propagating degrees of freedom: the two graviton
polarisations of GR, no propagating scalar mode.}

\subsection{Distinction from mimetic gravity}

The mimetic-gravity DOF controversy is the natural concern raised
by F01's structural similarity to mimetic
gravity~\cite{ChamseddineMukhanov2013}. In mimetic gravity, a
Lagrange multiplier $\lambda$ enforces
$g^{\mu\nu}\partial_\mu\phi\,\partial_\nu\phi = -1$ on the scalar
$\phi$. This constraint forces $\phi$ to play the role of a
clock-like coordinate and produces a propagating dust-like mode
whose kinetic structure has been analysed by
Barvinsky~\cite{Barvinsky2014},
Capela--Ramazanov~\cite{Capela2015}, and the
Firouzjahi--Gorji--Mansoori
sequence~\cite{Firouzjahi2017,Gorji2018}; several authors argue
that the constraint creates a hidden propagating degree of freedom
with non-trivial dispersion in the caustic regime.

\emph{The F01 action does not contain this constraint.} No
Lagrange multiplier $\lambda$ enforcing a derivative condition
on $\Psi$ appears in \eqref{eq:action}, and $\nabla\Psi$ is not
constrained to be timelike unit. The trace identity
$\Box\Psi=(8\pi/3)(1+f)T$ that emerges from $\dot\pi_\Psi=0$ does
not impose a kinematic condition on $\nabla\Psi$ analogous to
mimetic's $g^{\mu\nu}\partial_\mu\phi\,\partial_\nu\phi=-1$;
it is a Poisson-type sourced equation determining $\Psi$'s
gradient in terms of matter, with no preferred-frame structure
on $\Psi$. The mimetic-dust hidden-mode mechanism therefore does
not directly transfer to F01.

\textbf{Open: external verification.} The internal ADM count above
is consistent with the no-propagating-scalar reading of the F01
action and structurally distinct from the mimetic constraint.
However, given the mimetic-DOF literature's history of subtle
hidden modes in apparently constrained systems, an independent
ADM/Hamiltonian-constraint analysis by an external group --- with
explicit cross-comparison against the cited mimetic literature ---
is required to close the DOF question to the level a hostile
referee would demand. This external verification is identified
as a methodological priority (Sec.~\ref{sec:future}).

\subsection{What the trace-determined evolution does and does not mean}

The trace equation $\Box\Psi=(8\pi/3)(1+f)T$ is a second-order PDE
in $\Psi$. It is sourced by matter. In the PDE sense, $\Psi$
\emph{propagates}: hyperbolic equations on a Lorentzian background
admit characteristic propagation along null cones. The phrase
``non-dynamical $\Psi$'' is misleading at the PDE level and is
replaced throughout the paper by ``trace-determined,'' ``metrically
slaved,'' or ``constraint-propagated'' (Sec.~\ref{sec:psi_variational}).

What is true is the variational statement: $\Psi$ is not
independently extremised in the action --- there is no
$\delta S/\delta\Psi=0$ EOM. Its evolution is entirely determined
by the metric and matter equations through the trace constraint.
Whether this constraint-propagation constitutes an independent
dynamical degree of freedom is a Hamiltonian-constraint question
(this appendix), not a PDE-classification question on the trace
equation. The Hamiltonian answer above is that the
$(\pi_\Psi,R=0)$ pair is first-class, generating the gauge
transformation $\delta\Psi=\varepsilon$, and contributes zero
propagating DOF.

The body's ``no fifth force'' claim is therefore the narrow
statement that no propagating-scalar-mediated fifth force exists
in F01: the $\Psi$ sector contributes no propagating mode that
could mediate a Yukawa-type interaction between matter test
particles. The $(1+f)$ rescaling of the gravitational source is a
matter-side coupling, not a fifth-force mediator; it has no
propagator and no finite range. Conversely, the gradients of
$\Psi$ that do exist in space and time --- driven by matter sources
via the trace equation --- do enter local observables (the LLR
$\dot G/G$ prediction in Sec.~\ref{sec:solar} is exactly such an
effect) through boundary-condition drag on bound systems embedded
in an evolving cosmological background.

%==========================================================
\begin{thebibliography}{99}

% --- Foundations: scalar-tensor and information-theoretic ---

\bibitem{BransDicke1961} C.~Brans and R.~H.~Dicke, ``Mach's
principle and a relativistic theory of gravitation,''
\textit{Phys.~Rev.} \textbf{124}, 925 (1961).

\bibitem{OHanlonTupper1972} J.~O'Hanlon and B.~O.~J.~Tupper,
``Vacuum-field solutions in the Brans--Dicke theory,''
\textit{Nuovo Cim.~B} \textbf{7}, 305 (1972).

\bibitem{Horndeski1974} G.~W.~Horndeski, ``Second-order scalar-tensor
field equations in a four-dimensional space,''
\textit{Int.~J.~Theor.~Phys.} \textbf{10}, 363 (1974).

\bibitem{FujiiMaeda2003} Y.~Fujii and K.-i.~Maeda, \textit{The
Scalar--Tensor Theory of Gravitation} (Cambridge Univ.~Press, 2003).

\bibitem{Sotiriou2010} T.~P.~Sotiriou and V.~Faraoni, ``$f(R)$
theories of gravity,'' \textit{Rev.~Mod.~Phys.} \textbf{82}, 451
(2010).

\bibitem{Olmo2005} G.~J.~Olmo, ``The gravity Lagrangian according to
Solar System experiments,'' \textit{Phys.~Rev.~Lett.} \textbf{95},
261102 (2005).

\bibitem{Will2014} C.~M.~Will, ``The confrontation between general
relativity and experiment,'' \textit{Living Rev.~Rel.} \textbf{17},
4 (2014).

\bibitem{DamourEspositoFarese1992} T.~Damour and G.~Esposito-Farese,
``Tensor-multi-scalar theories of gravitation,''
\textit{Class.~Quant.~Grav.} \textbf{9}, 2093 (1992).

\bibitem{Khoury2004} J.~Khoury and A.~Weltman, ``Chameleon
cosmology,'' \textit{Phys.~Rev.~D} \textbf{69}, 044026 (2004).

\bibitem{Shore1980} J.~E.~Shore and R.~W.~Johnson, ``Axiomatic
derivation of the principle of maximum entropy and the principle of
minimum cross-entropy,'' \textit{IEEE Trans.~Inf.~Theory} \textbf{26},
26 (1980).

\bibitem{Aczel1966} J.~Acz\'el, \textit{Lectures on Functional
Equations and Their Applications} (Academic Press, New York, 1966).

\bibitem{Jaynes1957} E.~T.~Jaynes, ``Information theory and
statistical mechanics,'' \textit{Phys.~Rev.} \textbf{106}, 620 (1957).

\bibitem{Bekenstein1993} J.~D.~Bekenstein, ``Relation between
physical and gravitational geometry,'' \textit{Phys.~Rev.~D}
\textbf{48}, 3641 (1993).

\bibitem{Verlinde2011} E.~Verlinde, ``On the origin of gravity and
the laws of Newton,'' \textit{JHEP} \textbf{04}, 029 (2011)
[arXiv:1001.0785].

\bibitem{Verlinde2017} E.~Verlinde, ``Emergent gravity and the dark
universe,'' \textit{SciPost Phys.} \textbf{2}, 016 (2017).

\bibitem{Padmanabhan2010} T.~Padmanabhan, ``Thermodynamical aspects
of gravity: new insights,''
\textit{Rep.~Prog.~Phys.} \textbf{73}, 046901 (2010)
[arXiv:0911.5004].

% --- Branch / static-fluid / lensing internal ---

\bibitem{F03F08foundations} S.~Brailsford, F03--F08 Branch theorems,
ISST repository, \texttt{foundations/02\_branch\_theorem/} (2026).

\bibitem{MASTERledger} S.~Brailsford, F-task ledger,
ISST repository, \texttt{MASTER.md} (2026).

\bibitem{ISSTRepository} S.~Brailsford and Lily, ISST companion site
and live Derivation Passport, \url{https://lily-labs.co.uk/isst}
(2026).

% --- Background / averaging / cosmology ---

\bibitem{Buchert2000} T.~Buchert, ``On average properties of
inhomogeneous cosmologies,'' \textit{Gen.~Rel.~Grav.} \textbf{32}, 105
(2000).

\bibitem{BuchertRasanen2008} T.~Buchert and S.~R\"as\"anen,
``Backreaction in late-time cosmology,''
\textit{Annu.~Rev.~Nucl.~Part.~Sci.} \textbf{62}, 57 (2012).

\bibitem{Wiltshire2007} D.~L.~Wiltshire, ``Cosmic clocks, cosmic
variance and cosmic averages,'' \textit{New J.~Phys.} \textbf{9},
377 (2007).

\bibitem{Wiltshire2009} D.~L.~Wiltshire, ``Average observational
quantities in the timescape cosmology,'' \textit{Phys.~Rev.~D}
\textbf{80}, 123512 (2009).

\bibitem{WiltshirePRL2007} D.~L.~Wiltshire, ``Exact solution to the
averaging problem in cosmology,'' \textit{Phys.~Rev.~Lett.}
\textbf{99}, 251101 (2007).

\bibitem{Dam2017} L.~H.~Dam, A.~Heinesen and D.~L.~Wiltshire,
``Apparent cosmic acceleration from Type Ia supernovae,''
\textit{Mon.~Not.~Roy.~Astron.~Soc.} \textbf{472}, 835 (2017);
arXiv:1706.07236.

\bibitem{NazerWiltshire2015} M.~A.~Nazer and D.~L.~Wiltshire, ``Cosmic
microwave background anisotropies in the timescape cosmology,''
\textit{Phys.~Rev.~D} \textbf{91}, 063519 (2015).

\bibitem{LewisChallinor2011} A.~Lewis and A.~Challinor, ``CAMB: Code
for Anisotropies in the Microwave Background,'' Astrophysics Source
Code Library, ascl:1102.026 (2011); A.~Lewis, A.~Challinor and
A.~Lasenby, ``Efficient computation of CMB anisotropies in closed FRW
models,'' \textit{ApJ} \textbf{538}, 473 (2000).

\bibitem{Sheth2004} R.~K.~Sheth and R.~van~de~Weygaert, ``A hierarchy
of voids: much ado about nothing,''
\textit{MNRAS} \textbf{350}, 517 (2004).

% --- Cosmological observations ---

\bibitem{Planck2018} Planck Collaboration, ``Planck 2018 results.~VI.
Cosmological parameters,'' \textit{A\&A} \textbf{641}, A6 (2020).

\bibitem{Riess2022} A.~G.~Riess \emph{et al.}, ``A comprehensive
measurement of the local value of the Hubble constant with
1~km/s/Mpc uncertainty from the {\it Hubble Space Telescope} and the
SH0ES team,'' \textit{Astrophys.~J.~Lett.} \textbf{934}, L7 (2022).

\bibitem{VerdeTreuRiess2019} L.~Verde, T.~Treu and A.~G.~Riess,
``Tensions between the early and late universe,''
\textit{Nature Astron.} \textbf{3}, 891 (2019).

\bibitem{DiValentino2021} E.~Di~Valentino \emph{et al.}, ``In the
realm of the Hubble tension---a review of solutions,''
\textit{Class.~Quant.~Grav.} \textbf{38}, 153001 (2021).

\bibitem{DESIDR1} A.~G.~Adame \emph{et al.} (DESI Collaboration),
``DESI 2024~VI: Cosmological constraints from the measurements of
baryon acoustic oscillations,'' arXiv:2404.03002 (2024).

\bibitem{DESIDR2} M.~Abdul-Karim \emph{et al.} (DESI Collaboration),
``DESI DR2 results~II: Measurements of baryon acoustic oscillations
and cosmological constraints,'' \textit{Phys.~Rev.~D} \textbf{112},
083515 (2025).

\bibitem{KiDS1000} C.~Heymans \emph{et al.}, ``KiDS-1000 cosmology:
Multi-probe weak gravitational lensing and spectroscopic galaxy
clustering constraints,'' \textit{A\&A} \textbf{646}, A140 (2021).

\bibitem{DESY3} DES Collaboration, ``Dark Energy Survey Year~3
results: Cosmological constraints from galaxy clustering and weak
lensing,'' \textit{Phys.~Rev.~D} \textbf{105}, 023520 (2022).

\bibitem{HSCY3} S.~More \emph{et al.}, ``Hyper Suprime-Cam Year~3
results: Cosmology from galaxy clustering and weak lensing,''
\textit{Phys.~Rev.~D} \textbf{108}, 123520 (2023).

\bibitem{EuclidForecast} L.~Amendola \emph{et al.} (Euclid
Theory Working Group), ``Cosmology and fundamental physics with the
Euclid satellite,'' \textit{Living Rev.~Rel.} \textbf{21}, 2 (2018).

\bibitem{EuclidYellow} R.~Laureijs \emph{et al.} (Euclid Collaboration),
``Euclid definition study report,'' arXiv:1110.3193 (2011).

% --- Growth / RSD / structure ---

\bibitem{BeutlerEtAl2012} F.~Beutler \emph{et al.}, ``The 6dF galaxy
survey: $z\approx 0$ measurements of the growth rate and $\sigma_8$,''
\textit{MNRAS} \textbf{423}, 3430 (2012).

\bibitem{HowlettEtAl2015} C.~Howlett \emph{et al.}, ``2MTF VI.
Measuring the velocity power spectrum,''
\textit{MNRAS} \textbf{449}, 848 (2015).

\bibitem{AlamEtAl2017} S.~Alam \emph{et al.}, ``The clustering of
galaxies in the completed SDSS-III Baryon Oscillation Spectroscopic
Survey,'' \textit{MNRAS} \textbf{470}, 2617 (2017).

\bibitem{BlakeEtAl2012} C.~Blake \emph{et al.}, ``The WiggleZ Dark
Energy Survey: joint measurements of the expansion and growth history
at $z<1$,'' \textit{MNRAS} \textbf{425}, 405 (2012).

\bibitem{PezzottaEtAl2017} A.~Pezzotta \emph{et al.}, ``The VIMOS
public extragalactic redshift survey: the growth of structure at
$0.5<z<1.2$,'' \textit{A\&A} \textbf{604}, A33 (2017).

\bibitem{Adil2024} S.~A.~Adil, \"O.~Akarsu, E.~Di~Valentino,
R.~C.~Nunes, E.~\"Ozulker, A.~A.~Sen and E.~Specogna, ``Hints of
a physical origin for the $S_8$ tension,''
\textit{Nature Astron.} \textbf{8}, 1115 (2024);
\href{https://doi.org/10.1038/s41550-024-02262-3}{doi:10.1038/s41550-024-02262-3}.

\bibitem{ChevallierPolarski2001} M.~Chevallier and D.~Polarski,
``Accelerating universes with scaling dark matter,''
\textit{Int.~J.~Mod.~Phys.~D} \textbf{10}, 213 (2001).

\bibitem{Linder2003} E.~V.~Linder, ``Exploring the expansion history
of the universe,'' \textit{Phys.~Rev.~Lett.} \textbf{90}, 091301
(2003).

% --- Galactic dynamics ---

\bibitem{Milgrom1983} M.~Milgrom, ``A modification of the Newtonian
dynamics as a possible alternative to the hidden mass hypothesis,''
\textit{ApJ} \textbf{270}, 365 (1983).

\bibitem{LyndenBell1967} D.~Lynden-Bell, ``Statistical mechanics of
violent relaxation in stellar systems,''
\textit{MNRAS} \textbf{136}, 101 (1967).

\bibitem{Lelli2016} F.~Lelli, S.~S.~McGaugh and J.~M.~Schombert,
``SPARC: Mass models for 175 disk galaxies with Spitzer photometry
and accurate rotation curves,'' \textit{AJ} \textbf{152}, 157 (2016).

\bibitem{McGaugh2012} S.~S.~McGaugh, ``The baryonic Tully--Fisher
relation of gas-rich galaxies as a test of $\Lambda$CDM and MOND,''
\textit{AJ} \textbf{143}, 40 (2012).

\bibitem{McGaughLelliSchombert2016} S.~S.~McGaugh, F.~Lelli and
J.~M.~Schombert, ``Radial acceleration relation in rotationally
supported galaxies,'' \textit{Phys.~Rev.~Lett.} \textbf{117}, 201101
(2016).

\bibitem{LelliEtAl2017} F.~Lelli \emph{et al.}, ``One law to bind
them all: the radial acceleration relation of galaxies,''
\textit{ApJ} \textbf{836}, 152 (2017).

\bibitem{Bekenstein1984} J.~D.~Bekenstein and M.~Milgrom, ``Does the
missing mass problem signal the breakdown of Newtonian gravity?,''
\textit{ApJ} \textbf{286}, 7 (1984).

\bibitem{BekensteinTeVeS2004} J.~D.~Bekenstein, ``Relativistic
gravitation theory for the modified Newtonian dynamics paradigm,''
\textit{Phys.~Rev.~D} \textbf{70}, 083509 (2004).

\bibitem{FamaeyMcGaugh2012} B.~Famaey and S.~S.~McGaugh, ``Modified
Newtonian dynamics (MOND): observational phenomenology and
relativistic extensions,'' \textit{Living Rev.~Rel.} \textbf{15}, 10
(2012).

\bibitem{BanikZhao2022} I.~Banik and H.~Zhao, ``From galactic bars to
the Hubble tension: weighing up the astrophysical evidence for
Milgromian gravity,'' \textit{Symmetry} \textbf{14}, 1331 (2022).

\bibitem{SandersMcGaugh2002} R.~H.~Sanders and S.~S.~McGaugh,
``Modified Newtonian dynamics as an alternative to dark matter,''
\textit{Annu.~Rev.~Astron.~Astrophys.} \textbf{40}, 263 (2002).

\bibitem{Tully2008} R.~B.~Tully, ``Light-to-mass variations with
environment,'' \textit{AJ} \textbf{135}, 1488 (2008).

\bibitem{MadauDickinson2014} P.~Madau and M.~Dickinson, ``Cosmic
star-formation history,''
\textit{Annu.~Rev.~Astron.~Astrophys.} \textbf{52}, 415 (2014).

% --- Clusters ---

\bibitem{Markevitch2004} M.~Markevitch \emph{et al.}, ``Direct
constraints on the dark matter self-interaction cross-section from
the merging galaxy cluster 1E~0657-56,''
\textit{ApJ} \textbf{606}, 819 (2004).

\bibitem{Clowe2006} D.~Clowe \emph{et al.}, ``A direct empirical proof
of the existence of dark matter,'' \textit{ApJ} \textbf{648}, L109
(2006).

\bibitem{Paraficz2016} D.~Paraficz \emph{et al.}, ``The Bullet
cluster at its best: weighing stars, gas, and dark matter,''
\textit{A\&A} \textbf{594}, A121 (2016).

\bibitem{Cha2025} S.~Cha, B.~Y.~Cho, H.~Joo, \emph{et al.},
``A high-caliber view of the Bullet Cluster through JWST strong and
weak lensing analyses,''
\textit{ApJ} \textbf{987}, L15 (2025); arXiv:2503.21870.

\bibitem{Rihtarsic2026} G.~Rihtar\v{s}i\v{c}, M.~Brada\v{c},
G.~Desprez, \emph{et al.}, ``Mapping dark matter in the Bullet
Cluster using JWST imaging and spectroscopy,''
arXiv:2601.22245 (2026).

\bibitem{Hernandez2026} X.~Hernandez, ``A consistent MOND modelling
of the Bullet Cluster,'' arXiv:2604.10811 (2026).

\bibitem{Mahdavi2007} A.~Mahdavi \emph{et al.}, ``A dark core in
Abell 520,'' \textit{ApJ} \textbf{668}, 806 (2007).

\bibitem{Jee2012} M.~J.~Jee \emph{et al.}, ``A study of the dark core
in A520 with the Hubble Space Telescope: the mystery deepens,''
\textit{ApJ} \textbf{747}, 96 (2012).

\bibitem{Clowe2012} D.~Clowe \emph{et al.}, ``On dark peaks and
missing mass: a weak-lensing mass reconstruction of the merging
cluster system A520,'' \textit{ApJ} \textbf{758}, 128 (2012).

\bibitem{Jee2014} M.~J.~Jee \emph{et al.}, ``Hubble Space
Telescope/Advanced Camera for Surveys confirmation of the dark
substructure in A520,'' \textit{ApJ} \textbf{783}, 78 (2014).

\bibitem{Deshev2017} B.~Deshev \emph{et al.}, ``Galaxy evolution in
merging clusters: the passive core of the train wreck cluster of
galaxies, A520,'' \textit{A\&A} \textbf{607}, A131 (2017).

\bibitem{Loubser2020} S.~I.~Loubser \emph{et al.}, ``Dynamical masses
of brightest cluster galaxies~I: stellar velocity anisotropy and
mass-to-light ratios,'' \textit{MNRAS} \textbf{496}, 1857 (2020).

\bibitem{ChandrasekharDF} S.~Chandrasekhar, ``Dynamical friction.~I.
General considerations: the coefficient of dynamical friction,''
\textit{ApJ} \textbf{97}, 255 (1943).

\bibitem{BinneyTremaine2008} J.~Binney and S.~Tremaine,
\textit{Galactic Dynamics}, 2nd ed.\ (Princeton Univ.~Press, 2008).

% --- GW and Solar System ---

\bibitem{LIGO2017} B.~P.~Abbott \emph{et al.}\ (LIGO/Virgo),
``GW170817: Observation of gravitational waves from a binary neutron
star inspiral,'' \textit{Phys.~Rev.~Lett.} \textbf{119}, 161101 (2017).

\bibitem{LIGOGW170817EM} B.~P.~Abbott \emph{et al.}\
(LIGO/Virgo/Fermi/INTEGRAL), ``Gravitational waves and gamma-rays
from a binary neutron star merger: GW170817 and GRB~170817A,''
\textit{Astrophys.~J.~Lett.} \textbf{848}, L13 (2017).

\bibitem{Creminelli2017} P.~Creminelli and F.~Vernizzi, ``Dark
energy after GW170817,'' \textit{Phys.~Rev.~Lett.} \textbf{119},
251302 (2017).

\bibitem{Bertotti2003} B.~Bertotti, L.~Iess and P.~Tortora, ``A test
of general relativity using radio links with the Cassini spacecraft,''
\textit{Nature} \textbf{425}, 374 (2003).

\bibitem{Hofmann2018} F.~Hofmann and J.~M\"uller, ``Relativistic tests
with lunar laser ranging,'' \textit{Class.~Quant.~Grav.} \textbf{35},
035015 (2018).

\bibitem{Williams2004} J.~G.~Williams, S.~G.~Turyshev and
D.~H.~Boggs, ``Progress in lunar laser ranging tests of relativistic
gravity,'' \textit{Phys.~Rev.~Lett.} \textbf{93}, 261101 (2004).

\bibitem{Damour2002} T.~Damour, F.~Piazza and G.~Veneziano,
``Runaway dilaton and equivalence principle violations,''
\textit{Phys.~Rev.~Lett.} \textbf{89}, 081601 (2002).

% --- Direct detection nulls ---

\bibitem{XENONnT2023} XENON Collaboration, ``First dark matter search
with nuclear recoils from the XENONnT experiment,''
\textit{Phys.~Rev.~Lett.} \textbf{131}, 041003 (2023).

\bibitem{LUXZEPLIN2023} LZ Collaboration, ``First dark matter search
results from the LUX-ZEPLIN (LZ) experiment,''
\textit{Phys.~Rev.~Lett.} \textbf{131}, 041002 (2023).

\bibitem{ADMX2021} ADMX Collaboration, ``Search for invisible axion
dark matter in the 3.3--4.2~$\mu$eV mass range,''
\textit{Phys.~Rev.~Lett.} \textbf{127}, 261803 (2021).

% --- Halo / structure tensions ---

\bibitem{BullockBoylanKolchin2017} J.~S.~Bullock and
M.~Boylan-Kolchin, ``Small-scale challenges to the $\Lambda$CDM
paradigm,'' \textit{Annu.~Rev.~Astron.~Astrophys.} \textbf{55}, 343
(2017).

\bibitem{LabbeJWST2023} I.~Labb\'e \emph{et al.}, ``A population of
red candidate massive galaxies $\sim 600$~Myr after the Big Bang,''
\textit{Nature} \textbf{616}, 266 (2023).

\bibitem{BoylanKolchin2023} M.~Boylan-Kolchin, ``Stress-testing
$\Lambda$CDM with high-redshift galaxy candidates,''
\textit{Nature Astron.} \textbf{7}, 731 (2023).

\bibitem{Lesgourgues2011} J.~Lesgourgues, ``The Cosmic Linear
Anisotropy Solving System (CLASS).~I.~Overview,''
arXiv:1104.2932 (2011); J.~Lesgourgues and T.~Tram, ``The Cosmic
Linear Anisotropy Solving System (CLASS).~IV.~Efficient
implementation of non-cold relics,''
\textit{JCAP} \textbf{09}, 032 (2011).

\bibitem{ChamseddineMukhanov2013} A.~H.~Chamseddine and V.~Mukhanov,
``Mimetic dark matter,''
\textit{JHEP} \textbf{11}, 135 (2013) [arXiv:1308.5410].

\bibitem{vanDokkum2018} P.~van Dokkum \emph{et al.}, ``A galaxy
lacking dark matter,''
\textit{Nature} \textbf{555}, 629 (2018) [arXiv:1803.10237].

\bibitem{Trujillo2019} I.~Trujillo \emph{et al.}, ``A distance of
13~Mpc resolves the claimed anomalies of the galaxy lacking dark
matter NGC~1052--DF2,''
\textit{MNRAS} \textbf{486}, 1192 (2019) [arXiv:1806.10141].

\bibitem{Lelli2015} F.~Lelli, P.-A.~Duc, E.~Brinks, F.~Bournaud,
S.~S.~McGaugh, U.~Lisenfeld, T.~X.~Thuan, V.~Belles, V.~Boquien,
P.~M.~Weilbacher, L.~Boquien, J.~Braine, E.~Hibbard, J.~Iglesias-P\'aramo
and L.~Verdes-Montenegro,
``Gas dynamics in tidal dwarf galaxies: disc formation at $z=0$,''
\textit{A\&A} \textbf{584}, A113 (2015) [arXiv:1509.05415].

\bibitem{Mancera2022} P.~E.~Mancera Pi\~na \emph{et al.}, ``Off the
baryonic Tully--Fisher relation: A population of baryon-dominated
ultra-diffuse galaxies,''
\textit{ApJL} \textbf{921}, L9 (2021) [arXiv:2107.08029].

\bibitem{Penrose1979} R.~Penrose, ``Singularities and time-asymmetry,''
in \textit{General Relativity: an Einstein Centenary Survey},
S.~W.~Hawking and W.~Israel, eds. (Cambridge University Press, 1979),
pp.~581--638.

\bibitem{Jacobson1995} T.~Jacobson, ``Thermodynamics of spacetime:
the Einstein equation of state,''
\textit{Phys.~Rev.~Lett.} \textbf{75}, 1260 (1995)
[arXiv:gr-qc/9504004].

\bibitem{Israel1966} W.~Israel, ``Singular hypersurfaces and thin
shells in general relativity,''
\textit{Nuovo Cim.~B} \textbf{44}, 1 (1966); erratum
\textit{ibid.} \textbf{48}, 463 (1967).

\bibitem{Hartle1967} J.~B.~Hartle, ``Slowly rotating relativistic
stars.~I.~Equations of structure,''
\textit{ApJ} \textbf{150}, 1005 (1967).

\bibitem{BorsanyiQCD2016} S.~Bors\'anyi \emph{et al.}, ``Calculation
of the axion mass based on high-temperature lattice quantum
chromodynamics,''
\textit{Nature} \textbf{539}, 69 (2016) [arXiv:1606.07494]; see
Fig.~3 for $g_{*,s}(T)$ uncertainties through the QCD crossover.

\bibitem{Barvinsky2014} A.~O.~Barvinsky, ``Dark matter as a
ghost-free conformal extension of Einstein theory,''
\textit{JCAP} \textbf{01}, 014 (2014) [arXiv:1311.3111].

\bibitem{Capela2015} F.~Capela and S.~Ramazanov, ``Modified dust and
the small scale crisis in CDM,''
\textit{JCAP} \textbf{04}, 051 (2015) [arXiv:1412.2051].

\bibitem{Firouzjahi2017} H.~Firouzjahi, M.~A.~Gorji and
S.~A.~Hosseini Mansoori, ``Instabilities in mimetic matter
perturbations,''
\textit{JCAP} \textbf{07}, 031 (2017) [arXiv:1703.02923].

\bibitem{Gorji2018} M.~A.~Gorji, S.~Mukohyama, H.~Firouzjahi and
S.~A.~Hosseini Mansoori, ``Gauge-invariant cosmological perturbation
theory in mimetic gravity,''
\textit{JCAP} \textbf{08}, 047 (2018) [arXiv:1802.05920].

\end{thebibliography}

\end{document}
